% AA Classe Del Documento
\documentclass[11pt, a4paper]{article}

% AA Tipografia E Lingua
% BB Lingua E Codifica Font
\usepackage[italian]{babel}
\usepackage[T1]{fontenc}
\usepackage{lmodern}

% BB Geometria Pagina E Microtipografia
\usepackage[margin=2cm, top=1.5cm, bottom=1.5cm]{geometry}
\usepackage{parskip}
\usepackage{microtype}

% AA Gestione Dei Colori
% BB Pacchetto Colori
\usepackage[table]{xcolor}

% BB Palette Colori Personalizzata
\definecolor{darkgreen}{RGB}{0,100,0}
\definecolor{theoryblue}{RGB}{0,84,147}
\definecolor{backcolour}{RGB}{245,245,245}
\definecolor{codegray}{rgb}{0.5,0.5,0.5}

\definecolor{codebg}{HTML}{F8F9FA}
\definecolor{codeframe}{HTML}{E2E8F0}
\definecolor{codekeywords}{HTML}{7C3AED}
\definecolor{codedocker}{HTML}{0284C7}
\definecolor{codecomments}{HTML}{64748B}
\definecolor{codestrings}{HTML}{0D9488}
\definecolor{codenumbers}{HTML}{94A3B8}
\definecolor{codetext}{HTML}{1E293B}

% AA Matematica E Notazione Scientifica
\usepackage{amsmath, amssymb, amsfonts}
\usepackage{mathtools}
\usepackage{siunitx}

% BB OTTIMIZZAZIONE MATRICI LUNGHE PER LA STAMPA
\setlength{\arraycolsep}{10pt}
\usepackage{bm}

% AA Tabelle Ed Allineamenti
\usepackage{booktabs}
\usepackage{array}
\usepackage{multirow}
\usepackage{tabularx}
\usepackage{arydshln}

% AA Grafica E Disegno Vettoriale
% BB Gestione Figure E Posizionamento
\usepackage{graphicx}
\graphicspath{{./Images/RipassoNetworking/}}
\usepackage{subcaption}
\usepackage{float}

% BB TikZ E Pgfplots
\usepackage{tikz}
\usetikzlibrary{arrows.meta, patterns, angles, quotes, calc}
\usetikzlibrary{positioning}
\usepackage{pgfplots}
\pgfplotsset{compat=1.18}

% AA Listati Di Codice (Listings)
% BB Configurazione Generale
\usepackage{listings}
\usepackage{upquote}

% BB Configurazione Generale
\lstset{
    basicstyle=\ttfamily\footnotesize\color{codetext},
    backgroundcolor=\color{codebg},
    frame=single,
    rulecolor=\color{codeframe},
    framesep=6pt,
    xleftmargin=8pt,
    xrightmargin=8pt,
    breaklines=true,
    breakatwhitespace=true,
    showstringspaces=false,
    keepspaces=true,
    columns=flexible,
    upquote=true,
    tabsize=2,
    captionpos=b,
    % Stili e Colori Predefiniti Globali
    keywordstyle=\bfseries\color{codekeywords},
    commentstyle=\itshape\color{codecomments},
    stringstyle=\color{codestrings}
}

% BB Definizione Dei Linguaggi

% CC YAML
\lstdefinelanguage{yaml}{
    keywords={true,false,null,yes,no,on,off},
    keywordstyle=\bfseries\color{codekeywords},
    comment=[l]{\#},
    commentstyle=\itshape\color{codecomments},
    stringstyle=\color{codestrings},
    morestring=[b]',
    morestring=[b]",
    sensitive=false
}

% CC Dockerfile
\lstdefinelanguage{dockerfile}{
    keywords={FROM, RUN, CMD, LABEL, EXPOSE, ENV, ADD, COPY, ENTRYPOINT,
              VOLUME, USER, WORKDIR, ARG, ONBUILD, STOPSIGNAL, HEALTHCHECK,
              SHELL, AS},
    keywordstyle=\bfseries\color{codedocker},
    comment=[l]{\#},
    commentstyle=\itshape\color{codecomments},
    stringstyle=\color{codestrings},
    morestring=[b]",
    morestring=[b]',
    sensitive=true
}

% CC INI
\lstdefinelanguage{ini}{
    comment=[l]{\#},
    morecomment=[l]{;},
    commentstyle=\itshape\color{codecomments},
    stringstyle=\color{codestrings},
    morestring=[b]",
    morestring=[b]',
    keywordstyle=\bfseries\color{codekeywords},
    sensitive=false
}

% CC JSON
\lstdefinelanguage{json}{
    keywords={true,false,null},
    keywordstyle=\bfseries\color{codekeywords},
    comment=[l]{//},
    morecomment=[s]{/*}{*/},
    commentstyle=\itshape\color{codecomments},
    stringstyle=\color{codestrings},
    morestring=[b]",
    sensitive=true
}

% CC PYTHON (Dialetto nativo pulito)
\lstset{defaultdialect=[3]Python}
\lstdefinelanguage{python}[]{Python}{
    morekeywords={as, async, await, True, False, None},
    sensitive=true
}

% CC BASH / SHELL
\lstdefinelanguage{bash}[]{sh}{
    morekeywords={docker, cd, run, build, stop, rm, rmi, push, pull, network, compose, openstack, rq, worker, keystone-manage, time, sleep, echo, exit, sudo, export, systemctl},
    sensitive=true
}

% CC HTML (Autonomo, previene il loop infinito del parser XML/HTML)
\lstdefinelanguage{html}{
    keywords={img, src, alt, class, id, div, span, table, tr, td, th, href, a, p, h1, h2, h3, form, input, button},
    keywordstyle=\bfseries\color{codekeywords},
    comment=[s]{<!--}{-->},
    commentstyle=\itshape\color{codecomments},
    stringstyle=\color{codestrings},
    morestring=[b]",
    morestring=[b]',
    sensitive=false
}

% CC SQL / CYPHER (Autonomo con supporto per // di Cypher)
\lstdefinelanguage{sql}{
    keywords={MATCH, RETURN, DETACH, DELETE, WHERE, CREATE, SET, WITH, ORDER, BY, ASC, DESC, LIMIT, SELECT, INSERT, INTO, VALUES, FROM, UPDATE, DROP, TABLE, JOIN, ON, AND, OR, NOT},
    keywordstyle=\bfseries\color{codekeywords},
    comment=[l]{//},
    morecomment=[l]{--},
    commentstyle=\itshape\color{codecomments},
    stringstyle=\color{codestrings},
    morestring=[b]',
    morestring=[b]",
    sensitive=false
}

% AA Algoritmi E Pseudocodice
\usepackage[ruled,vlined,linesnumbered]{algorithm2e}
\SetKwData{Dati}{Input}
\SetKwData{Risultato}{Output}
\SetKw{Inizializza}{Inizializza}
\SetKwFor{Per}{per}{fare}{fine per}
\SetKwComment{Commento}{// }{}

% BB Costrutti Di Controllo Aggiuntivi
\SetKwIF{Se}{AltrSe}{Altrimenti}{Se}{Allora}{Altrimenti Se}{Altrimenti}{Fine Se}
\SetKwFor{Mentre}{mentre}{Fare}{Fine Mentre}

% AA Struttura, Titoli Ed Ambienti
% BB Personalizzazione Titoli Sezioni
\usepackage{titlesec}
\titleformat{\section}{\normalfont\LARGE\bfseries}{\thesection}{1em}{}
\titleformat{\subsection}{\normalfont\Large\bfseries}{\thesubsection}{1em}{}
\titleformat{\subsubsection}{\normalfont\large\bfseries}{\thesubsubsection}{1em}{}

% BB Teoremi Ed Ambienti Ausiliari
\usepackage{comment}
\newtheorem{theorem}{Definizione}

% BB Liste Ed Elenchi Puntati
\usepackage{enumitem}
\setlist{noitemsep}
\usepackage{pifont}

% AA Riferimenti, Link Ipertestuali E Segnalibri
\usepackage[unicode]{hyperref}
\hypersetup{
    colorlinks=true,
    linkcolor=blue,
    urlcolor=blue,
    citecolor=blue,
    pdftitle={1) Ripasso Networking IP E Sicurezza},
    pdfauthor={Francesco Lo Verde},
    pdfpagemode=FullScreen,
}
\usepackage{bookmark}

% AA Metadati Del Documento
\title{Ripasso Networking IP E Sicurezza}

\date{\today}

\author{Francesco Lo Verde}
\begin{document}
\maketitle

\section{Architettura Di Rete E Modelli Di Riferimento}

La progettazione dei moderni sistemi telematici, delle infrastrutture \emph{cloud} e delle reti virtualizzate affonda le proprie radici teoriche nei modelli di rete a strati. L'astrazione a livelli permette di scomporre un problema complesso — la trasmissione affidabile ed efficiente di flussi informativi tra processi eterogenei eseguiti su nodi fisicamente distribuiti — in sotto-problemi indipendenti, standardizzati e cooperanti.

\subsection{Il Modello Di Riferimento OSI E La Comunicazione Tra Nodi}

Il modello di riferimento OSI (\emph{Open Systems Interconnection}), standardizzato dall'\textbf{ISO} (\emph{International Organization for Standardization}), definisce un'architettura gerarchica a sette livelli, numerati dal basso verso l'alto:
\begin{enumerate}[noitemsep]
    \item \textbf{Livello Fisico} (\emph{Physical Layer}): trasmissione della sequenza binaria grezza sul mezzo trasmissivo;
    \item \textbf{Livello Collegamento Dati} (\emph{Data Link Layer}): trasferimento affidabile o non affidabile di trame tra nodi adiacenti appartenenti al medesimo segmento fisico o logico;
    \item \textbf{Livello Di Rete} (\emph{Network Layer}): instradamento (\emph{routing}) e commutazione logica dei pacchetti attraverso topologie interconnesse;
    \item \textbf{Livello Di Trasporto} (\emph{Transport Layer}): comunicazione affidabile o non affidabile \emph{end-to-end} (processo-a-processo);
    \item \textbf{Livello Di Sessione} (\emph{Session Layer}): apertura, sincronizzazione e gestione del dialogo tra entità applicative;
    \item \textbf{Livello Di Presentazione} (\emph{Presentation Layer}): negoziazione delle sintassi, codifica, cifratura e serializzazione del dato;
    \item \textbf{Livello Applicativo} (\emph{Application Layer}): punto di accesso per i servizi software e i protocolli di livello utente.
\end{enumerate}

Nella trasmissione tra due terminali d'estremità (\emph{end systems} o \emph{hosts}) interconnessi tramite una rete geografica o un'infrastruttura di commutazione intermedia, esiste una netta separazione tra i protocolli attivi \emph{end-to-end} e i protocolli attivi \emph{hop-by-hop}.

% [IMMAGINE DEL FILE DI SLIDE CON NOME "1) Ripasso_1"] Modello di riferimento ISO/OSI con attraversamento di router intermedi: livelli 1-3 attivi nei nodi interni di commutazione, livelli 4-7 confinati agli host terminali — Slide 2 / Pagina 2 del PDF

Mentre i livelli superiori (Applicazione, Presentazione, Sessione e Trasporto) sono instanziati esclusivamente nei terminali sorgente e destinazione, i dispositivi intermedi lungo il cammino di transito — tipicamente router e commutatori — implementano unicamente i livelli inferiori dello stack (Fisico, Collegamento e Rete). In tale contesto, un router intermedio opera a livello 3 de-incapsulando la trama ricevuta dal livello 2, ispezionando l'intestazione di rete per consultare la tabella di instradamento (\emph{routing table}), modificando il parametro di durata di vita del pacchetto e re-incapsulando il datagramma all'interno di una nuova trama di livello 2 verso il nodo successivo (\emph{next hop}).

\subsection{Confronto Tra Il Modello OSI E Lo Stack TCP/IP}

Nonostante l'eccellenza concettuale e l'esaustività formale del modello OSI, lo sviluppo concreto della rete Internet ha decretato il primato operativo della famiglia di protocolli TCP/IP (\emph{Transmission Control Protocol/Internet Protocol}).

La complessità prescrittiva dello stack OSI, caratterizzato da un sovraccarico elaborativo (\emph{overhead}) considerevole introdotto dai livelli di Sessione e Presentazione, ne ha limitato l'adozione commerciale. Al contrario, lo stack TCP/IP ha adottato un approccio pragmatico e snello, raggruppando i livelli 5, 6 e 7 dell'OSI in un unico Livello Applicativo, demandando le relative funzionalità di serializzazione e mantenimento dello stato direttamente al codice applicativo o a librerie condivise.

% [IMMAGINE DEL FILE DI SLIDE CON NOME "1) Ripasso_1"] Mappatura concettuale tra il modello OSI a 7 livelli, il modello TCP/IP standard a 4 livelli e il modello TCP/IP equivalente a 5 livelli — Slide 3 / Pagina 3 del PDF

Nel modello TCP/IP standard formulato originariamente da DARPA e descritto nelle RFC dell'IETF, l'architettura si articola in quattro strati:
\begin{itemize}[noitemsep]
    \item \textbf{Application Layer};
    \item \textbf{Transport Layer};
    \item \textbf{Internet Layer};
    \item \textbf{Network Access Layer} (o \emph{Host-to-Network Layer}).
\end{itemize}

Tuttavia, l'accorpamento del livello di collegamento (\emph{Data Link}) e del livello fisico all'interno di un generico strato di accesso alla rete risulta distorsivo rispetto alla reale architettura hardware e software dei calcolatori e delle apparecchiature di rete. Nel silicio e nei driver del sistema operativo, la strutturazione delle trame logiche (gestita da schede NIC e controllori MAC) è nettamente distinta dalla modulazione di tensione, radiofrequenza o impulso ottico (trasduttori fisici PHY). Per tale ragione, nella trattazione ingegneristica e nel proseguo di questo testo si adotta il Modello TCP/IP Equivalente a 5 Livelli, il quale integra il rigore funzionale della separazione fisica/collegamento dell'OSI con la suite di protocolli TCP/IP:
\begin{equation}
    \text{Stack Equivalente} = \{\text{Applicativo},\, \text{Trasporto},\, \text{Rete},\, \text{Collegamento Dati},\, \text{Fisico}\}.
\end{equation}

\subsection{Tassonomia Dei Principali Protocolli Dello Stack TCP/IP}

Ogni livello del modello TCP/IP definisce specifiche entità protocollari standardizzate, progettate per rispondere a precisi vincoli di comunicazione.

% [IMMAGINE DEL FILE DI SLIDE CON NOME "1) Ripasso_1"] Panoramica dei protocolli comuni distribuiti sui livelli dello stack TCP/IP equivalente — Slide 4 / Pagina 4 del PDF

\subsubsection{Livello Applicativo}
Fornisce i protocolli per l'interazione diretta con i servizi utente e i sistemi distribuiti:
\begin{itemize}
    \item \textbf{HTTP} (\emph{Hypertext Transfer Protocol}): protocollo di tipo client-server basato su paradigma richiesta/risposta per il reperimento di risorse sul Web.
    \item \textbf{FTP} (\emph{File Transfer Protocol}): protocollo per il trasferimento bidirezionale di file tra calcolatori, fondato sull'apertura disaccoppiata di due canali logici distinti per controllo (porta TCP 21) e dati (porta TCP 20).
    \item \textbf{TFTP} (\emph{Trivial File Transfer Protocol}): versione elementare di trasferimento file basata su UDP (porta 69), priva di autenticazione e a bassissimo impiego di risorse, frequentemente impiegata per il \emph{bootstrap} di rete di apparati \emph{diskless} o il caricamento di firmware.
    \item \textbf{Telnet}: protocollo di terminale virtuale (porta TCP 23) per l'amministrazione remota in chiaro, oggi surrogato in produzione da SSH.
    \item \textbf{SNMP} (\emph{Simple Network Management Protocol}): protocollo per la telemetria, il monitoraggio dello stato operativo e la configurazione remota dei nodi di rete.
    \item \textbf{SMTP} (\emph{Simple Mail Transfer Protocol}): protocollo orientato al testo per l'inoltro e l'instradamento di messaggi di posta elettronica (porta TCP 25).
    \item \textbf{DNS} (\emph{Domain Name System}): servizio di risoluzione distribuita e gerarchica tra nomi a dominio mnemonici e indirizzi IP numerici (porta UDP/TCP 53).
    \item \textbf{DHCP} (\emph{Dynamic Host Configuration Protocol}): meccanismo di allocazione dinamica dei parametri di rete (indirizzo IP, maschera di sottorete, \emph{default gateway}, server DNS) all'accesso di un nodo al segmento locale.
\end{itemize}

\subsubsection{Livello Di Trasporto}
Realizza la comunicazione logica di processo end-to-end avvalendosi di identificatori noti come porte:
\begin{itemize}
    \item \textbf{TCP} (\emph{Transmission Control Protocol}, RFC 793): protocollo orientato alla connessione (\emph{connection-oriented}) che garantisce consegna affidabile, ordinata e priva di duplicazioni dei flussi di byte, implementando controllo di flusso (\emph{sliding window}) e controllo della congestione.
    \item \textbf{UDP} (\emph{User Datagram Protocol}, RFC 768): protocollo non orientato alla connessione (\emph{connectionless}), privo di meccanismi di riscontro (\emph{acknowledgement}), controllo di sequenza o ritrasmissione. Caratterizzato da bassissimo \emph{overhead} e latenze minime, delega interamente l'eventuale controllo di affidabilità alla logica applicativa.
\end{itemize}

\subsubsection{Livello Di Rete}
Governa l'indirizzamento logico universale e l'instradamento globale:
\begin{itemize}
    \item \textbf{IP} (\emph{Internet Protocol}): protocollo principale dello strato di rete; offre un servizio di trasporto a pacchetto di tipo \emph{best-effort}, non garantendo l'assenza di perdite, l'ordine di consegna o l'assenza di duplicazioni. Opera attraverso le versioni IPv4 e IPv6.
    \item \textbf{ICMP} (\emph{Internet Control Message Protocol}): protocollo di servizio a supporto di IP delegato al trasporto di messaggi di diagnostica, notifica di errori operativi (ad esempio \emph{Destination Unreachable}, \emph{Time Exceeded}) e strumenti di collaudo connettività (comandi \texttt{ping} e \texttt{traceroute}).
    \item \textbf{IGMP} (\emph{Internet Group Management Protocol}): protocollo utilizzato dagli host e dai router locali per gestire e sincronizzare l'adesione dei nodi ai gruppi di trasmissione multicast IPv4.
\end{itemize}

\subsubsection{Livello Collegamento Dati}
Governa il transito delle strutture informative all'interno del medesimo segmento di rete locale o su collegamenti punto-a-punto:
\begin{itemize}
    \item \textbf{Ethernet}: standard dominante per reti locali (LAN) ad accesso multiplo e diffusione (\emph{broadcast}).
    \item \textbf{PPP} (\emph{Point-to-Point Protocol}): protocollo per collegamenti dedicati diretti tra due soli nodi, ampiamente impiegato nelle reti ad estensione geografica (WAN) con supporto nativo all'autenticazione.
    \item \textbf{PPPoE} (\emph{Point-to-Point Protocol over Ethernet}): incapsula trame PPP all'interno di trame Ethernet, consentendo a un fornitore di servizi (ISP) di autenticare e tariffare singoli calcolatori collegati su apparati di commutazione condivisi tipici delle connessioni a banda larga residenziali.
\end{itemize}

\subsection{Meccanismo Di Incapsulamento E Protocol Data Unit (PDU)}

La cooperazione tra i livelli dello stack si concretizza attraverso il principio fondamentale dell'\textbf{incapsulamento}. A ciascun livello gerarchico, l'entità software/hardware locale riceve dal livello immediatamente superiore una porzione informativa denominata SDU (\emph{Service Data Unit}), vi appone la propria intestazione specifica (\emph{header}) ed eventualmente una coda di chiusura (\emph{trailer}), generando così la PDU (\emph{Protocol Data Unit}) del livello corrente.

% [IMMAGINE DEL FILE DI SLIDE CON NOME "1) Ripasso_1"] Diagramma del processo sequenziale di incapsulamento dati alla sorgente: dalla composizione HTTP all'emissione dei bit sul canale fisico — Slide 6 / Pagina 6 del PDF

Si consideri a titolo esemplificativo l'accesso a un sito web mediante browser:
\begin{enumerate}
    \item Livello Applicativo: Il browser genera una richiesta HTTP (\emph{payload} applicativo o Dato).
    \item Livello Di Trasporto: Il modulo TCP incapsula il dato aggiungendo l'intestazione TCP contenente le porte logiche sorgente e destinazione (ad esempio porta di destinazione 80) e i numeri di sequenza, generando il Segmento (\emph{Segment}).
    \item Livello Di Rete: Il modulo IP incapsula il segmento TCP aggiungendo l'intestazione IP contenente gli indirizzi logici sorgente e destinazione, definendo il Pacchetto o Datagramma (\emph{Packet}).
    \item Livello Di Collegamento Dati: Il driver della scheda di rete (NIC) incapsula il pacchetto IP aggiungendo un'intestazione Ethernet con gli indirizzi fisici (MAC) sorgente e destinazione e appende in coda un campo di verifica dell'integrità denominato FCS (\emph{Frame Check Sequence}), costituendo la Trama (\emph{Frame}).
    \item Livello Fisico: La trama binaria viene codificata e convertita in un flusso sequenziale di simboli fisici (impulsi elettrici, modulazioni luminose in fibra ottica o onde elettromagnetiche su canale radio), prendendo il nome di Flusso Di Bit (\emph{Bitstream}).
\end{enumerate}

% [IMMAGINE DEL FILE DI SLIDE CON NOME "1) Ripasso_1"] Mappatura esplicita tra livelli dello stack TCP/IP, terminologia standard delle PDU e campi costitutivi della trama incapsulata — Slide 28 / Pagina 28 del PDF

La tabella~\ref{tab:pdu_layers} sintetizza la corrispondenza biunivoca tra livello, unità informativa ed esempi tipici di protocollo.

\begin{table}[H]
    \centering
    \small % Migliora la resa tipografica per tabelle informative con acronimi
    \renewcommand{\arraystretch}{1.25} % Aumenta l'altezza delle righe per dare respiro al testo
    \caption{Corrispondenza tra livelli gerarchici, denominazione standard delle PDU e protocolli associati.}
    \label{tab:pdu_layers}

    % Tabella estesa al 100% (\linewidth) con allineamento a sinistra (raggedright) nelle celle di testo
    \begin{tabularx}{\linewidth}{l l >{\hsize=0.9\hsize\raggedright\arraybackslash}X >{\hsize=1.1\hsize\raggedright\arraybackslash}X}
        \toprule
        \textbf{Livello Logico} &
        \textbf{Nome PDU} &
        \textbf{Campi Caratterizzanti} &
        \textbf{Esempi Di Protocollo} \\
        \midrule

        Applicativo &
        Dato (\emph{Data}) &
        Metadati applicativi, codifiche MIME &
        HTTP, FTP, DNS, SMTP, Telnet \\

        Trasporto &
        Segmento (\emph{Segment}) &
        Porte logiche, numeri di sequenza/ACK &
        TCP, UDP \\

        Rete &
        Pacchetto (\emph{Packet}) &
        Indirizzi logici IP, TTL, Protocol &
        IPv4, IPv6, ICMP, IGMP \\

        Collegamento &
        Trama (\emph{Frame}) &
        Indirizzi fisici MAC, EtherType, FCS &
        Ethernet (802.3/II), PPP, PPPoE \\

        Fisico &
        Flusso di bit (\emph{Bitstream}) &
        Tensioni, frequenze, transizioni di fase &
        1000BASE-T, Fibra ottica, Wi-Fi \\

        \bottomrule
    \end{tabularx}
\end{table}


\newpage

\section{Tecnologia Ethernet E Livello Di Collegamento Dati}

Il livello di collegamento dati si posiziona come strato mediatore tra il piano logico globale del livello di rete e le peculiarità trasmissive del livello fisico. Esso è primariamente responsabile della frammentazione in trame, dell'indirizzamento fisico su scala locale e del controllo dell'integrità strutturale dei dati trasferiti.

\subsection{Evoluzione Delle Architetture Ethernet: Da Hub Condivisi A Switch Dedicati}

Le reti Ethernet, introdotte storicamente su cavo coassiale condiviso, nascono come architetture basate su canali trasmissivi ad accesso multiplo e diffusione passiva.

Nelle topologie Ethernet primitive (\emph{Early Ethernet}), i nodi erano interconnessi mediante un concentratore fisico noto come hub. L'hub opera esclusivamente a Livello 1 (Fisico), fungendo da semplice ripetitore di segnale elettrico: qualsiasi bit ricevuto da una porta viene amplificato e replicato su tutte le altre porte del dispositivo, senza alcuna elaborazione del contenuto o ispezione delle intestazioni.

% [IMMAGINE DEL FILE DI SLIDE CON NOME "1) Ripasso_1"] Confronto architetturale tra bus condiviso/hub (Early Ethernet) e commutazione dedicata a stella mediante switch (Switch Networking) — Slide 12 / Pagina 12 del PDF

Tale configurazione presenta intrinseci limiti di scalabilità e prestazioni dovuti al fenomeno della contesa del canale. Qualora due o più stazioni tentino di trasmettere contemporaneamente sul mezzo condiviso, i relativi segnali si sovrappongono, corrompendo l'informazione: si verifica una collisione.

L'evoluzione tecnologica ha introdotto lo switch (commutatore) a Livello 2. Lo switch implementa un'architettura di commutazione a pacchetto di tipo \emph{store-and-forward} (o \emph{cut-through}), analizzando l'indirizzo MAC di destinazione della trama per inoltrarla selettivamente sulla sola porta fisica a cui è collegato il nodo ricevente. Di conseguenza, ciascuna porta di uno switch costituisce un circuito di comunicazione indipendente che può operare in modalità bidirezionale simultanea (Full-Duplex), confinando ed eliminando del tutto le collisioni tra host dedicati.

\subsection{Domini Di Collisione E Il Meccanismo CSMA/CD}

Per formalizzare rigorosamente l'impatto topologico delle tecnologie di commutazione, si definiscono due concetti cardine dell'ingegneria di rete:

\begin{theorem}[\textbf{Dominio Di Collisione}]
Un \textbf{dominio di collisione} è definito come l'insieme di tutti i nodi e le interfacce fisiche connessi a un mezzo trasmissivo condiviso tali per cui, qualora due o più entità trasmettano contemporaneamente, i rispettivi segnali elettromagnetici collidono, corrompendosi a vicenda.
\end{theorem}

% [IMMAGINE DEL FILE DI SLIDE CON NOME "1) Ripasso_1"] Distinzione visiva di dominio di collisione singolo su mezzo condiviso rispetto a domini di collisione isolati su ciascuna interfaccia di uno switch — Slide 13 / Pagina 13 del PDF

Su un bus condiviso o su una rete imperniata su hub, l'intera infrastruttura appartiene a un unico dominio di collisione. All'aumentare del numero di nodi concorrenti, la probabilità di collisione cresce rapidamente, degradando drasticamente l'efficienza globale del canale (\emph{throughput} utile).

Per coordinare l'accesso su tali mezzi condivisi, Ethernet adotta l'algoritmo CSMA/CD (\emph{Carrier Sense Multiple Access with Collision Detection}), formalizzato dai seguenti passaggi logici:
\begin{enumerate}
    \item Ascolto preventivo (\emph{Carrier Sense}): Il terminale che intende trasmettere verifica costantemente se la linea condivisa sia occupata. Se la linea è occupata, attende; se la linea risulta libera per un intervallo temporale minimo prestabilito, avvia l'emissione del pacchetto.
    \item Monitoraggio durante la trasmissione (\emph{Listen While Sending}): Anche durante l'emissione, il trasmettitore campiona il mezzo per rilevare eventuali anomalie nei livelli di ampiezza o forma d'onda del segnale.
    \item Rilevamento della collisione e sequenza di disturbo (\emph{Jamming}): Qualora due stazioni rilevino il canale libero e trasmettano simultaneamente, i segnali collidono generando un'onda anomala. Non appena il trasmettitore rileva la collisione, interrompe immediatamente l'invio dei dati e trasmette un segnale di disturbo deliberato (\emph{jam signal} o impulsi di disturbo) per garantire che tutti gli altri nodi rilevino l'avvenuta collisione.
    \item Attesa casuale esponenziale (\emph{Exponential Backoff}): Dopo aver interrotto la trasmissione, la stazione entra in una fase di attesa determinata da un algoritmo di \emph{backoff} casuale esponenziale, al termine del quale riavvia la sequenza d'ascolto prima di ritrasmettere.
\end{enumerate}

Lo switch di rete segmenta fisicamente i domini di collisione: ogni interfaccia dello switch corrisponde a un dominio di collisione indipendente. Se a ogni porta è collegato un singolo host operativo in Full-Duplex, le collisioni vengono del tutto sradicate alla base.

\subsection{Domini Di Broadcast E Indirizzi Di Diffusione}

Mentre lo switch isola i domini di collisione, esso si comporta diversamente nei confronti del traffico di diffusione.

\begin{theorem}[\textbf{Dominio Di Broadcast}]
Un \textbf{dominio di broadcast} è definito come l'ambito topologico massimo entro cui una trama di diffusione (\emph{broadcast frame}) generata da un qualsiasi host viene propagata e recapitata a tutti gli altri nodi appartenenti alla medesima rete logica di livello 2.
\end{theorem}

% [IMMAGINE DEL FILE DI SLIDE CON NOME "1) Ripasso_1"] Comportamento di propagazione del frame di broadcast sia in ambiente hub sia attraverso le interfacce di switch appartenenti al medesimo dominio di broadcast L2 — Slide 14 / Pagina 14 del PDF

Una trama Ethernet indirizzata all'indirizzo fisico speciale con tutti i bit posti a uno (\texttt{FF-FF-FF-FF-FF-FF}) deve essere ricevuta e processata da tutte le schede di rete della rete locale.

Uno switch Ethernet privo di configurazioni VLAN propaga (\emph{flooding}) le trame di broadcast su tutte le porte attive, ad eccezione della sola interfaccia di ingresso dalla quale la trama è giunta. Di conseguenza, tutti i nodi attestati su uno switch (o su più switch interconnessi in cascata) condividono il medesimo dominio di broadcast.

Un traffico di broadcast incontrollato (generato ad esempio da continue richieste ARP o tempeste di broadcast indotte da anelli di rete) consuma risorse trasmissive e costringe le CPU di tutti gli host terminali a interrompere le proprie attività per esaminare e scartare trame irrilevanti. L'interruzione e il contenimento dei domini di broadcast richiedono l'impiego di dispositivi che operano a Livello 3 (router o switch Layer 3) o la partizione logica mediante VLAN (\emph{Virtual Local Area Network}).

\subsection{Formati Delle Trame: Ethernet II Ed IEEE 802.3 LLC/SNAP}

Nel panorama delle reti locali convivono storicamente due formati primari per la strutturazione della trama Ethernet: il formato industriale Ethernet II (sviluppato dal consorzio DEC-Intel-Xerox, o DIX) e lo standard internazionale IEEE 802.3.

% [IMMAGINE DEL FILE DI SLIDE CON NOME "1) Ripasso_1"] Struttura comparativa dettagliata dei campi costitutivi tra trame Ethernet II e trame IEEE 802.3 LLC/SNAP — Slide 15 / Pagina 15 del PDF

\subsubsection{Formato Ethernet II}
Rappresenta la stragrande maggioranza del traffico sulle moderne reti TCP/IP. La trama presenta la seguente organizzazione:
\begin{itemize}[noitemsep]
    \item \textbf{D.MAC} (6 byte): indirizzo MAC di destinazione (\emph{Destination MAC Address});
    \item \textbf{S.MAC} (6 byte): indirizzo MAC di sorgente (\emph{Source MAC Address});
    \item \textbf{Type} (2 byte): campo EtherType, indicante il protocollo di livello 3 a cui inoltrare il carico utile (\emph{payload}). Esempi: \texttt{0x0800} per IPv4, \texttt{0x0806} per ARP, \texttt{0x8137} per Novell IPX, \texttt{0x809B} per AppleTalk;
    \item \textbf{User Data} (da 46 a 1500 byte): carico utile proveniente dallo strato di rete;
    \item \textbf{FCS} (4 byte): sequenza di controllo di trama basata su codice di ridondanza ciclica a 32 bit (\textbf{CRC-32}).
\end{itemize}

\subsubsection{Formato IEEE 802.3 LLC/SNAP}
Lo standard IEEE 802.3 ha originariamente sostituito il campo EtherType con un campo Length (2 byte), indicante la lunghezza effettiva dei dati trasportati. Per supportare il demultiplexing dei protocolli, è stato introdotto il sottolivello LLC (\emph{Logical Link Control}, IEEE 802.2), composto da:
\begin{itemize}[noitemsep]
    \item \textbf{DSAP} (1 byte): \emph{Destination Service Access Point}; vale ad esempio \texttt{0x06} nel caso di incapsulamento IP;
    \item \textbf{SSAP} (1 byte): \emph{Source Service Access Point};
    \item \textbf{Ctrl} (1 byte): campo di controllo, tipicamente configurato a \texttt{0x03} per indicare frame informativo non numerato (\emph{unnumbered}) in modalità non orientata alla connessione.
\end{itemize}
Per ovviare all'esiguità dello spazio di indirizzamento a 8 bit dei SAP e mantenere compatibilità retroattiva con i codici DIX, lo standard prevede l'estensione opzionale SNAP (\emph{Subnetwork Access Protocol}, 5 byte):
\begin{itemize}[noitemsep]
    \item \textbf{Org Code} (3 byte): codice identificativo dell'organizzazione fornitrice (impostato a tutti zero per compatibilità);
    \item \textbf{Type} (2 byte): esprime il medesimo protocollo specificato nel campo EtherType della trama Ethernet II.
\end{itemize}

\subsubsection{Criterio Di Discriminazione Tra Ethernet II Ed IEEE 802.3}
Poiché la posizione a 12 byte dall'inizio della trama (subito dopo i due indirizzi MAC) è occupata dal campo \emph{Type} in Ethernet II e dal campo \emph{Length} in 802.3, i ricevitori adottano una semplice regola numerica di demultiplazione:
\begin{equation}
    \begin{cases}
        \text{Valore del campo} \le 1500 \; (\texttt{0x05DC}) & \implies \text{Trama interpretata come IEEE 802.3 (Length)} \\
        \text{Valore del campo} \ge 1536 \; (\texttt{0x0600}) & \implies \text{Trama interpretata come Ethernet II (Type)}
    \end{cases}
\end{equation}
L'intervallo compreso tra 1501 e 1535 è formalmente riservato ed escluso dalle specifiche standard.

\subsection{Vincoli Dimensionali Della Trama: Limiti Minimi, Massimi Ed Interpacket Gap}

La lunghezza complessiva di una trama Ethernet II (escludendo il preambolo fisico di 8 byte composto da bit 0 e 1 alternati e dal delimitatore di inizio trama SFD) è rigorosamente vincolata tra 64 byte e 1518 byte.

\subsubsection{Motivazione Della Lunghezza Minima (64 Byte)}
Il limite inferiore di 64 byte è imposto dalla relazione fisica esistente tra la massima distanza geografica coperta dal cavo coassiale condiviso, la velocità di propagazione del segnale e il funzionamento di CSMA/CD.

Affinché un nodo trasmettitore $A$ possa accorgersi con certezza di una collisione causata da un nodo remoto $B$ posizionato all'estremo opposto del segmento, la trasmissione della trama da parte di $A$ deve prolungarsi per un tempo non inferiore al tempo di andata e ritorno (\textbf{RTT}, \emph{Round Trip Time} o \emph{Slot Time}) necessario affinché il segnale raggiunga $B$ e l'anomalia di tensione ritorni verso $A$:
\begin{equation}
    T_{\text{tx}} \ge 2 \cdot \tau_{\text{prop}}, \qquad \tau_{\text{prop}} = \frac{d_{\max}}{v}.
\end{equation}
Se $A$ terminasse di trasmettere prima che il segnale di collisione facesse ritorno, $A$ riterrebbe erroneamente la trasmissione completata con successo e libererebbe il buffer, ignorando che la trama è andata distrutta.

Poiché le intestazioni della trama (6 byte D.MAC + 6 byte S.MAC + 2 byte Type) unitamente alla coda FCS (4 byte) occupano un totale fisso di $14 + 4 = 18\text{ byte}$, il carico utile di livello 3 (\emph{payload}) deve comprendere un quantitativo minimo di dati pari a:
\begin{equation}
    \text{Carico utile minimo} = 64 - 18 = 46\text{ byte}.
\end{equation}
Se il pacchetto IP generato dal livello di rete ha una dimensione inferiore a 46 byte, lo strato di collegamento provvede all'inserimento di bit di riempimento (\textbf{Padding}) per saturare la trama fino alla soglia dei 64 byte.

\subsubsection{Motivazione Della Lunghezza Massima (1518 Byte) E MTU}
All'estremo opposto, la dimensione massima di 1518 byte — corrispondente a una MTU (\emph{Maximum Transmission Unit}) a livello IP di 1500 byte — rappresenta un compromesso ingegneristico tra:
\begin{itemize}
    \item Efficienza trasmissiva: un pacchetto di grandi dimensioni riduce l'incidenza percentuale dell'\emph{overhead} dovuto alle intestazioni;
    \item Equità e latenza: su un mezzo condiviso a commutazione di pacchetto, trame eccessivamente lunghe monopolizzerebbero il canale monopolizzando i buffer di commutazione e introducendo un degrado inaccettabile (\emph{jitter} e ritardo di accodamento) per il traffico interattivo e sensibile alla latenza;
    \item Probabilità di errore: all'aumentare della lunghezza di un pacchetto, aumenta la probabilità che un singolo bit corrotto richieda la ritrasmissione dell'intero blocco.
\end{itemize}

\subsubsection{Interpacket Gap (IPG)}
Tra l'emissione di due trame consecutive, il livello fisico impone un intervallo temporale minimo di inattività del mezzo noto come Interpacket Gap (IPG). Nello standard Ethernet a 10 e 100 Mbps, l'IPG corrisponde ad almeno 96 tempi di bit (equivalenti a 12 byte o ottetti). Questo intervallo consente ai circuiti di ricezione delle schede di rete di recuperare lo stato di riposo, svuotare i registri interni di temporizzazione e prepararsi a sincronizzarsi sul preambolo della trama successiva.

\newpage

\section{Indirizzamento Fisico: MAC Address Ed Inoltro}

A differenza degli indirizzi logici IP, che riflettono la posizione topologica del nodo all'interno di una determinata sottorete gerarchica, gli indirizzi fisici identificano in modo permanente e univoco l'interfaccia hardware (NIC).

\subsection{Struttura Dell'Indirizzo MAC: OUI E CID}

Un indirizzo MAC (\emph{Media Access Control}) è composto da 48 bit (6 byte), convenzionalmente rappresentato mediante 12 cifre esadecimali raggruppate a coppie separate da trattini (\texttt{00-1E-10-DD-DD-02}), due punti (\texttt{00:1E:10:DD:DD:02}) o a gruppi di quattro cifre separate da punti (\texttt{001E-10DD-DD02}).

% [IMMAGINE DEL FILE DI SLIDE CON NOME "1) Ripasso_1"] Struttura dell'indirizzo MAC a 48 bit con ripartizione esadecimale e binaria tra OUI e identificativo fornitore — Slide 17 / Pagina 17 del PDF

La codifica a 48 bit si struttura in due campi da 24 bit ciascuno:
\begin{enumerate}
    \item OUI (\emph{Organizationally Unique Identifier}): i primi 24 bit (3 byte) rappresentano un codice univoco globale assegnato dall'ente IEEE a ciascun produttore hardware autorizzato.
    \item CID (\emph{Company ID} o identificativo del dispositivo): i restanti 24 bit (3 byte) vengono assegnati dal produttore in maniera sequenziale per garantire che ogni singola scheda di rete fabbricata possieda un identificatore univoco al mondo.
\end{enumerate}

% [IMMAGINE DEL FILE DI SLIDE CON NOME "1) Ripasso_1"] Classificazione tassonomica degli indirizzi MAC (Unicast, Multicast, Broadcast) e posizione del bit I/G — Slide 18 / Pagina 18 del PDF

\subsection{Classificazione Degli Indirizzi: Unicast, Multicast E Broadcast}

La destinazione operativa di una trama Ethernet è determinata dal valore assunto dai bit dell'indirizzo MAC, con particolare rilevanza per l'ottavo bit del primo byte (noto come bit I/G, \emph{Individual/Group}):

% [IMMAGINE DEL FILE DI SLIDE CON NOME "1) Ripasso_1"] Flusso di trasmissione di una trama Unicast da host sorgente ad host destinatario singolo su topologia a stella — Slide 20 / Pagina 20 del PDF

\begin{itemize}
    \item \textbf{Indirizzo MAC Unicast (Individual)}: Corrisponde a un singolo dispositivo fisico. Il bit I/G (l'ottavo bit del primo ottetto) è posto rigidamente a 0 (ad esempio: \texttt{00-1E-10-DD-DD-02}, dove \texttt{00}$_{16} = 0000\,000\mathbf{0}_2$). Può fungere sia da indirizzo sorgente che da indirizzo di destinazione. Due dispositivi con il medesimo indirizzo MAC unicast sulla stessa rete L2 causano un conflitto catastrofico che inibisce la corretta commutazione del traffico.

    % [IMMAGINE DEL FILE DI SLIDE CON NOME "1) Ripasso_1"] Flusso di trasmissione di una trama Broadcast verso l'indirizzo all-ones con ricezione obbligatoria da parte di tutti i nodi — Slide 21 / Pagina 21 del PDF

    \item \textbf{Indirizzo MAC Broadcast}: Indirizzo speciale composto da 48 bit tutti posti a uno, la cui notazione esadecimale è \texttt{FF-FF-FF-FF-FF-FF}. È inteso per recapitare il pacchetto indistintamente a tutti i terminali attivi sul segmento di rete. Può essere utilizzato solo come indirizzo di destinazione.

    % [IMMAGINE DEL FILE DI SLIDE CON NOME "1) Ripasso_1"] Flusso di trasmissione di una trama Multicast diretta a un gruppo specifico di nodi in ascolto — Slide 22 / Pagina 22 del PDF

    \item \textbf{Indirizzo MAC Multicast (Group)}: Indirizzato a un gruppo selezionato di stazioni in ascolto per uno specifico servizio o protocollo. È identificato dal bit I/G posto a 1 (ad esempio: \texttt{01-80-C2-00-00-01}, dove \texttt{01}$_{16} = 0000\,000\mathbf{1}_2$). Non può mai essere impiegato come indirizzo sorgente di una trama.
\end{itemize}

\subsection{Commutazione A Livello 2 E Gestione Della Tabella MAC}

A differenza di un hub, che ripete ciecamente i segnali, uno switch di Livello 2 esegue l'inoltro intelligente delle trame analizzando l'intestazione Ethernet.

% [IMMAGINE DEL FILE DI SLIDE CON NOME "1) Ripasso_1"] Posizionamento architetturale dei commutatori di Livello 2 e commutatori di Livello 3 all'interno della gerarchia di rete di campus — Slide 23 / Pagina 23 del PDF

Per operare, il commutatore mantiene in memoria volatile una tabella degli indirizzi MAC (\emph{MAC Address Table} o memoria CAM), la quale associa a ciascun indirizzo MAC noto la specifica porta fisica dello switch su cui tale dispositivo risulta connesso.

% [IMMAGINE DEL FILE DI SLIDE CON NOME "1) Ripasso_1"] Struttura della tabella di corrispondenza MAC-interfaccia fisica all'interno di uno switch — Slide 24 / Pagina 24 del PDF

L'algoritmo di apprendimento e inoltro (\emph{Learning and Forwarding}) è governato dalla logica descritta nell'Algoritmo~\ref{alg:switch_forwarding}.

\begin{algorithm}[H]
    \caption{Apprendimento e Inoltro di Livello 2 nello Switch}
    \label{alg:switch_forwarding}
    \Dati{Trama Ethernet in ingresso con indirizzo sorgente $\text{MAC}_S$, destinazione $\text{MAC}_D$, su interfaccia $I_{\text{in}}$}
    \Risultato{Inoltro della trama sulla o sulle interfacce di uscita}
    \Inizializza\;
    \Commento{1. Fase di Apprendimento (Source Learning)}
    \Se{$\text{MAC}_S$ non appartiene alla tabella MAC oppure la porta associata è diversa da $I_{\text{in}}$}{
        Aggiorna la tabella con la voce $(\text{MAC}_S \to I_{\text{in}})$\;
        Inizializza il timer di invecchiamento (\emph{aging timer}) per tale voce\;
    }
    \Commento{2. Fase di Inoltro (Destination Forwarding)}
    \Se{$\text{MAC}_D == \text{FF-FF-FF-FF-FF-FF}$ (\emph{Broadcast}) oppure $\text{MAC}_D$ è Multicast}{
        Esegui il \textbf{Flooding}: trasmetti la trama su tutte le porte tranne $I_{\text{in}}$\;
    }
    \Altrimenti{
        \Se{$\text{MAC}_D$ è presente nella tabella MAC con porta $I_{\text{out}}$}{
            \Se{$I_{\text{out}} \ne I_{\text{in}}$}{
                Inoltra la trama selettivamente su $I_{\text{out}}$ (\textbf{Unicast Forwarding})\;
            }
            \Altrimenti{
                Scarta la trama (\emph{Filtering}), risiede sullo stesso segmento\;
            }
        }
        \Altrimenti{
            Esegui il \textbf{Flooding} della trama (\emph{Unknown Unicast Flooding})\;
        }
    }
\end{algorithm}

% [IMMAGINE DEL FILE DI SLIDE CON NOME "1) Ripasso_1"] Stato iniziale con tabella MAC vuota prima dell'inizio delle comunicazioni tra host — Slide 25 / Pagina 25 del PDF

All'accensione dello switch, la tabella è vuota. Quando l'Host 1 invia una prima trama verso l'Host 2, la destinazione risulta sconosciuta: lo switch memorizza la porta di ingresso associandola al MAC dell'Host 1 e replica la trama su tutte le porte (\emph{unknown unicast flooding}). Quando l'Host 2 risponde, lo switch apprende anche la sua porta, trasformando ogni scambio successivo tra i due nodi in una commutazione punto-a-punto diretta.

% [IMMAGINE DEL FILE DI SLIDE CON NOME "1) Ripasso_1"] Meccanismo di unknown unicast flooding e contestuale apprendimento della porta sorgente — Slide 26 / Pagina 26 del PDF

Per evitare la saturazione della tabella e garantire l'adattamento ai cambiamenti topologici (ad esempio lo spostamento fisico di un cavo o la disconnessione di una macchina), le voci apprese dinamicamente hanno una validità limitata da un tempo di decadimento (\textbf{Aging Time}). Se una voce non viene rinfrescata da nuove trame in ingresso entro tale lasso temporale (il cui valore predefinito negli apparati aziendali è tipicamente di 300 secondi), essa viene eliminata dalla memoria.

\newpage

\section{Il Livello Di Rete Ed Il Protocollo IPv4}

Mentre il livello di collegamento garantisce la consegna delle trame entro i confini dello stesso segmento locale, il Livello Di Rete presiede all'interconnessione globale eterogenea, astraendo le peculiarità trasmissive sottostanti mediante un'infrastruttura di indirizzamento universale.

% [IMMAGINE DEL FILE DI SLIDE CON NOME "1) Ripasso_1"] Schema gerarchico del livello di rete e protocolli ausiliari (IP, ICMP, IPX, IGMP) — Slide 27 / Pagina 27 del PDF

\subsection{Funzioni Cardinali Del Livello Di Rete}

Le responsabilità primarie dello strato di rete includono:
\begin{itemize}
    \item Indirizzamento logico gerarchico: assegnazione di identificatori numerici strutturati univoci (\emph{network ID} e \emph{host ID}) svincolati dal costruttore hardware, consentendo l'aggregazione delle rotte;
    \item Instradamento (\emph{Routing}): selezione computata del cammino ottimale per convogliare un pacchetto dalla sorgente alla destinazione remota attraverso molteplici nodi di commutazione;
    \item Commutazione e inoltro (\emph{Forwarding}): trasferimento interno dei pacchetti dall'interfaccia di ingresso all'opportuna interfaccia di uscita del router sulla scorta dell'ispezione della tabella di instradamento;
    \item Frammentazione e riassemblaggio: frammentazione controllata del datagramma per adeguarlo a collegamenti caratterizzati da una MTU inferiore rispetto a quella del nodo sorgente.
\end{itemize}

\subsection{Formato Dell'Intestazione IPv4}

Il protocollo base della rete Internet è IPv4 (\emph{Internet Protocol Version 4}). Un pacchetto IPv4 è strutturato in una sezione di intestazione (\emph{header}) e una sezione dati (\emph{payload}). L'intestazione ha una dimensione minima prefissata di 20 byte, estensibile fino a un massimo di 60 byte mediante campi opzionali.

% [IMMAGINE DEL FILE DI SLIDE CON NOME "1) Ripasso_1"] Diagramma dettagliato della struttura a 32 bit dei campi costitutivi dell'intestazione IPv4 — Slide 29 / Pagina 29 del PDF

L'intestazione si sviluppa su parole da 32 bit e comprende i campi analiticamente descritti di seguito:

\begin{itemize}
    \item \textbf{Version} (4 bit): specifica la versione del protocollo impiegato. Assume il valore \texttt{4} (\texttt{0100}$_2$) per datagrammi IPv4 e \texttt{6} (\texttt{0110}$_2$) per IPv6.
    \item \textbf{Header Length (IHL)} (4 bit): esprime la lunghezza totale dell'intestazione misurata in parole da 32 bit (blocchi da 4 byte). Poiché la dimensione minima dell'header IPv4 privo di opzioni è di 20 byte, il valore minimo di IHL è pari a:
    \begin{equation}
        \text{IHL}_{\min} = \frac{20\text{ byte}}{4\text{ byte/parola}} = 5 \quad (\texttt{0101}_2).
    \end{equation}
    Il valore massimo rappresentabile su 4 bit è $15$ ($1111_2$), che determina una lunghezza massima dell'header pari a:
    \begin{equation}
        \text{Header}_{\max} = 15 \times 4\text{ byte} = 60\text{ byte}.
    \end{equation}
    \item \textbf{Type of Service / DSCP} (8 bit): originariamente concepito per indicare requisiti di priorità, ritardo e affidabilità, è oggi riconfigurato per supportare i modelli di qualità del servizio DiffServ (\emph{Differentiated Services}), dedicando 6 bit al \emph{Differentiated Services Code Point} (DSCP) per la classificazione del traffico e 2 bit all'\emph{Explicit Congestion Notification} (ECN).
    \item \textbf{Total Length} (16 bit): specifica la lunghezza complessiva dell'intero datagramma IP (intestazione più carico utile applicativo), espressa rigorosamente in byte. Essendo un campo a 16 bit, la dimensione teorica massima consentita per un singolo pacchetto IPv4 risulta pari a:
    \begin{equation}
        \text{Total Length}_{\max} = 2^{16} - 1 = 65535\text{ byte}.
    \end{equation}
    \item \textbf{Identification} (16 bit): identificatore univoco generato dall'host sorgente, assegnato a tutti i frammenti derivati dal medesimo datagramma originale per permetterne il corretto riassemblaggio alla destinazione.
    \item \textbf{Flags} (3 bit): campo di controllo per la gestione della frammentazione, composto da tre flag binari:
    \begin{itemize}[noitemsep]
        \item Bit 0: Riservato, obbligatoriamente impostato a 0;
        \item Bit 1: DF (\emph{Don't Fragment}). Se posto a 1, proibisce ai router intermedi di frammentare il pacchetto; qualora il pacchetto ecceda la MTU del link uscente, viene scartato e il router invia un messaggio di errore ICMP;
        \item Bit 2: MF (\emph{More Fragments}). Se posto a 1, indica che il frammento corrente non è l'ultimo; se posto a 0, specifica che il frammento è l'ultimo della serie o che il pacchetto non è stato frammentato.
    \end{itemize}
    \item \textbf{Fragment Offset} (13 bit)\footnote{Sebbene in alcuni documenti sintetici venga indicata un'ampiezza di 12 bit per via di refusi redazionali, lo standard RFC 791 prescrive in modo non ambiguo l'ampiezza di 13 bit.}: specifica la posizione relativa dei dati trasportati dal frammento corrente rispetto al datagramma iniziale non frammentato. L'offset è misurato in unità di 8 byte (blocchi da 64 bit), vincolando la dimensione di ogni frammento intermedio a un multiplo esatto di 8 byte.
    \item \textbf{Time to Live (TTL)} (8 bit): contatore intero decremenatato di almeno un'unità da ciascun router attraversato lungo la catena di instradamento. Qualora il valore raggiunga lo zero ($\text{TTL} = 0$), il pacchetto viene scartato e viene generato un messaggio ICMP \emph{Time Exceeded} verso la sorgente, impedendo la circolazione indefinita di pacchetti orfani in presenza di anelli di instradamento (\emph{routing loops}).
    \item \textbf{Protocol} (8 bit): identifica il protocollo di livello superiore a cui il payload deve essere consegnato al termine della decapsulazione. I valori standard più diffusi includono:
    \begin{itemize}[noitemsep]
        \item \texttt{1}: ICMP (\emph{Internet Control Message Protocol});
        \item \texttt{2}: IGMP (\emph{Internet Group Management Protocol});
        \item \texttt{6}: TCP (\emph{Transmission Control Protocol});
        \item \texttt{17}: UDP (\emph{User Datagram Protocol}).
    \end{itemize}
    \item \textbf{Header Checksum} (16 bit): codice di rilevazione d'errore calcolato esclusivamente sui campi dell'intestazione IPv4 (escludendo il carico utile). Viene ricalcolato obbligatoriamente a ogni salto di instradamento a causa della continua variazione del campo TTL.
    \item \textbf{Source IP Address} (32 bit): indirizzo logico del nodo trasmettitore originario.
    \item \textbf{Destination IP Address} (32 bit): indirizzo logico del nodo finale a cui il pacchetto è destinato.
    \item \textbf{Options} (dimensione variabile, da 0 a 40 byte): parametri facoltativi utilizzati per diagnostica, sicurezza, tracciamento del cammino (\emph{Record Route}) o instradamento esplicito da sorgente (\emph{Source Routing}).
    \item \textbf{Padding} (variabile): sequenza di bit posti a zero inseriti in coda alle opzioni al solo scopo di allineare la lunghezza totale dell'header su un multiplo intero di 32 bit (4 byte).
\end{itemize}

\newpage

\section{Quadro Comparativo Dei Protocolli Di Trasporto}

A completamento dell'infrastruttura di comunicazione, il Livello Di Trasporto si incarica di colmare il divario funzionale intercorrente tra il servizio di inoltro a pacchetto non garantito fornito da IP a livello di rete e le eterogenee necessità di affidabilità richieste dalle applicazioni distribuite.

% [IMMAGINE DEL FILE DI SLIDE CON NOME "1) Ripasso_1"] Rappresentazione concettuale del livello di trasporto orientato alle porte e connessioni end-to-end — Slide 8 / Pagina 8 del PDF

\subsection{Confronto Architetturale: TCP Versus UDP}

I due protocolli fondamentali dello strato di trasporto riflettono filosofie progettuali radicalmente disaccoppiate:

\begin{itemize}
    \item \textbf{TCP (\emph{Transmission Control Protocol})}: Standardizzato nella RFC 793, è un protocollo rigorosamente orientato alla connessione (\emph{connection-oriented}). Prima che qualunque scambio di dati abbia luogo, i due endpoint stabiliscono un canale logico bidirezionale mediante la sequenza a tre vie nota come Three-Way Handshake (\texttt{SYN} $\to$ \texttt{SYN-ACK} $\to$ \texttt{ACK}). TCP garantisce:
    \begin{itemize}[noitemsep]
        \item Trasferimento ordinato e integro dei dati mediante numeri di sequenza progressivi e riscontri (\emph{acknowledgements});
        \item Rilevazione e ritrasmissione automatica dei segmenti perduti o corrotti;
        \item Controllo di flusso basato su finestra scorrevole (\emph{sliding window}) per non saturare il ricevitore;
        \item Controllo di congestione dinamico per non sovraccaricare i router di transito.
    \end{itemize}
    Tali proprietà introducono un costo in termini di latenza addizionale e un \emph{overhead} d'intestazione minimo di 20 byte per ciascun segmento.

    \item \textbf{UDP (\emph{User Datagram Protocol})}: Standardizzato nella RFC 768, è un protocollo non orientato alla connessione (\emph{connectionless}). Non richiede alcuna fase preliminare di inizializzazione della sessione e non mantiene traccia dello stato interno. I pacchetti (datagrammi) vengono inviati in modo indipendente:
    \begin{itemize}[noitemsep]
        \item Nessuna garanzia di consegna né di preservazione dell'ordine temporale;
        \item Nessun meccanismo integrato di controllo di flusso o congestione;
        \item Intestazione estremamente compatta di soli 8 byte.
    \end{itemize}
    UDP costituisce la scelta preferenziale per servizi interattivi in tempo reale (streaming audio/video, videoconferenza, gaming online, telefonia VoIP) in cui la tempestività prevale sulla ritrasmissione tardiva del dato, o per protocolli transazionali rapidi (DNS, SNMP, DHCP, TFTP).
\end{itemize}

La tabella~\ref{tab:tcp_udp_comparison} formalizza il confronto sinottico tra le due filosofie trasmissive.

\begin{table}[H]
    \centering
    \caption{Confronto esaustivo tra le caratteristiche operative dei protocolli TCP e UDP.}
    \label{tab:tcp_udp_comparison}
    \begin{tabularx}{\textwidth}{l X X}
        \toprule
        \textbf{Caratteristica} & \textbf{TCP (RFC 793)} & \textbf{UDP (RFC 768)} \\
        \midrule
        Orientamento alla connessione & Sì (richiede Three-Way Handshake) & No (emissione immediata senza negoziazione) \\
        Affidabilità di consegna & Garantita (ritrasmissioni con timeout) & \emph{Best-effort} (nessuna ritrasmissione nativa) \\
        Ordinamento dei dati & Garantito tramite numeri di sequenza & Non garantito (possibile ricezione disordinata) \\
        Controllo di flusso & Presente (finestra dinamica del ricevitore) & Assente \\
        Controllo di congestione & Presente (algoritmi \emph{Slow Start}, \emph{AIMD}) & Assente \\
        Dimensione intestazione & Minimo 20 byte (fino a 60 byte con opzioni) & Fisso a 8 byte \\
        Comunicazione supportata & Esclusivamente punto-a-punto (Unicast) & Unicast, Broadcast e Multicast \\
        Applicazioni tipiche & HTTP, HTTPS, FTP, SSH, SMTP & DNS, DHCP, TFTP, RTP/RTCP, VoIP \\
        \bottomrule
    \end{tabularx}
\end{table}

L'assimilazione puntuale di questi concetti fondativi — incapsulamento, discriminazione dei formati Ethernet, demultiplazione a livelli mediante porte e campi protocollo, topologia a commutazione L2 e instradamento L3 — costituisce il prerequisito indispensabile per comprendere le architetture evolute di virtualizzazione di rete, l'incapsulamento \emph{overlay} (quale VxLAN) e i meccanismi di disaccoppiamento tra piano dati e piano di controllo tipici del paradigma \emph{Software-Defined Networking} (SDN).

% AA ---------------------------------------

\section{Architettura Dell'Intestazione IPv4 E Meccanismi Di Trasporto}

\subsection{Richiamo Fondamentale: Il Livello Di Rete Nel Modello A Strati}
Nel modello architetturale di riferimento TCP/IP, il livello di rete (\emph{Network Layer} o \emph{Internet Layer}) ha il compito primario di fornire un servizio di consegna dei dati \emph{best-effort}, non orientato alla connessione e non affidabile, tra host terminali potenzialmente situati su reti geograficamente eterogenee e topologicamente distanti. L'unità informativa fondamentale scambiata a questo strato è il datagramma IP (o pacchetto IP).

A differenza del livello di collegamento (\emph{Data Link Layer}), che gestisce la trasmissione di trame (frame) all'interno di un singolo segmento fisico o logico condiviso (dominio di collisione o di broadcast L2), il livello di rete si occupa dell'instradamento (\emph{routing}) multi-hop. Per consentire ai commutatori di livello 3 (router) di analizzare, frammentare, instradare e recapitare correttamente il traffico, il protocollo IPv4 antepone al carico utile (\emph{payload}, tipicamente fornito dai protocolli di livello di trasporto quali TCP o UDP) un'intestazione strutturata (\emph{IP header}), i cui campi governano la semantica della trasmissione.

\subsection{Frammentazione Del Datagramma IP (\emph{Packet Fragmentation})}
Uno dei compiti critici svolti dal livello di rete risiede nella gestione delle disomogeneità fisiche dei mezzi trasmissivi attraversati dal pacchetto lungo il percorso da sorgente a destinazione. Ogni tecnologia di livello di collegamento stabilisce un limite superiore alla dimensione massima della trama dati che può essere fisicamente trasportata; tale parametro prende il nome di Maximum Transmission Unit (MTU). Ad esempio, la tecnologia Ethernet standard impone un vincolo di MTU pari a \SI{1500}{\byte}.

Quando la dimensione complessiva di un pacchetto IP (comprendente intestazione e carico utile) supera l'MTU del collegamento di rete successivo lungo il cammino di instradamento, il router non è in grado di inoltrarlo come blocco unitario. In assenza di meccanismi alternativi, il pacchetto verrebbe scartato. Per evitare tale perdita, il protocollo IPv4 integra nativamente la funzione di frammentazione: il datagramma originale viene suddiviso in due o più frammenti autonomi di dimensioni inferiori o uguali all'MTU della tratta da impegnare, ciascuno dotato di una propria intestazione IP opportunamente configurata per consentirne il futuro riassemblaggio.

% [IMMAGINE DEL FILE DI SLIDE CON NOME "1) Ripasso_2"] Schema del pacchetto IP con focus sui campi di frammentazione (Identification, Flags, Fragment Offset) e diagramma di Host A che frammenta un pacchetto verso Host B — Slide 1 / Pagina 1 del PDF

La ricostruzione del datagramma originario non avviene solitamente nodo per nodo (il che imporrebbe un pesante overhead elaborativo e di memorizzazione sui router di transito), bensì esclusivamente presso l'host di destinazione finale (\emph{end-to-end reassembly}). Per governare in modo deterministico questo meccanismo, l'intestazione IPv4 impiega tre campi dedicati:
\begin{enumerate}
    \item \textbf{Identification} (16 bit): questo campo contiene un valore numerico progressivo o pseudocasuale generato dall'host mittente per ogni pacchetto originario. Tutti i frammenti generati a partire dal medesimo pacchetto ereditano lo stesso valore di Identification. In fase di ricezione, la destinazione sfrutta questo identificativo, congiuntamente all'indirizzo IP sorgente, per discriminare e raggruppare i frammenti appartenenti alla medesima unità informativa.
    \item \textbf{Flags} (3 bit): campo di controllo formato da tre flag logici:
    \begin{itemize}
        \item \emph{Reserved bit} (Bit 0): bit riservato per sviluppi futuri, vincolato per specifica allo zero logico ($0$).
        \item \emph{Don't Fragment} (DF, Bit 1): flag che controlla l'autorizzazione alla frammentazione. Se impostato a $1$, indica esplicitamente che il pacchetto non deve essere frammentato da alcun router di transito; se tale pacchetto incontra un link con MTU inferiore alla propria dimensione, il router è costretto a scartarlo e a notificare l'errore al mittente mediante un apposito messaggio ICMP. Se impostato a $0$, la frammentazione è formalmente autorizzata.
        \item \emph{More Fragment} (MF, Bit 2): flag indicatore di continuazione. Il valore $1$ segnala che il frammento corrente non è l'ultimo del flusso, indicando alla destinazione che deve attendere ulteriori segmenti; il valore $0$ contraddistingue in modo univoco l'ultimo frammento della catena (o un pacchetto non frammentato).
    \end{itemize}
    \item \textbf{Fragment Offset} (13 bit da standard RFC 791, talvolta schematizzato didatticamente come campo a 12/13 bit a seconda del tracciamento): definisce la posizione relativa (\emph{offset}) del frammento di dati corrente all'interno del payload del pacchetto originale non frammentato. Per massimizzare l'intervallo rappresentabile con un numero ridotto di bit, l'offset non misura i singoli byte, ma opera in blocchi elementari di 8 byte (64 bit). Pertanto, il valore numerico inserito nell'header deve essere moltiplicato per 8 per risalire alla reale posizione in byte:
    \[
    \text{Posizione del byte iniziale} = \text{Fragment Offset} \times 8
    \]
    La combinazione coordinata del campo \emph{Fragment Offset} e del bit \emph{More Fragment} permette al ricevitore di ricomporre il payload originario collocando ogni porzione di dati nella corretta sequenza ordinata, rilevando tempestivamente eventuali buchi dovuti a perdite di frammenti.
\end{enumerate}

\subsection{Il Meccanismo Di Time To Live (TTL)}
In un'infrastruttura di rete a commutazione di pacchetto, anomalie transitorie di convergenza nei protocolli di routing (ad esempio durante l'esecuzione dell'algoritmo di Bellman-Ford in protocolli \emph{distance-vector} o per errori umani nella definizione di rotte statiche) possono dar luogo a anelli di instradamento (\emph{routing loops}). In una situazione di loop, due o più router instradano ciclicamente lo stesso pacchetto l'uno verso l'altro. Senza un vincolo esplicito sulla vita del pacchetto, esso continuerebbe a orbitare indefinitamente, consumando banda trasmissiva e risorse computazionali fino a saturare i canali di trasmissione e causare congestione e collasso del servizio.

% [IMMAGINE DEL FILE DI SLIDE CON NOME "1) Ripasso_2"] Decremento del valore TTL da 255 a 254 e 253 attraverso il passaggio tra router intermedi da Host A a Host B — Slide 2 / Pagina 2 del PDF

Per scongiurare il fenomeno del blocco permanente, l'intestazione IPv4 implementa il campo Time to Live (TTL):
\begin{itemize}
    \item \textbf{Dimensione e funzione}: il campo ha una dimensione di 8 bit, consentendo di rappresentare un intero compreso nell'intervallo $[0, 255]$. Sebbene originariamente pensato per misurare il tempo di vita residuo in secondi, nell'ingegneria di rete moderna il TTL funge da contatore di transiti (\emph{hop counter}).
    \item \textbf{Regola di decremento}: l'host sorgente inizializza il campo a un valore predefinito (tipicamente 64, 128 o 255 a seconda dello stack del sistema operativo). Ogni volta che un pacchetto attraversa un apparato di livello 3 (router), tale dispositivo elabora l'intestazione e decrementa il valore del TTL esattamente di 1:
    \[
    \text{TTL}_{\text{nuovo}} = \text{TTL}_{\text{precedente}} - 1
    \]
    \item \textbf{Scarto del pacchetto}: non appena un router riceve un pacchetto e, applicando la riduzione, rileva che $\text{TTL} = 0$, il pacchetto viene scartato (\emph{dropped}) ed eliminato dai buffer di memoria.
    \item \textbf{Segnalazione dell'evento}: contestualmente allo scarto per azzeramento del TTL, il router che ha eliminato il frame genera un messaggio diagnostico di errore ICMP (\emph{Time Exceeded}) e lo trasmette all'indirizzo IP della sorgente originaria ricavato dall'header del pacchetto terminato, a patto che la generazione di messaggi ICMP di errore non sia stata disabilitata per criteri di sicurezza o di \emph{rate limiting} sulle interfacce dell'apparato.
\end{itemize}

\subsection{Demultiplazione Con Il Campo Protocol}
L'architettura TCP/IP poggia sul principio della stratificazione e del multiplexing/demultiplexing dei flussi informativi. Un singolo host o interfaccia di rete gestisce contemporaneamente sessioni applicative concorrenti incapsulate in protocolli differenti di livello superiore. Il compito di indirizzare il payload estratto dal datagramma IP al corretto modulo software o processo dello strato superiore è demandato al campo Protocol.

% [IMMAGINE DEL FILE DI SLIDE CON NOME "1) Ripasso_2"] Demultiplazione del pacchetto IP verso lo strato superiore basata sui valori del campo Protocol (1 per ICMP, 6 per TCP, 17 per UDP) — Slide 3 / Pagina 3 del PDF

Il campo \textbf{Protocol} (dimensione di 8 bit) contiene un codice numerico univoco standardizzato dalla IANA (\emph{Internet Assigned Numbers Authority}):
\begin{itemize}
    \item Può identificare un protocollo ausiliario cooperante allo stesso livello di rete, come ad esempio l'ICMP (\emph{Internet Control Message Protocol}), associato al valore esadecimale \texttt{0x01} (valore decimale 1).
    \item Può specificare un protocollo di trasporto di livello superiore:
    \begin{itemize}
        \item Transmission Control Protocol (TCP): identificato dal valore \texttt{0x06} (decimale 6).
        \item User Datagram Protocol (UDP): identificato dal valore \texttt{0x11} (decimale 17).
    \end{itemize}
\end{itemize}
All'arrivo del datagramma, lo strato IP della destinazione esamina questo campo e passa direttamente il blocco dati (\emph{User data}) al buffer di ricezione del driver di protocollo corrispondente.

\newpage
\section{Struttura Dell'Indirizzamento IPv4 E Metodologie Di Subnetting}

\subsection{Anatomia Di Un Indirizzo IPv4 E Maschera Di Rete}
Un indirizzo IPv4 è formalmente definito come un identificativo numerico binario della lunghezza di 32 bit (4 byte), convenzionalmente rappresentato mediante la notazione decimale puntata (\emph{dotted decimal notation}), nella quale ciascun byte (ottetto) viene convertito in un numero intero compreso tra 0 e 255, separato da punti (ad esempio \texttt{192.168.10.1}).

Dal punto di vista topologico e funzionale, un indirizzo IPv4 si suddivide in due componenti gerarchiche disgiunte:
\begin{enumerate}
    \item \textbf{Network Part} (\emph{Network ID}): la porzione più significativa di bit, deputata a identificare in modo univoco il segmento di rete logico a cui l'apparato appartiene. Gli indirizzi IP non incorporano nativamente alcuna informazione geografica o fisica: tutti gli host che condividono il medesimo identificativo di rete si considerano posizionati all'interno della medesima rete logica, indipendentemente dalla loro collocazione fisica.
    \item \textbf{Host Part} (\emph{Host ID}): la porzione meno significativa di bit, utilizzata per discriminare e identificare univocamente le singole interfacce di rete attestate su quel determinato segmento di rete.
\end{enumerate}

% [IMMAGINE DEL FILE DI SLIDE CON NOME "1) Ripasso_2"] Struttura binaria dell'indirizzo 192.168.10.1 con maschera 255.255.255.0 che separa la Network part dalla Host part (/24) — Slide 4 / Pagina 4 del PDF

Per consentire a un dispositivo di rete o a un algoritmo di routing di determinare la linea di demarcazione tra la parte di rete e la parte host, si associa all'indirizzo una maschera di sottorete (\emph{network mask} o \emph{subnet mask}). La maschera è una sequenza di 32 bit costituita rigorosamente da una successione ininterrotta di bit pari a $1$ seguita da una successione ininterrotta di bit pari a $0$:
\begin{itemize}
    \item I bit impostati a $1$ mappano le posizioni della componente di rete (\emph{Network bits}).
    \item I bit impostati a $0$ mappano le posizioni della componente host (\emph{Host bits}).
\end{itemize}
La lunghezza della maschera corrisponde semplicemente al numero di bit $1$ presenti. Nella prassi operativa contemporanea si utilizza la notazione compatta CIDR (\emph{Classless Inter-Domain Routing}), anteponendo una barra obliqua seguita dal numero di bit a 1 (es. \texttt{/24}). Ad esempio, la maschera \texttt{255.255.255.0} corrisponde alla notazione \texttt{/24}, mentre una maschera minimale come \texttt{252.0.0.0} è costituita da 6 bit ad uno seguiti da 26 zeri, corrispondendo a un prefisso \texttt{/6}. L'espressione complessiva assume la forma compatta:
\[
\texttt{192.168.10.1 255.255.255.0} \iff \texttt{192.168.10.1/24}
\]

\subsection{Tipologie Di Indirizzi All'Interno Di Un Segmento Di Rete}
All'interno di un qualsiasi segmento di rete IPv4 definito da una maschera con $n$ bit destinati alla porzione host, si determinano in linea teorica $2^n$ combinazioni binarie distinte. Tuttavia, non tutti gli indirizzi appartenenti a questo spazio possono essere assegnati alle interfacce fisiche o logiche dei dispositivi terminali.

% [IMMAGINE DEL FILE DI SLIDE CON NOME "1) Ripasso_2"] Classificazione dei tipi di indirizzi: Network address (tutti 0), Broadcast address (tutti 1) e Available addresses (da 1 a 2^n - 2) — Slide 8 / Pagina 8 del PDF

Si identificano tre categorie fondamentali:
\begin{enumerate}
    \item \textbf{Indirizzo di Rete} (\emph{Network Address}): è la configurazione in cui tutti i bit della porzione host sono posti contemporaneamente a zero logico ($0$). Tale indirizzo identifica formalmente l'intero segmento di rete all'interno delle tabelle di routing e non può in alcun caso essere assegnato all'interfaccia di un host. Ad esempio, per la rete \texttt{192.168.10.0/24}, l'indirizzo di rete è \texttt{192.168.10.0}.
    \item \textbf{Indirizzo di Broadcast Diretto} (\emph{Broadcast Address}): è la configurazione in cui tutti i bit della porzione host sono posti contemporaneamente a uno logico ($1$). Rappresenta un indirizzo speciale adibito alla trasmissione broadcast locale: qualsiasi datagramma inviato con questo indirizzo di destinazione viene captato e consegnato a tutte le stazioni collegate a quel segmento. Anche in questo caso, è categoricamente proibito assegnarlo a una singola interfaccia. Nell'esempio \texttt{/24}, il broadcast corrisponde a \texttt{192.168.10.255}.
    \item \textbf{Indirizzi Assegnabili agli Host} (\emph{Available Host Addresses}): sono tutte le combinazioni binarie intermedie in cui i bit di host non sono né tutti $0$ né tutti $1$. Solo questi indirizzi possono essere attribuiti alle schede di rete di server, router, client e stampanti.
\end{enumerate}

Da questa strutturazione discende la formula cardine per il dimensionamento delle sottoreti IP: se la porzione host comprende $n$ bit, si ha:
\begin{align*}
N_{\text{totali}} &= 2^n \\
\boxed{N_{\text{utilizzabili}} = 2^n - 2}
\end{align*}
La sottrazione di 2 tiene conto dell'indisponibilità strutturale dell'indirizzo di rete e dell'indirizzo di broadcast.

\subsection{Indirizzamento Con Classi (\emph{Classful Addressing}) E Limiti}
Agli albori dell'architettura Internet, lo spazio degli indirizzi IPv4 (pari complessivamente a $2^{32} = 4\,294\,967\,296$ indirizzi) venne ripartito in modo rigido e prefissato in cinque classi logiche, identificate dall'analisi dei bit più significativi (\emph{Most Significant Bits}, MSB) del primo ottetto.

% [IMMAGINE DEL FILE DI SLIDE CON NOME "1) Ripasso_2"] Tabella riepilogativa della classificazione Classful degli indirizzi IP (Classi da A a E) con prefissi binari, intervalli e maschere di default — Slide 5 / Pagina 5 del PDF

\begin{table}[H]
\centering
\small % Riduce leggermente il corpo del testo per alloggiare comodamente gli indirizzi in \texttt
\renewcommand{\arraystretch}{1.25} % Migliora la spaziatura verticale tra le righe
\setlength{\tabcolsep}{5pt} % Ottimizza lo spazio tra le colonne a favore della colonna descrittiva
\caption{Classificazione \emph{Classful} dello spazio di indirizzamento IPv4.}
\label{tab:classful}

\begin{tabularx}{\linewidth}{l c l l >{\raggedright\arraybackslash}X}
\toprule
\textbf{Classe} & \textbf{Prefisso MSB} & \textbf{Intervallo Indirizzi} & \textbf{Maschera Base} & \textbf{Destinazione Funzionale} \\
\midrule
\textbf{Classe A} & \texttt{0} & \texttt{0.0.0.0 -- 127.255.255.255} & \texttt{/8} ($255.0.0.0$) & Reti geografiche globali (Unicast) \\
\textbf{Classe B} & \texttt{10} & \texttt{128.0.0.0 -- 191.255.255.255} & \texttt{/16} ($255.255.0.0$) & Medie/grandi organizzazioni (Unicast) \\
\textbf{Classe C} & \texttt{110} & \texttt{192.0.0.0 -- 223.255.255.255} & \texttt{/24} ($255.255.255.0$) & Piccole reti locali (Unicast) \\
\textbf{Classe D} & \texttt{1110} & \texttt{224.0.0.0 -- 239.255.255.255} & Non applicabile & Trasmissioni Multicast uno-a-molti \\
\textbf{Classe E} & \texttt{1111} & \texttt{240.0.0.0 -- 255.255.255.255} & Non applicabile & Riservata per scopi sperimentali e ricerca \\
\bottomrule
\end{tabularx}
\end{table}

Solo le classi A, B e C definiscono indirizzi unicast, ossia indirizzi che possono essere configurati sulle interfacce fisiche di apparati e stazioni di lavoro. In questo contesto, il termine host denota formalmente qualsiasi apparato dotato di interfaccia IP, includendo calcolatori, server o interfacce di instradamento di un router.

Le capacità dimensionali teoriche delle classi unicast risultano drasticamente differenziate:
\begin{itemize}
    \item \textbf{Classe A}: riserva 8 bit alla rete e 24 bit agli host. Permette la creazione di 128 identificatori di rete teorici ($2^7 = 128$, con la rete \texttt{0.0.0.0} e \texttt{127.0.0.0} riservate, lasciando 126 o 127 reti pratiche), ciascuna in grado di accogliere $2^{24} = 16\,777\,216$ indirizzi ($16\,777\,214$ host attivi).
    \item \textbf{Classe B}: impiega 16 bit per la rete e 16 bit per gli host. Fornisce $2^{14} = 16\,384$ reti distinte, ognuna con $2^{16} = 65\,536$ indirizzi totali ($65\,534$ host utilizzabili).
    \item \textbf{Classe C}: assegna 24 bit alla rete e 8 bit agli host. Mette a disposizione ben $2^{21} = 2\,097\,152$ reti, ma ciascuna può ospitare al più $2^8 = 256$ indirizzi ($254$ host attivi).
\end{itemize}

L'architettura basata su classi presentava due debolezze strutturali non sostenibili per la crescita della rete:
\begin{enumerate}
    \item \textbf{Granularità eccessivamente rigida e spreco massivo di indirizzi}: una media impresa con un fabbisogno di 1000 nodi non poteva essere soddisfatta da una rete di Classe C (massimo 254 nodi), rendendo necessaria l'allocazione di un blocco di Classe B intero (65534 host). Tale scelta condannava oltre il 98\% dello spazio assegnato a rimanere inutilizzato.
    \item \textbf{Dimensioni ingestibili dei domini di broadcast}: attestare $65\,000$ nodi su una singola rete fisica di livello 2 genera un dominio di broadcast gigantesco. Il traffico di servizio generato spontaneamente dai protocolli di rete (es. ARP) provoca tempeste di broadcast (\emph{broadcast storm}), saturando la banda trasmissiva e sovraccaricando le CPU di ogni singola stazione attestata sul link.
\end{enumerate}

\subsection{Subnetting E Tecnica Del Prestito Di Bit (VLSM)}
Per risolvere le criticità dell'approccio \emph{classful}, è stata introdotta la tecnologia di subnetting (\emph{suddivisione in sottoreti}), integrata formalmente dal paradigma VLSM (\emph{Variable Length Subnet Mask}). Il subnetting consiste nel frammentare un'ampia rete logica originaria in molteplici sottoreti più piccole, ciascuna assegnata a un dominio di broadcast indipendente interconnesso mediante router.

% [IMMAGINE DEL FILE DI SLIDE CON NOME "1) Ripasso_2"] Confronto concettuale tra una rete di Classe B non partizionata (dominio di broadcast da 65536 IP saturo) e una partizionata in sottoreti isolate — Slide 20 / Pagina 20 del PDF

La ripartizione interna si attua attraverso il meccanismo del prestito di bit (\emph{bit borrowing}): si sottraggono bit alla porzione host dell'indirizzo originario destinandoli a una nuova porzione intermedia definita Subnet ID.

Analizziamo il procedimento analitico e formale partendo da un segmento di Classe C convenzionale, come mostrato nelle slide:
\begin{itemize}
    \item Rete di partenza: \texttt{192.168.10.0/24} con maschera di default \texttt{255.255.255.0}.
    \item Configurazione iniziale: 24 bit di rete, 8 bit di host ($n = 8$).
    \item Capacità: $2^8 = 256$ indirizzi totali, di cui $254$ utilizzabili.
\end{itemize}

% [IMMAGINE DEL FILE DI SLIDE CON NOME "1) Ripasso_2"] Analisi binaria della rete di partenza 192.168.10.0/24 con dettagli su Host bits, Network address e Broadcast address — Slide 21 / Pagina 21 del PDF

Supponiamo ora di voler segmentare questa rete prestando $k = 1$ bit dalla porzione host:
\begin{itemize}
    \item Numero di sottoreti ottenute: $2^k = 2^1 = 2$ sottoreti.
    \item Nuova lunghezza del prefisso di rete: $24 + 1 = 25$ bit (\texttt{/25}).
    \item Nuova maschera di sottorete binaria: i primi 25 bit sono posti a $1$, i restanti 7 a $0$:
    \[
    \texttt{11111111.11111111.11111111.10000000}_2 \iff \texttt{255.255.255.128}
    \]
    \item Bit residui di host: $n' = 8 - 1 = 7$ bit.
    \item Capacità di ciascuna sottorete:
    \begin{align*}
    N_{\text{totali}} &= 2^7 = 128 \\
    N_{\text{utilizzabili}} &= 2^7 - 2 = 126 \text{ host per ciascuna sottorete}
    \end{align*}
\end{itemize}

% [IMMAGINE DEL FILE DI SLIDE CON NOME "1) Ripasso_2"] Prestito di 1 bit dalla porzione host e generazione delle due sottoreti /25 con calcolo dei totali e disponibili — Slide 22 / Pagina 22 del PDF

Il singolo bit preso in prestito (il bit più significativo del quarto ottetto) può assumere i due valori alternativi $0$ o $1$:
\begin{enumerate}
    \item Se il bit di subnet è posto a $0$:
    \begin{itemize}
        \item \textbf{Indirizzo di Rete}: ponendo tutti i 7 bit di host a $0$:
        \[
        \texttt{192.168.10.00000000}_2 \iff \texttt{192.168.10.0/25}
        \]
        \item \textbf{Indirizzo di Broadcast}: ponendo tutti i 7 bit di host a $1$:
        \[
        \texttt{192.168.10.01111111}_2 \iff \texttt{192.168.10.127/25}
        \]
        \item Intervallo host validi: da \texttt{192.168.10.1} a \texttt{192.168.10.126}.
    \end{itemize}
    \item Se il bit di subnet è posto a $1$:
    \begin{itemize}
        \item \textbf{Indirizzo di Rete}: ponendo tutti i 7 bit di host a $0$:
        \[
        \texttt{192.168.10.10000000}_2 \iff \texttt{192.168.10.128/25}
        \]
        \item \textbf{Indirizzo di Broadcast}: ponendo tutti i 7 bit di host a $1$:
        \[
        \texttt{192.168.10.11111111}_2 \iff \texttt{192.168.10.255/25}
        \]
        \item Intervallo host validi: da \texttt{192.168.10.129} a \texttt{192.168.10.254}.
    \end{itemize}
\end{enumerate}

% [IMMAGINE DEL FILE DI SLIDE CON NOME "1) Ripasso_2"] Dettaglio formale del calcolo dell'indirizzo di rete per le sottoreti /25 impostando i bit di host a zero — Slide 23 / Pagina 23 del PDF

% [IMMAGINE DEL FILE DI SLIDE CON NOME "1) Ripasso_2"] Dettaglio formale del calcolo dell'indirizzo di broadcast per le sottoreti /25 impostando i bit di host a uno — Slide 24 / Pagina 24 del PDF

\subsection{CIDR E Aggregazione Delle Rotte (\emph{Route Summarization})}
Mentre il subnetting divide uno spazio di rete in frazioni gerarchiche inferiori per scopi di ingegneria locale, la tecnica speculare definita CIDR (\emph{Classless Inter-Domain Routing}) opera in ottica inter-dominio e di transito globale. Il CIDR svincola completamente l'instradamento dalla rigida nozione di classe, introducendo la possibilità di definire prefissi continui di qualsiasi lunghezza arbitraria.

Una delle applicazioni più rilevanti del CIDR risiede nella Route Summarization (o \emph{Supernetting}). Quando un router di confine (\emph{Autonomous System Border Router} o router di core) gestisce instradamenti verso molteplici segmenti di rete contigui, non è efficiente annunciare ciascuna sottorete come voce distinta nella tabella di routing. Attraverso l'aggregazione, molteplici prefissi contigui vengono collassati in una singola macro-rotta cumulativa, riducendo drasticamente il numero di voci nella tabella di inoltro (\emph{Forwarding Information Base}, FIB) e l'occupazione di memoria TCAM sugli apparati.

% [IMMAGINE DEL FILE DI SLIDE CON NOME "1) Ripasso_2"] Diagramma del CIDR e Route Summarization con allineamento binario di quattro reti per ricavare la rotta aggregata 192.168.8.0/21 — Slide 7 / Pagina 7 del PDF

Consideriamo l'esempio mostrato nella documentazione di riferimento:
\begin{align*}
\text{Rete 1: } & \texttt{192.168.12.0/22} \\
\text{Rete 2: } & \texttt{192.168.10.0/23} \\
\text{Rete 3: } & \texttt{192.168.9.0/21} \\
\text{Rete 4: } & \texttt{192.168.14.0/23}
\end{align*}
Per procedere alla summarization, convertiamo in forma binaria la componente variabile degli indirizzi, ossia il terzo ottetto:
\begin{align*}
\texttt{192.168.} & \mathbf{00001}100 \texttt{.0} \quad (/22) \\
\texttt{192.168.} & \mathbf{00001}010 \texttt{.0} \quad (/23) \\
\texttt{192.168.} & \mathbf{00001}001 \texttt{.0} \quad (/21) \\
\texttt{192.168.} & \mathbf{00001}110 \texttt{.0} \quad (/23)
\end{align*}
Osservando le stringhe binarie allineate da sinistra verso destra, si constata che tutti e quattro i blocchi di indirizzi condividono un prefisso comune identico per i primi 21 bit:
\[
\overbrace{192.168.}^{\text{16 bit}} \overbrace{\mathbf{00001}}^{\text{5 bit}} = 21 \text{ bit coincidenti}
\]
I restanti 11 bit (3 bit del terzo ottetto e gli 8 bit del quarto ottetto) coprono in modo continuo le combinazioni binarie comprese tra il valore decimale $8$ ($\texttt{00001000}_2$) e il valore decimale $15$ ($\texttt{00001111}_2$). Azzerando tutti i bit al di fuori del prefisso comune di 21 bit si ricava l'indirizzo di supernetting:
\[
\texttt{192.168.00001000.00000000}_2 \implies \boxed{\texttt{192.168.8.0/21}}
\]
La singola voce aggregata \texttt{192.168.8.0/21} sostituisce le quattro rotte individuali nell'annuncio di routing, ottimizzando la scalabilità globale della rete.

\newpage
\section{Indirizzi Privati E Tecnologie NAT (Network Address Translation)}

\subsection{Indirizzamento Privato (RFC 1918) E Ruolo Architetturale Del NAT}
Con la proliferazione globale dei calcolatori e dei servizi di rete verificatasi a partire dagli anni '90, l'esaurimento dello spazio di indirizzamento a 32 bit di IPv4 è divenuto una minaccia concreta per la scalabilità di Internet. Per mitigare tale problematica in attesa dell'adozione estesa di IPv6, la IANA ha standardizzato la separazione formale tra:
\begin{itemize}
    \item \textbf{Indirizzi IP Pubblici}: blocchi di indirizzi registrati e assegnati in modo rigorosamente univoco a livello mondiale. Sono direttamente instradabili lungo le dorsali globali di Internet.
    \item \textbf{Indirizzi IP Privati}: blocchi di indirizzi riservati ad uso esclusivo di reti interne aziendali, campus universitari o ambienti domestici. Non possiedono univocità globale e possono essere riutilizzati in modo concorrente e sovrapposto da una miriade di reti private distinte senza generare conflitti, purché tali reti rimangano confinate e isolate dalla rete pubblica.
\end{itemize}

% [IMMAGINE DEL FILE DI SLIDE CON NOME "1) Ripasso_2"] Connessione di domini privati sovrapposti (10.0.0.0/8 e 192.168.1.0/24) a Internet mediante router NAT — Slide 9 / Pagina 9 del PDF

I blocchi privati standardizzati (RFC 1918) sono distribuiti sulle tre classi principali:
\begin{itemize}
    \item \textbf{Classe A}: da \texttt{10.0.0.0} a \texttt{10.255.255.255} (notazione CIDR: \texttt{10.0.0.0/8}, pari a 1 rete da oltre 16 milioni di host).
    \item \textbf{Classe B}: da \texttt{172.16.0.0} a \texttt{172.31.255.255} (notazione CIDR: \texttt{172.16.0.0/12}, pari a 16 reti di Classe B continue).
    \item \textbf{Classe C}: da \texttt{192.168.0.0} a \texttt{192.168.255.255} (notazione CIDR: \texttt{192.168.0.0/16}, pari a 256 reti di Classe C continue).
\end{itemize}

I router posizionati lungo la rete pubblica Internet sono programmaticamente configurati per scartare qualsiasi pacchetto recante come indirizzo sorgente o di destinazione un IP appartenente a questi intervalli, poiché non esistono rotte globali di inoltro per essi. Qualora un dispositivo situato all'interno di un dominio privato necessiti di accedere a una risorsa situata su Internet (o viceversa), è obbligatorio interporre sull'apparato di confine (\emph{egress gateway}, router perimetrale o firewall) il meccanismo di Network Address Translation (NAT).

% [IMMAGINE DEL FILE DI SLIDE CON NOME "1) Ripasso_2"] Flusso di traduzione bidirezionale del NAT tra host privato 192.168.1.10, router con IP pubblico 122.1.2.1 e web server pubblico 200.1.2.3 — Slide 10 / Pagina 10 del PDF

La dinamica base del NAT illustrata nelle slide procede attraverso le seguenti fasi sequenziali:
\begin{enumerate}
    \item L'host interno (ad es. un PC con indirizzo \texttt{192.168.1.10}) emette un pacchetto diretto al web server pubblico \texttt{200.1.2.3}. Il pacchetto ha:
    \[
    \text{Source IP: } \texttt{192.168.1.10}, \quad \text{Destination IP: } \texttt{200.1.2.3}
    \]
    \item Il router di confine riceve la trama sulla propria interfaccia interna. Prima di instradarla sull'interfaccia WAN, esegue una traduzione a livello di intestazione, sostituendo l'indirizzo sorgente privato con il proprio indirizzo pubblico autorizzato (es. \texttt{122.1.2.1}), tenendo traccia dell'operazione in una tabella di conversione:
    \[
    \text{Source IP: } \texttt{122.1.2.1}, \quad \text{Destination IP: } \texttt{200.1.2.3}
    \]
    \item Il server remoto riceve il pacchetto e trasmette la risposta instradandola verso l'IP pubblico del router di confine:
    \[
    \text{Source IP: } \texttt{200.1.2.3}, \quad \text{Destination IP: } \texttt{122.1.2.1}
    \]
    \item Il router di confine intercetta il pacchetto in ingresso, consulta la propria tabella di stato NAT, individua la corrispondenza con l'host privato originatore, riscrive l'indirizzo di destinazione ripristinando \texttt{192.168.1.10} e immette il pacchetto nella rete locale:
    \[
    \text{Source IP: } \texttt{200.1.2.3}, \quad \text{Destination IP: } \texttt{192.168.1.10}
    \]
\end{enumerate}

\subsection{Tassonomia Delle Modalità Di Implementazione Del NAT}
A seconda delle policy di allocazione, della tipologia di vincolo numerico e del livello dello stack analizzato, il NAT si declina in molteplici varianti architetturali.

\subsubsection{Static NAT (NAT Statico)}
Lo Static NAT implementa una mappatura fissa, permanente e biunivoca ($1:1$) tra un indirizzo IP privato interno e un determinato indirizzo IP pubblico esterno.

% [IMMAGINE DEL FILE DI SLIDE CON NOME "1) Ripasso_2"] Architettura di Static NAT con tabella di mappatura 1:1 permanente tra host interni 192.168.1.x e IP pubblici 122.1.2.x — Slide 11 / Pagina 11 del PDF

Caratteristiche fondamentali:
\begin{itemize}
    \item \textbf{Accesso Bidirezionale}: poiché la correlazione è codificata in modo deterministico e immutabile nella tabella di traduzione, la connessione può essere instaurata in entrambe le direzioni. Sia l'host interno può accedere a Internet, sia un apparato esterno può contattare direttamente la stazione interna inviando traffico all'indirizzo pubblico abbinato.
    \item \textbf{Limite intrinseco}: non attua alcun effettivo risparmio di indirizzi pubblici, poiché per ogni host privato che deve accedere all'esterno è richiesto un indirizzo IP pubblico dedicato. Viene utilizzato primariamente per assegnare identità pubbliche fisse a server interni (es. server web o mail).
\end{itemize}

\subsubsection{Dynamic NAT (NAT Dinamico)}
Il Dynamic NAT supera l'inconveniente della fissità dello Static NAT definendo un insieme di indirizzi pubblici a disposizione, denominato NAT Address Pool.

% [IMMAGINE DEL FILE DI SLIDE CON NOME "1) Ripasso_2"] Schema del Dynamic NAT con pool di IP pubblici marcati dinamicamente nello stato In Use o Not Use — Slide 12 / Pagina 12 del PDF

Funzionamento operativo:
\begin{itemize}
    \item Gli indirizzi pubblici appartenenti al pool sono inizialmente posti nello stato di riposo (\emph{Not Use}).
    \item Quando una stazione privata genera un flusso destinato all'esterno, il router di confine attinge un indirizzo libero dal pool, crea temporaneamente una voce nella tabella NAT contrassegnando l'IP pubblico come occupato (\emph{In Use}), e traduce il traffico.
    \item Terminata la sessione e decorso un tempo di inattività (\emph{timeout}), l'IP pubblico viene rilasciato e torna disponibile per qualsiasi altra stazione.
    \item \textbf{Limite}: il Dynamic NAT standard opera ancora in modalità senza traduzione di porta (\emph{No-Port Address Translation}, o No-PAT). La corrispondenza resta istantaneamente di tipo $1:1$: se il pool contiene 3 indirizzi pubblici, al massimo 3 stazioni interne possono accedere contemporaneamente all'esterno; eventuali richieste contemporanee di un quarto host vengono scartate per esaurimento del pool.
\end{itemize}

\subsubsection{NAPT (\emph{Network Address And Port Translation})}
Il NAPT (indicato storicamente anche come PAT o NAT Overload) costituisce l'evoluzione decisiva per risolvere la penuria di indirizzi. Esso non si limita alla manipolazione degli indirizzi di livello 3, ma esamina e traduce congiuntamente le porte di livello di trasporto (TCP e UDP) a livello 4.

% [IMMAGINE DEL FILE DI SLIDE CON NOME "1) Ripasso_2"] Architettura del NAPT con multiplexing 1:n su singolo IP pubblico mediato da porte Layer 4 — Slide 13 / Pagina 13 del PDF

Con il NAPT si ottiene una mappatura molti-a-uno o molti-a-pochi ($1:n$ o $m:n$): decine o centinaia di host interni possono condividere in contemporanea il medesimo indirizzo IP pubblico esterno. Il router distingue i singoli flussi associando a ciascun dispositivo interno una specifica porta di trasporto sorgente pubblica (\emph{translated port}).

Esaminiamo la dinamica di funzionamento analitica mostrata nei dettagli esecutivi delle slide:

% [IMMAGINE DEL FILE DI SLIDE CON NOME "1) Ripasso_2"] Esempio NAPT Fase 1: Host sorgente emette richiesta verso web server, il router seleziona porta libera e crea voce in tabella — Slide 14 / Pagina 14 del PDF

\begin{enumerate}
    \item \textbf{Flusso Outbound (Richiesta da client interno)}:
    \begin{itemize}
        \item L'host privato \texttt{192.168.1.1} apre una connessione HTTP verso il web server \texttt{200.1.2.3:80}, generando il pacchetto con parametri:
        \[
        \text{Socket Sorgente: } \texttt{192.168.1.1:10321}, \quad \text{Socket Destinazione: } \texttt{200.1.2.3:80}
        \]
        \item Il router NAT intercetta la trama, attinge un IP pubblico dal pool (es. \texttt{122.1.2.2}) e alloca una porta sorgente pubblica libera (es. \texttt{1025}).
        \item Crea una riga di stato nella tabella di traduzione (\emph{NAT mapping table}):
        \[
        [\texttt{192.168.1.1:10321}] \longleftrightarrow [\texttt{122.1.2.2:1025}]
        \]
        \item Riscrive l'intestazione IP e TCP e immette il pacchetto su Internet:
        \[
        \text{Sorgente: } \texttt{122.1.2.2:1025}, \quad \text{Destinazione: } \texttt{200.1.2.3:80}
        \]
    \end{itemize}

% [IMMAGINE DEL FILE DI SLIDE CON NOME "1) Ripasso_2"] Esempio NAPT Fase 2: Risposta dal web server, corrispondenza nella tabella NAPT e ripristino di IP e porta interna — Slide 15 / Pagina 15 del PDF

    \item \textbf{Flusso Inbound (Risposta da server remoto)}:
    \begin{itemize}
        \item Il web server evade la richiesta e invia la risposta diretta alla socket pubblica:
        \[
        \text{Sorgente: } \texttt{200.1.2.3:80}, \quad \text{Destinazione: } \texttt{122.1.2.2:1025}
        \]
        \item Il router NAT riceve il pacchetto dall'esterno, ne analizza l'intestazione L3/L4 rilevando la coppia $[\texttt{IP:Port}]_{\text{dest}} = [\texttt{122.1.2.2:1025}]$.
        \item Esegue una ricerca (\emph{lookup}) nella tabella NAPT: riscontra l'associazione attiva con la socket interna $[\texttt{192.168.1.1:10321}]$.
        \item Riscrive la destinazione dell'intestazione e inoltra il frame sul segmento privato, consegnandolo all'host generatore.
    \end{itemize}
\end{enumerate}

\subsubsection{Easy IP}
La modalità Easy IP è una specializzazione architetturale del NAPT concepita per contesti in cui la rete periferica non dispone di un pool di indirizzi statici, né tantomeno di un IP pubblico fisso. Tipicamente, nelle connessioni residenziali o per piccole filiali, l'indirizzo IP pubblico viene assegnato dinamicamente dall'operatore di telecomunicazioni (ISP) all'interfaccia WAN del router tramite protocolli come DHCP (\emph{Dynamic Host Configuration Protocol}) o PPPoE (\emph{Point-to-Point Protocol over Ethernet}).

% [IMMAGINE DEL FILE DI SLIDE CON NOME "1) Ripasso_2"] Architettura Easy IP in cui la traduzione NAPT avviene impiegando direttamente l'indirizzo dinamico dell'interfaccia WAN — Slide 16 / Pagina 16 del PDF

In Easy IP non viene configurato alcun NAT pool. Il meccanismo di traduzione intercetta dinamicamente l'indirizzo associato in quel momento all'interfaccia di uscita (\emph{egress interface}) e mappa le socket di tutte le stazioni private sull'indirizzo di tale interfaccia, differenziandole unicamente per numero di porta.

\subsubsection{NAT Server (\emph{Port Forwarding})}
Nelle implementazioni NAPT ed Easy IP descritte, le sessioni possono essere initiate unicamente in direzione \emph{outbound} (dall'interno verso l'esterno); un client situato sulla rete pubblica non ha alcuna possibilità di avviare una sessione verso una macchina interna, poiché in assenza di un flusso previo non esiste alcuna voce nella tabella di traduzione.

% [IMMAGINE DEL FILE DI SLIDE CON NOME "1) Ripasso_2"] Schema del meccanismo NAT Server con mappatura statica tra socket pubblica 122.1.2.1:80 e server interno 192.168.1.10:80 — Slide 17 / Pagina 17 del PDF

Per consentire l'erogazione di servizi pubblici (es. server web, mail, o istanze applicative nel cloud privato) posizionati all'interno della DMZ o della rete interna privata, si adotta la funzionalità NAT Server (nota comunemente come \emph{Port Forwarding} o Destination NAT statico):
\begin{itemize}
    \item L'amministratore inserisce a priori una mappatura statica nella tabella del router:
    \[
    [\text{Public IP}:\text{Public Port}] \longleftrightarrow [\text{Private IP}:\text{Private Port}]
    \]
    \item Ad esempio, abbinando la coppia \texttt{122.1.2.1:80} all'host interno \texttt{192.168.1.10:80}.
    \item Qualsiasi client esterno che punti il browser all'indirizzo pubblico \texttt{122.1.2.1} sulla porta 80 vedrà i propri pacchetti intercettati dal router, traslati nell'indirizzo di destinazione (\texttt{Destination IP} riscritta da \texttt{122.1.2.1} a \texttt{192.168.1.10}) e recapitati trasparentemente al server privato interno.
\end{itemize}

\newpage
\section{Indirizzi Speciali E Protocollo ICMP Per La Diagnostica Di Rete}

\subsection{Indirizzamenti IPv4 Con Funzione Dedicata}
Oltre alle canoniche ripartizioni \emph{classful} o private, lo standard IPv4 riserva specifici blocchi di indirizzi a funzionalità di segnalazione, autoconfigurazione e test locale del software di comunicazione.

% [IMMAGINE DEL FILE DI SLIDE CON NOME "1) Ripasso_2"] Tabella riepilogativa degli indirizzi IP speciali (Limited Broadcast, Any IP, Loopback, Link-local) e rispettive funzioni — Slide 18 / Pagina 18 del PDF

\begin{table}[H]
\centering
\caption{Quadro sinottico degli indirizzi speciali in IPv4.}
\label{tab:special_ip}
\begin{tabularx}{\textwidth}{l l X}
\toprule
\textbf{Tipologia Indirizzo} & \textbf{Ambito / Prefisso} & \textbf{Ruolo E Semantica Operativa} \\
\midrule
\textbf{Limited Broadcast} & \texttt{255.255.255.255} & Indirizzo di broadcast limitato al segmento fisico locale. Utilizzabile solo come IP di destinazione. I router bloccano categoricamente l'inoltro di questi pacchetti all'esterno del segmento. \\
\textbf{Any IP / Unspecified} & \texttt{0.0.0.0} & Identifica ``questa rete'' o l'assenza di un indirizzo valido. Utilizzato come IP sorgente nelle fasi di bootstrap o richiesta DHCP da host non ancora configurati; funge inoltre da rotta di default (\texttt{0.0.0.0/0}). \\
\textbf{Loopback Address} & \texttt{127.0.0.0/8} & Blocco riservato al collaudo software interno dello stack TCP/IP. I pacchetti emessi verso tale spazio non transitano mai sull'hardware di rete fisico. \\
\textbf{Link-Local Address} & \texttt{169.254.0.0/16} & Indirizzamento automatico per comunicazioni d'emergenza punto-a-punto su scala locale (\emph{APIPA}) qualora la configurazione dinamica DHCP fallisca. \\
\bottomrule
\end{tabularx}
\end{table}

Analizziamo il comportamento architetturale di ciascuna voce:
\begin{itemize}
    \item \textbf{Limited Broadcast (\texttt{255.255.255.255})}: a differenza del broadcast diretto (es. \texttt{192.168.10.255}), che presuppone la conoscenza della maschera di rete, il broadcast limitato consente a un host di comunicare con tutte le stazioni del proprio link fisico senza conoscere a priori la configurazione di sottorete. Il gateway d'inoltro (router) interrompe tassativamente la propagazione di queste trame per scongiurare la diffusione incontrollata di broadcast su scala inter-rete.
    \item \textbf{Any IP Address (\texttt{0.0.0.0})}: quando una macchina avvia il proprio sistema operativo ed è impostata per acquisire i parametri di rete via DHCP, la scheda di rete non possiede alcun indirizzo IP valido. La stazione emette una richiesta \emph{DHCP Request} incapsulata in un frame broadcast, inserendo come indirizzo IP di destinazione il limited broadcast (\texttt{255.255.255.255}) e come indirizzo IP sorgente il marcatore fittizio \texttt{0.0.0.0}.
    \item \textbf{Loopback Address (\texttt{127.0.0.0/8})}: l'indirizzo canonico più utilizzato è \texttt{127.0.0.1} (\emph{localhost}). Qualsiasi datagramma generato da un'applicazione con destinazione all'interno di questo blocco viene intercettato direttamente al livello software dallo stack IP e reindirizzato immediatamente al buffer di ricezione interno (\emph{loopback interface}), consentendo lo sviluppo e il test di servizi di rete indipendentemente dalla presenza di connettività fisica.
    \item \textbf{Link-Local (\texttt{169.254.0.0/16})}: quando un client configurato in modalità dinamica tenta ripetutamente la transazione di bootstrap senza ricevere risposta (ad esempio per assenza o guasto del server DHCP), interviene il meccanismo di autoconfigurazione (RFC 3927 / APIPA). L'host estrae uno pseudocasuale all'interno del segmento \texttt{169.254.0.0/16}, verifica mediante richieste ARP che non sia già in uso da un'altra stazione adiacente e lo assegna provvisoriamente alla propria interfaccia per preservare la connettività di base all'interno del link locale.
\end{itemize}

\subsection{Internet Control Message Protocol (ICMP)}
L'ICMP è un protocollo ausiliario del livello di rete concepito specificamente per scambiare informazioni di controllo, tracciare anomalie operative e notificare errori riscontrati durante il trattamento e l'inoltro dei datagrammi IP. Sebbene dal punto di vista concettuale appartenga al livello 3, i messaggi ICMP vengono incapsulati direttamente all'interno di un pacchetto IP convenzionale, come evidenziato dal campo \emph{Protocol} valorizzato a 1 (\texttt{0x01}).

% [IMMAGINE DEL FILE DI SLIDE CON NOME "1) Ripasso_2"] Struttura di incapsulamento dell'ICMP nel pacchetto IP e tabella dei codici Type/Code fondamentali — Slide 25 / Pagina 25 del PDF

Un messaggio ICMP generico presenta un'intestazione costituita da tre campi essenziali:
\begin{enumerate}
    \item \textbf{Type} (8 bit): definisce la tipologia semantica generale del messaggio.
    \item \textbf{Code} (8 bit): specifica la causale o il sottocaso diagnostico associato a quel determinato \emph{Type}.
    \item \textbf{Checksum} (16 bit): campo di controllo integrità algoritmico calcolato sull'intero messaggio ICMP.
\end{enumerate}
A questi campi segue un'area dati variabile a 32 bit, il cui significato dipende strettamente dal tipo di messaggio inviato:
\begin{itemize}
    \item Nei messaggi di tipo \emph{Redirect}, tale campo trasporta l'indirizzo IP del nuovo gateway a cui il mittente dovrà inoltrare il traffico futuro.
    \item Nei messaggi di tipo \emph{Echo Request} ed \emph{Echo Reply}, il campo contiene due valori identificativi cruciali: l'Identifier (16 bit) e il Sequence Number (16 bit), impiegati per associare deterministicamente ciascuna risposta alla corretta richiesta inviata.
\end{itemize}

\begin{table}[H]
\centering
\caption{Tipologie canoniche di messaggi ICMP e rispettiva codifica.}
\label{tab:icmp_types}
\begin{tabularx}{\textwidth}{c c l X}
\toprule
\textbf{Type} & \textbf{Code} & \textbf{Denominazione} & \textbf{Descrizione Funzionale} \\
\midrule
\texttt{0} & \texttt{0} & \emph{Echo Reply} & Risposta affermativa emessa a fronte della ricezione di un'istanza di Echo Request. \\
\texttt{3} & \texttt{0} & \emph{Network Unreachable} & Mancanza di rotta verso la rete di destinazione all'interno della FIB del router. \\
\texttt{3} & \texttt{1} & \emph{Host Unreachable} & Mancanza di risposta a livello L2 (es. risoluzione ARP fallita) per l'host finale. \\
\texttt{3} & \texttt{2} & \emph{Protocol Unreachable} & Il modulo del protocollo di trasporto richiesto nell'header IP non è attivo o supportato. \\
\texttt{3} & \texttt{3} & \emph{Port Unreachable} & La porta UDP specificata non corrisponde a nessun socket di processo in ascolto. \\
\texttt{5} & \texttt{0} & \emph{Redirect} & Notifica a un host dell'esistenza di un percorso di primo salto migliore sulla medesima rete. \\
\texttt{8} & \texttt{0} & \emph{Echo Request} & Richiesta di controllo connettività inviata a una destinazione remota. \\
\texttt{11} & \texttt{0} & \emph{Time Exceeded} & Segnalazione di esaurimento del contatore di transito (TTL decrementato a zero). \\
\bottomrule
\end{tabularx}
\end{table}

\subsection{Meccanismo Di Reindirizzamento: ICMP Redirection}
Quando un host non dispone di un protocollo di routing dinamico attivo, si affida a una rotta statica preconfigurata, ovvero al proprio Default Gateway. Può accadere tuttavia che, per motivi architetturali, una sottorete locale ospiti più router di uscita verso direttrici diverse.

% [IMMAGINE DEL FILE DI SLIDE CON NOME "1) Ripasso_2"] Flusso sequenziale del meccanismo di ICMP Redirection tra Host A, RTB (default gateway) ed RTA verso Server A — Slide 27 / Pagina 27 del PDF

Esaminiamo la dinamica descritta nello scenario applicativo:
\begin{itemize}
    \item L'host terminale \texttt{Host A} (\texttt{10.0.0.1/24}) deve trasmettere un flusso di pacchetti al server \texttt{Server A} (\texttt{20.0.0.1/24}).
    \item La configurazione IP di \texttt{Host A} imposta come default gateway il router \texttt{RTB} (\texttt{10.0.0.100/24}).
    \item Sulla medesima sottorete locale esiste un secondo apparato, il router \texttt{RTA} (\texttt{10.0.0.200/24}), il quale è direttamente collegato alla rete di destinazione \texttt{20.0.0.0/24} (con interfaccia \texttt{20.0.0.2/24}).
\end{itemize}
La procedura di segnalazione ed ottimizzazione avviene secondo i seguenti passaggi:
\begin{enumerate}
    \item \texttt{Host A} invia il pacchetto destinato a \texttt{20.0.0.1} instradandolo verso il proprio default gateway \texttt{RTB}.
    \item \texttt{RTB} riceve la trama, esamina la propria tabella di routing e constata che il percorso ottimale verso \texttt{20.0.0.0/24} impone come prossimo salto (\emph{next-hop}) il router \texttt{RTA}.
    \item Notando che l'interfaccia da cui il pacchetto è entrato e quella attraverso cui deve essere ritrasmesso appartengono al medesimo segmento di rete locale, \texttt{RTB} rileva che l'host mittente sta adottando un instradamento sub-ottimale.
    \item Pertanto, mentre inoltra il pacchetto transitorio verso \texttt{RTA}, \texttt{RTB} genera e invia a \texttt{Host A} un pacchetto ICMP Redirect (Type 5, Code 0), recante nel proprio campo dati l'IP di \texttt{RTA} (\texttt{10.0.0.200}).
    \item Ricevuto il messaggio di Redirect, \texttt{Host A} aggiorna dinamicamente la propria tabella di instradamento locale: tutti i successivi pacchetti destinati a \texttt{Server A} verranno inviati direttamente a \texttt{RTA}, eliminando il passaggio intermedio su \texttt{RTB}.
\end{enumerate}

\subsection{Applicazioni Diagnostiche Fondamentali: Ping E Tracert}
I messaggi ICMP costituiscono le fondamenta operative di due tra i più importanti strumenti diagnostici della suite Internet: il comando `ping` e il comando `traceroute` (denominato `tracert` sui sistemi operativi Windows o apparati di rete specifici).

\subsubsection{Diagnostica Di Raggiungibilità: Ping}
Il comando `ping` (Packet Internet Groper) rappresenta un'applicazione leggera e universale, disponibile in qualsiasi sistema operativo, preposta a verificare la raggiungibilità a livello 3 di un nodo remoto e a stimare la qualità del collegamento calcolando il Round-Trip Time (RTT), ossia il tempo impiegato dal messaggio per compiere il percorso di andata e ritorno.

% [IMMAGINE DEL FILE DI SLIDE CON NOME "1) Ripasso_2"] Output di sessione di Ping ed emissione ciclica di Echo Request ed Echo Reply per stimare la latenza RTT — Slide 28 / Pagina 28 del PDF

L'interazione si articola su principi rigorosi:
\begin{enumerate}
    \item Il client generatore emette un datagramma IP contenente un messaggio ICMP Echo Request (Type 8, Code 0), popolando l'intestazione con un identificatore di processo (\emph{Identifier}) e un numero di sequenza progressivo (\emph{Sequence Number}). Contestualmente, memorizza il marcatore temporale d'invio (\emph{timestamp}).
    \item Il destinatario intercetta il messaggio, commuta il tipo in ICMP Echo Reply (Type 0, Code 0), ricopia immutati l'Identifier e il Sequence Number originari e reinvia il frame al mittente.
    \item Il mittente rileva l'arrivo della risposta, confronta il Sequence Number con la richiesta emessa e, calcolando la differenza tra il timestamp corrente e quello di trasmissione, misura puntualmente l'RTT. L'esecuzione ciclica di più pacchetti consente di tracciare la percentuale di pacchetti persi (\emph{packet loss}) e i valori statistici minimo, medio e massimo di latenza (come evidenziato nel log di slide: 5 pacchetti trasmessi, 5 ricevuti, $0.00\%$ packet loss, RTT $\min/\text{avg}/\max = 30/40/70 \text{ ms}$).
\end{enumerate}

\subsubsection{Mappatura Topologica Multi-Hop: Tracert}
Lo strumento `traceroute` (`tracert`) assolve a una funzione più complessa: mappare sequenzialmente l'elenco esatto di tutti i router intermedi (\emph{hops}) attraversati lungo il cammino di instradamento che separa la sorgente dalla destinazione, rilevando latenze locali ed evidenziando eventuali anelli di instradamento.

% [IMMAGINE DEL FILE DI SLIDE CON NOME "1) Ripasso_2"] Funzionamento del comando Tracert con sonde basate su decremento ed esaurimento incrementale del campo TTL — Slide 29 / Pagina 29 del PDF

Il meccanismo sfrutta in modo ingegnoso la combinazione del campo TTL dell'intestazione IPv4 e dei messaggi di errore ICMP Time Exceeded:
\begin{enumerate}
    \item \textbf{Prima sonda}: la sorgente genera un pacchetto verso la destinazione finale impostando deliberatamente $\text{TTL} = 1$.
    \item Quando il primo router di transito lungo la direttrice riceve il pacchetto, applica la regola di elaborazione L3 e decrementa il TTL: $\text{TTL} = 1 - 1 = 0$. Rilevando l'azzeramento, il router scarta il pacchetto e restituisce alla sorgente un messaggio di errore ICMP Time Exceeded (Type 11, Code 0).
    \item La sorgente intercetta il messaggio ICMP: l'indirizzo IP del mittente di tale notifica rivela l'identità del primo nodo di transito (\emph{Hop 1}). La misura temporale tra invio e ricezione fornisce la stima della latenza di quella tratta.
    \item \textbf{Sonde successive}: il processo itera incrementando di volta in volta il campo TTL:
    \begin{itemize}
        \item Pacchetto con $\text{TTL} = 2$: supererà il primo nodo (che lo inoltra con $\text{TTL} = 1$) e scadrà sul secondo router, inducendo la trasmissione di un nuovo messaggio ICMP Time Exceeded che svela l'identità del secondo nodo (\emph{Hop 2}).
        \item Pacchetto con $\text{TTL} = 3$: svela il terzo nodo, e così via progressivamente.
    \end{itemize}
    \item \textbf{Conclusione della traccia}: la sequenza si arresta nel momento in cui il pacchetto raggiunge l'effettivo host terminale. In base all'implementazione (ad esempio l'invio di probe UDP su porte volutamente inesistenti superiori a 30000, oppure richieste ICMP Echo), l'host finale non risponderà con un Time Exceeded, bensì con un ICMP Port Unreachable (Type 3, Code 3) o con un ICMP Echo Reply (Type 0, Code 0), segnalando inequivocabilmente che la destinazione è stata raggiunta e che l'esplorazione dell'albero di instradamento è completata.
\end{enumerate}

% AA ---------------------------------------

% ==============================================================================
% CAPITOLO: FONDAMENTI DI NETWORKING: INDIRIZZAMENTO, INSTRADAMENTO E TRASPORTO
% ==============================================================================

\section{Livello Di Rete E Risoluzione Degli Indirizzi: Il Protocollo ARP}

Nel contesto delle reti a commutazione di pacchetto basate sullo stack TCP/IP, la corretta comunicazione tra nodi richiede una cooperazione sinergica tra differenti livelli di astrazione. Prima di analizzare i meccanismi di interconnessione su larga scala e l'incapsulamento avanzato tipico delle reti virtualizzate (come i tunnel overlay VXLAN adoperati nel cloud computing), è fondamentale consolidare i principi di risoluzione degli indirizzi locali e di instradamento inter-sottorete.

A livello di rete (Layer 3 del modello ISO/OSI), i nodi sono identificati da indirizzi logici, tipicamente espressi nel formato \textbf{IPv4} a $32\text{ bit}$. Tuttavia, la reale trasmissione delle trame a livello di collegamento dati (Layer 2) richiede l'impiego di indirizzi fisici hardware, noti come indirizzi \textbf{MAC} (\emph{Media Access Control}), lunghi $48\text{ bit}$. Poiché non sussiste alcun legame algebrico deterministico tra l'indirizzo logico IP assegnato a un'interfaccia e il suo indirizzo fisico MAC memorizzato sulla scheda di rete (\emph{Network Interface Card}, NIC), si rende indispensabile un meccanismo dinamico di mappatura.

Tale funzione è assolta dall'\textbf{Address Resolution Protocol} (\textbf{ARP}), specificato originariamente nella RFC~826, che opera a cavallo tra il livello di collegamento e il livello di rete per consentire la corretta incapsulazione delle trame Ethernet.

\begin{theorem}[Protocollo Di Risoluzione Degli Indirizzi (ARP)]
L'\textbf{Address Resolution Protocol} (ARP) è un protocollo dello stack TCP/IP deputato a identificare dinamicamente l'indirizzo a livello di collegamento dati (\emph{MAC address}) associato a un determinato indirizzo logico di rete (\emph{IP address}) all'interno del medesimo dominio di broadcast (Layer 2).
\end{theorem}

Nel panorama delle architetture IPv4, il protocollo ARP fornisce tre funzionalità essenziali:
\begin{enumerate}[noitemsep]
    \item \textbf{Risoluzione dinamica degli indirizzi}: determina il MAC address corrispondente a un indirizzo IP di destinazione noto.
    \item \textbf{Gestione della cache ARP}: memorizza localmente e temporaneamente le corrispondenze risolte all'interno di una tabella denominata \emph{ARP table} o \emph{ARP cache}, evitando l'emissione ripetitiva di richieste broadcast per ogni singolo pacchetto da trasmettere.
    \item \textbf{Rilevamento dei conflitti d'indirizzo}: attraverso il meccanismo di \emph{Gratuitous ARP}, un host può verificare la presenza di indirizzi IP duplicati all'interno del segmento di rete locale prima di attestarne l'assegnazione alla propria interfaccia.
\end{enumerate}

\subsection{Meccanismo Operativo: Richiesta Broadcast E Risposta Unicast}

Consideriamo due stazioni appartenenti alla medesima sottorete IP ($192.168.1.0/24$), indicate come Host~A e Host~B:
\begin{itemize}[noitemsep]
    \item \textbf{Host A}: IP $192.168.1.1/24$, MAC \texttt{3C-52-82-49-7E-9D}
    \item \textbf{Host B}: IP $192.168.1.2/24$, MAC \texttt{48-A4-72-1C-8F-4F}
\end{itemize}

Quando un'applicazione su Host~A desidera inviare un datagramma IP a Host~B, il sottosistema di rete di Host~A verifica anzitutto la propria \emph{ARP cache}. Se non è presente alcuna voce valida per l'indirizzo IP $192.168.1.2$, si innesca la seguente sequenza di risoluzione:

% [IMMAGINE DEL FILE DI SLIDE CON NOME "1) Ripasso_3"] Schema del protocollo ARP tra Host A e Host B con ARP Request broadcast e ARP Reply unicast — Slide 1 / Pagina 1 del PDF

\begin{enumerate}
    \item \textbf{Generazione dell'ARP Request}: Host~A costruisce un pacchetto contenente il proprio indirizzo IP e MAC sorgente, l'indirizzo IP target ($192.168.1.2$) e un campo MAC target lasciato indefinito ($00\text{:}00\text{:}00\text{:}00\text{:}00\text{:}00$). Questo pacchetto viene incapsulato in una trama di livello 2 con indirizzo MAC di destinazione impostato al valore di broadcast globale:
    \[
        \text{MAC}_{\text{dest}} = \texttt{FF:FF:FF:FF:FF:FF}
    \]
    Tale trama raggiunge indistintamente tutti gli host attestati sul dominio di collisione/broadcast locale (switch Ethernet).

    \item \textbf{Ricezione ed elaborazione da parte degli host}: Tutti i nodi ricevono la trama broadcast e la inoltrano al rispettivo stack di rete. Ciascun host analizza il campo \emph{Target IP}: i nodi il cui indirizzo IP differisce da quello interrogato scartano silenziosamente la richiesta. Host~B, riscontrando la corrispondenza con il proprio indirizzo IP locale, accetta la richiesta e aggiorna contestualmente la propria cache ARP con i parametri di Host~A.

    \item \textbf{Generazione dell'ARP Reply}: Host~B risponde costruendo un messaggio ARP Reply contenente come mittente il proprio IP ($192.168.1.2$) e il proprio MAC fisico (\texttt{48-A4-72-1C-8F-4F}). Poiché l'indirizzo MAC di Host~A è ora noto a Host~B, questo pacchetto viene inviato tramite una trama \textbf{unicast} indirizzata esclusivamente a \texttt{3C-52-82-49-7E-9D}.

    \item \textbf{Aggiornamento della cache}: Host~A riceve la replica unicast, estrae l'associazione indirizzo IP--MAC e la inserisce nella propria tabella ARP con un tempo di vita prefissato (\emph{aging timer}). A questo punto, il datagramma IP in sospeso può essere incapsulato nella trama Ethernet definitiva e inviato sul mezzo trasmissivo.
\end{enumerate}

\newpage

\section{Fondamenti Di Instradamento Inter-Subnet E Struttura Delle Tabelle Di Routing}

Il meccanismo di risoluzione ARP è intrinsecamente confinato al singolo dominio di broadcast Layer 2. Tuttavia, le reti telematiche globali e le infrastrutture cloud sono strutturate come un insieme interconnesso di sottoreti logiche (\emph{subnets}).

% [IMMAGINE DEL FILE DI SLIDE CON NOME "1) Ripasso_3"] Topologia di interconnessione tra sottoreti distribuite N ed M tramite router intermedi — Slide 2 / Pagina 2 del PDF

Ogni nodo è univocamente identificato da un indirizzo IP appartenente a una determinata sottorete. Quando un nodo situato nella sottorete $N$ deve comunicare con una destinazione attestata su una sottorete distante $M$, la trasmissione diretta a livello 2 non è praticabile: il pacchetto deve attraversare dispositivi intermedi specializzati, denominati \textbf{router} o \textbf{gateway}.

\subsection{Il Principio Di Instradamento Basato Su Rotte (\emph{Route-Based Forwarding})}

Un router è un dispositivo di livello 3 che esamina l'indirizzo IP di destinazione incapsulato nell'\emph{header} di ciascun datagramma e, consultando la propria base di dati locale denominata \textbf{IP Routing Table}, determina la migliore interfaccia di uscita e il nodo successivo verso cui indirizzare il traffico.

% [IMMAGINE DEL FILE DI SLIDE CON NOME "1) Ripasso_3"] Diagramma del Route-based Packet Forwarding da Gateway R1 attraverso nodi intermedi R2, R4 fino a R3 — Slide 3 / Pagina 3 del PDF

Il processo di inoltro basato su rotta (\emph{route-based packet forwarding}) segue una precisa catena logica:
\begin{enumerate}[noitemsep]
    \item Il nodo sorgente rileva, confrontando l'indirizzo di destinazione con la propria subnet mask, che l'host bersaglio non risiede sulla rete locale. Inoltra quindi la trama Layer 2 al proprio \textbf{default gateway} (router $R_1$).
    \item Il gateway esamina l'header IP, decrementa il campo \emph{Time to Live} (TTL), ricalcola il checksum e ricerca una corrispondenza nella propria tabella di instradamento.
    \item Il pacchetto viene instradato lungo una sequenza di nodi intermedi (hop successivi, es. $R_2$, $R_4$), ciascuno dei quali seleziona il percorso ritenuto ottimale.
    \item L'ultimo router della catena (il cosiddetto \emph{last-hop router}, $R_3$), direttamente connesso alla sottorete di destinazione $M$, esegue l'indirizzamento Layer 2 tramite risoluzione ARP locale e deposita la trama direttamente sulla scheda dell'host finale.
\end{enumerate}

\subsection{Anatomia Di Una Voce Di Routing}

Affinché un router possa prendere una decisione di inoltro coerente, ogni voce (\emph{routing entry}) presente nella tabella di instradamento deve contenere un nucleo informativo fondamentale:

% [IMMAGINE DEL FILE DI SLIDE CON NOME "1) Ripasso_3"] Struttura analitica dei campi di una tabella di instradamento con interfaccia GE0/0/0 e Next Hop — Slide 4 / Pagina 4 del PDF

\begin{itemize}
    \item \textbf{Destination}: identifica il prefisso di rete o sottorete di destinazione.
    \item \textbf{Mask}: la maschera di sottorete (espressa in formato decimale puntato o in notazione CIDR con barra diagonale), che definisce congiuntamente all'indirizzo la porzione identificativa della rete e la porzione riservata agli host.
    \item \textbf{Outbound Interface}: specifica l'interfaccia fisica o logica locale del router attraverso la quale il pacchetto deve essere trasmesso.
    \item \textbf{Next Hop}: specifica l'indirizzo IP del router successivo verso cui inviare il pacchetto affinché prosegua lungo il cammino stabilito.
\end{itemize}

\begin{theorem}[Necessità Congiunta Di Interfaccia E Next Hop]
Nei collegamenti aperti ad accesso multiplo (quali Ethernet con switch), l'indicazione della sola interfaccia di uscita è intrinsecamente insufficiente. Poiché su tale interfaccia possono coesistere molteplici dispositivi adiacenti, è indispensabile specificare l'indirizzo IP del \textbf{Next Hop} per permettere al protocollo ARP di risolvere il corretto MAC address del router successivo a cui affidare la trama.
\end{theorem}

\subsection{Organizzazione Della Tabella Di Routing IP}

La tabella di instradamento non racchiude l'intero grafo della topologia globale, bensì sintetizza un insieme selezionato di percorsi ottimali. I router apprendono le rotte tramite metodologie distinte, le filtrano in base a criteri di preferenza metrica e le installano nella memoria operativa.

% [IMMAGINE DEL FILE DI SLIDE CON NOME "1) Ripasso_3"] Topologia di rete con router R2 connesso a R1, R3, R4 e relativa IP Routing Table — Slide 5 / Pagina 5 del PDF

Prendendo a riferimento il nodo centrale $R_2$ collegato tramite l'interfaccia \texttt{GE0/0} alla rete $2.2.2.0/30$, \texttt{GE0/1} a $3.3.3.0/30$ e \texttt{GE0/2} a $1.1.1.0/30$, la tabella di routing presenta la seguente fisionomia:

\begin{table}[H]
\centering
\caption{Esempio di voci di instradamento registrate nel router $R_2$.}
\label{tab:routing_table_r2}
\begin{tabularx}{\textwidth}{l l X}
\toprule
\textbf{Destination/Mask} & \textbf{Next Hop} & \textbf{Outbound Interface} \\
\midrule
$11.0.0.0/8$ & $2.2.2.2$ & \texttt{GE0/0} \\
$13.0.0.0/8$ & $3.3.3.2$ & \texttt{GE0/1} \\
$14.0.0.0/8$ & $1.1.1.2$ & \texttt{GE0/2} \\
\dots & \dots & \dots \\
$1.1.1.0/30$ & $1.1.1.1$ & \texttt{GE0/2} \\
$1.1.1.1/32$ & $127.0.0.1$ & \texttt{GE0/2} \\
\bottomrule
\end{tabularx}
\end{table}

\subsection{Tipologie Di Generazione Delle Rotte}

Le informazioni di instradamento possono pervenire alla tabella di un dispositivo secondo tre modalità complementari:

% [IMMAGINE DEL FILE DI SLIDE CON NOME "1) Ripasso_3"] Confronto visivo tra rotte dirette, rotte statiche e rotte dinamiche — Slide 6 / Pagina 6 del PDF

\begin{enumerate}
    \item \textbf{Rotte Dirette (\emph{Direct Routes})}: generate in modo automatico dal sistema operativo di rete non appena un'interfaccia fisica o logica viene configurata con un indirizzo IP valido e portata nello stato operativo \emph{up}. Esse puntano direttamente alle sottoreti localmente connesse.
    \item \textbf{Rotte Statiche (\emph{Static Routes})}: introdotte manualmente dall'amministratore di rete mediante comandi espliciti di configurazione. Garantiscono pieno controllo e determinismo, non consumano banda né cicli CPU per scambi di segnalazione, ma presentano una scalabilità limitata e non reagiscono autonomamente a guasti topologici.
    \item \textbf{Rotte Dinamiche (\emph{Dynamic Routes})}: scoperte e aggiornate in tempo reale mediante protocolli di instradamento dinamico (quali OSPF, IS-IS e BGP). I router comunicano tra loro scambiandosi messaggi di stato, calcolano i cammini ottimali e riconvergono automaticamente in caso di variazione della topologia.
\end{enumerate}

\subsection{Interpretazione Della Tabella Di Routing Nel Piano Di Controllo}

Nelle piattaforme di rete commerciali (come i sistemi operativi Huawei VRP visualizzati nelle slide), l'ispezione della tabella di instradamento avviene mediante comandi CLI dedicati:

\begin{lstlisting}[language=bash, caption={Output del comando display ip routing-table.}]
<Quidway> display ip routing-table
Route Flags: R - relay, D - download to fib
----------------------------------------------------------------------
Routing Tables: Public
 Destinations : 6 Routes : 6
 Destination/Mask   Proto   Pre Cost Flags NextHop     Interface
 1.1.1.1/32         Static  60  0    D     0.0.0.0     NULL0
 2.2.2.2/32         Static  60  0    D     100.0.0.2   Vlanif100
 100.0.0.0/24       Direct  0   0    D     100.0.0.1   Vlanif100
 100.0.0.1/32       Direct  0   0    D     127.0.0.1   Vlanif100
 127.0.0.0/8        Direct  0   0    D     127.0.0.1   InLoopBack0
 127.0.0.1/32       Direct  0   0    D     127.0.0.1   InLoopBack0
\end{lstlisting}

% [IMMAGINE DEL FILE DI SLIDE CON NOME "1) Ripasso_3"] Esplosione analitica dei campi dell'output CLI display ip routing-table — Slide 7 / Pagina 7 del PDF

Ciascun campo dell'output possiede una precisa semantica:
\begin{itemize}
    \item \textbf{Destination/Mask}: indirizzo di rete e relativa maschera. L'indirizzo di sottorete corrispondente viene calcolato mediante l'operazione di \emph{AND bit a bit} tra l'indirizzo IP considerato e la maschera:
    \[
        \text{Subnet} = \text{IP} \;\land\; \text{Mask}
    \]
    Ad esempio, applicando la maschera $255.255.255.0$ all'indirizzo $1.1.1.1$, si ottiene l'indirizzo di rete $1.1.1.0$.
    \item \textbf{Proto (\emph{Protocol})}: specifica l'origine o il protocollo che ha generato la rotta (\texttt{Direct}, \texttt{Static}, \texttt{OSPF}, ecc.).
    \item \textbf{Pre (\emph{Preference})}: indica la distanza amministrativa o preferenza del protocollo. Qualora esistano percorsi concorrenti verso la stessa identica destinazione provenienti da protocolli differenti, il sistema installa nella tabella la rotta con il \textbf{valore di preferenza più basso} (la preferenza più alta corrisponde numericamente al valore minimo).
    \item \textbf{Cost (\emph{Metric})}: definisce il costo associato al cammino. Se due rotte verso la medesima destinazione presentano identico valore di preferenza, viene preferita quella a costo inferiore. Se anche i costi coincidono, è possibile attuare un bilanciamento di carico (\emph{Equal-Cost Multi-Path}, ECMP).
    \item \textbf{Flags}: attributi di inoltro del kernel di rete, ad esempio \texttt{D} indica che la rotta è stata scaricata con successo nella \textbf{FIB} (\emph{Forwarding Information Base}) a livello di piano dati hardware, mentre \texttt{R} indica una risoluzione iterativa (\emph{route relay}).
    \item \textbf{NextHop}: l'indirizzo del nodo successivo lungo il percorso.
    \item \textbf{Interface}: l'interfaccia locale di uscita attraverso la quale il traffico deve essere irradiato.
\end{itemize}

I valori di preferenza predefiniti variano a seconda del costruttore; i riferimenti standard tipici delle slide sono riassunti nella seguente tabella:

\begin{table}[H]
\centering
\caption{Valori di default della Route Preference.}
\label{tab:route_preference}
\begin{tabular}{llr}
\toprule
\textbf{Protocollo} & \textbf{Tipologia di Rotta} & \textbf{Preferenza di Default} \\
\midrule
Direct & Rotta locale direttamente connessa & $0$ \\
OSPF & Rotta interna OSPF & $10$ \\
Static & Rotta statica configurata & $60$ \\
OSPF & Rotta esterna importata in OSPF & $150$ \\
\bottomrule
\end{tabular}
\end{table}

\subsection{Meccanismo Di Selezione: Regola Del Prefisso Più Lungo (\emph{Longest Prefix Match})}

Quando un pacchetto giunge al piano di inoltro, l'indirizzo IP di destinazione viene confrontato con le voci della tabella. Spesso un medesimo indirizzo IP può soddisfare contemporaneamente molteplici regole di instradamento aventi lunghezze di prefisso differenti.

\begin{theorem}[Regola Del Prefisso Più Lungo (\emph{Longest Match Rule})]
Dato un insieme di rotte corrispondenti all'indirizzo IP di destinazione, il router seleziona invariabilmente la rotta dotata della \textbf{maschera di sottorete più lunga} (ovvero con il maggior numero di bit di rete continui a $1$), in quanto tale prefisso descrive in modo più specifico e circoscritto la locazione del nodo target.
\end{theorem}

% [IMMAGINE DEL FILE DI SLIDE CON NOME "1) Ripasso_3"] Rappresentazione bit a bit dell'algoritmo di Longest Matching per l'indirizzo 172.16.2.1 — Slide 10 / Pagina 10 del PDF

Consideriamo un pacchetto recante come destinazione l'indirizzo $172.16.2.1$, da confrontare con tre voci di instradamento presenti nella tabella:
\begin{align*}
    \text{IP Destinazione: } & 172.16.2.1 \implies \texttt{10101100.00010000.00000010.00000001}_2 \\
    \text{Entry 1: } & 172.16.1.0/24 \implies \texttt{10101100.00010000.00000001.xxxxxxxx}_2 \quad (\text{Mismatch al 24-esimo bit}) \\
    \text{Entry 2: } & 172.16.2.0/24 \implies \texttt{10101100.00010000.00000010.xxxxxxxx}_2 \quad (\text{Match su 24 bit}) \\
    \text{Entry 3: } & 172.16.0.0/16 \implies \texttt{10101100.00010000.xxxxxxxx.xxxxxxxx}_2 \quad (\text{Match su 16 bit})
\end{align*}

Sia l'Entry~2 sia l'Entry~3 comprendono l'indirizzo $172.16.2.1$. Tuttavia, poiché la maschera dell'Entry~2 è estesa su $24\text{ bit}$ mentre l'Entry~3 su soli $16\text{ bit}$, il router elegge l'Entry~2 come cammino vincente (\emph{Longest Match}).

% [IMMAGINE DEL FILE DI SLIDE CON NOME "1) Ripasso_3"] Esempi applicativi di Longest Matching per gli indirizzi 192.168.2.2 e 192.168.3.2 sul router RTA — Slide 11 e 12 / Pagine 11 e 12 del PDF

Analogamente, consideriamo la tabella del router RTA:
\begin{itemize}[noitemsep]
    \item $192.168.0.0/16 \to \text{Next Hop } 10.1.1.2$
    \item $192.168.2.0/24 \to \text{Next Hop } 20.1.1.2$
    \item $192.168.3.0/24 \to \text{Next Hop } 30.1.1.2$
\end{itemize}
Se il pacchetto è diretto a $192.168.2.2$, esso soddisfa contemporaneamente sia il prefisso $/16$ sia $/24$; viene preferita la rotta $192.168.2.0/24$ inoltrando a $20.1.1.2$. Se invece il pacchetto è destinato a $192.168.3.2$, esso matcha la voce $/24$ specifica inoltrando a $30.1.1.2$.

\subsection{Esempio Completo Di Inoltro Hop-By-Hop Tra Sottoreti}

Analizziamo il tragitto compiuto da un pacchetto originato nella sottorete $10.0.1.0/24$ e diretto all'host $40.0.1.2$ attestato sulla sottorete $40.0.1.0/24$, attraversando la dorsale formata dai router $R_1$, $R_2$ e $R_3$.

% [IMMAGINE DEL FILE DI SLIDE CON NOME "1) Ripasso_3"] Flusso sequenziale end-to-end con routing tables complete di R1, R2 e R3 per la rete 40.0.1.0/24 — Slide 13 / Pagina 13 del PDF

\begin{enumerate}
    \item \textbf{Ingresso nel Gateway $R_1$}: Il gateway $R_1$ riceve la trama sulla propria interfaccia locale \texttt{GE0/1} ($10.0.1.1$). Analizzando la propria tabella di routing per il prefisso $40.0.1.0/24$, riscontra:
    \[
        \text{Dest: } 40.0.1.0/24 \implies \text{Next Hop: } 20.0.1.2, \quad \text{Outbound Intf: } \texttt{GE0/0}
    \]
    $R_1$ incapsula il pacchetto in una trama Layer 2 indirizzata al MAC di $R_2$ ($20.0.1.2$) e lo trasmette da \texttt{GE0/0}.

    \item \textbf{Transito attraverso $R_2$}: Il router intermedio $R_2$ riceve il pacchetto sulla porta \texttt{GE0/0} ($20.0.1.2$), esamina l'header IP ed effettua la ricerca per $40.0.1.0/24$:
    \[
        \text{Dest: } 40.0.1.0/24 \implies \text{Next Hop: } 30.0.1.2, \quad \text{Outbound Intf: } \texttt{GE0/1}
    \]
    $R_2$ reinvia il pacchetto attraverso \texttt{GE0/1} verso l'interfaccia $30.0.1.2$ di $R_3$.

    \item \textbf{Inoltro finale su $R_3$}: Il router $R_3$ riceve il pacchetto sulla propria porta \texttt{GE0/0} ($30.0.1.2$). La consultazione della tabella rivela che la sottorete $40.0.1.0/24$ corrisponde a una rotta diretta:
    \[
        \text{Dest: } 40.0.1.0/24 \implies \text{Next Hop: } 40.0.1.1, \quad \text{Outbound Intf: } \texttt{GE0/1}
    \]
    Poiché l'interfaccia \texttt{GE0/1} è direttamente affacciata sul segmento bersaglio, $R_3$ non inoltra a un router ulteriore, ma esegue una risoluzione ARP locale per determinare il MAC dell'host $40.0.1.2$ e consegna direttamente la trama al destinatario finale.
\end{enumerate}

\newpage

\section{Architettura Dei Sistemi Autonomi E Protocolli Di Routing Dinamico}

L'infrastruttura globale di Internet e le grandi reti \emph{enterprise} o \emph{data center cloud} sono organizzate su base gerarchica per prevenire l'esplosione delle dimensioni delle tabelle di routing.

\begin{theorem}[Sistema Autonomo (\emph{Autonomous System}, AS)]
Come formalizzato nella RFC~1030 e ripreso nella RFC~1930, un \textbf{Sistema Autonomo} (AS) è costituito da un insieme connesso di prefissi IP gestiti da uno o più operatori di rete sotto il controllo di un'amministrazione tecnica unificata, che esprime una \textbf{politica di instradamento singola e chiaramente definita} (\emph{Single and Clearly Defined Routing Policy}).
\end{theorem}

% [IMMAGINE DEL FILE DI SLIDE CON NOME "1) Ripasso_3"] Rappresentazione concettuale di un Autonomous System con magliatura interna di router — Slide 14 / Pagina 14 del PDF

Sebbene la formulazione originaria prevedesse l'utilizzo di un unico protocollo di instradamento per AS, le reti moderne supportano la coesistenza di molteplici protocolli coordinati tramite l'iniezione bidirezionale delle rotte (\emph{route redistribution}). La politica di routing (\emph{routing policy}) rappresenta l'insieme di vincoli operativi e contrattuali che determinano la propagazione e la manipolazione dei prefissi all'interno e ai confini dell'AS.

\subsection{Instradamento Interno Ed Esterno: Nodi Di Bordo E Nodi Core}

Nella topologia di Internet si identificano due categorie funzionali distinte di router:

% [IMMAGINE DEL FILE DI SLIDE CON NOME "1) Ripasso_3"] Architettura di Internet Routing: Inter-AS border router ed Intra-AS interior router — Slide 15 / Pagina 15 del PDF

\begin{itemize}
    \item \textbf{Intra-AS Interior Router}: apparati confinati all'interno del perimetro dell'AS, dedicati all'inoltro ad alta efficienza del traffico interno.
    \item \textbf{Inter-AS Border Router} (o \emph{Exterior Gateway}): nodi di confine posizionati al margine dell'AS, responsabili dell'interconnessione con altri domini autonomi e dell'applicazione delle policy commerciali ed economiche di transito o \emph{peering}.
\end{itemize}

\subsection{Tassonomia Dei Protocolli Di Instradamento Dinamico}

I protocolli dinamici si classificano secondo due assi ortogonali: il raggio d'azione applicativo e l'algoritmo matematico computazionale impiegato.

% [IMMAGINE DEL FILE DI SLIDE CON NOME "1) Ripasso_3"] Tassonomia dei protocolli di routing dinamico per ambito applicativo e per algoritmo — Slide 16 / Pagina 16 del PDF

\begin{enumerate}
    \item \textbf{Classificazione per ambito applicativo (\emph{Application Scope})}:
    \begin{itemize}
        \item \textbf{Interior Gateway Protocol} (\textbf{IGP}): operano all'interno del singolo AS con l'obiettivo primario di ottimizzare le prestazioni e convergere rapidamente (es. RIP, OSPF, IS-IS).
        \item \textbf{Exterior Gateway Protocol} (\textbf{EGP}): operano tra Sistemi Autonomi distinti; lo standard universale è rappresentato dal protocollo \textbf{BGP} (\emph{Border Gateway Protocol}), focalizzato sull'applicazione di policy amministrative complesse piuttosto che sulla mera convergenza a costo minimo.
    \end{itemize}

    \item \textbf{Classificazione per principio algoritmico}:
    \begin{itemize}
        \item \textbf{Distance-Vector}: basati sull'algoritmo di Bellman-Ford (es. RIP). I router inviano periodicamente l'intera tabella ai soli vicini diretti; la conoscenza topologica è indiretta (``\emph{routing by rumor}'').
        \item \textbf{Link-State}: basati sull'algoritmo di Dijkstra (es. OSPF, IS-IS). I router distribuiscono informazioni sullo stato dei propri collegamenti a tutti i nodi dell'area, costruendo una mappa topologica completa e identica in ciascun apparato.
        \item \textbf{Path-Vector}: evoluzione del modello distance-vector adottata da BGP, in cui la metrica include l'intera sequenza di AS attraversati (\emph{AS-Path}), consentendo la prevenzione dei loop su scala planetaria.
    \end{itemize}
\end{enumerate}

\subsection{Meccanismo Operativo Dei Protocolli Link-State}

I protocolli basati sullo stato dei collegamenti (OSPF, IS-IS) articolano il proprio funzionamento in quattro fasi sequenziali rigidamente strutturate:

% [IMMAGINE DEL FILE DI SLIDE CON NOME "1) Ripasso_3"] Fasi sequenziali del protocollo Link-State: Neighbor Setup, Flooding LSDB, SPF Computation e Caricamento RIB — Slide 17 / Pagina 17 del PDF

\begin{enumerate}
    \item \textbf{Fase 1 -- Instaurazione del vicinato (\emph{Neighbor Relationship Setup})}: I router adiacenti scoprono la presenza reciproca scambiandosi pacchetti di controllo leggeri (\emph{Hello packets}) e stabiliscono una relazione formale di adiacenza una volta verificata la compatibilità dei parametri configurati.

    \item \textbf{Fase 2 -- Scambio dello stato dei link e sincronizzazione del database (\emph{LSDB Synchronization})}: Ciascun nodo genera descrizioni dettagliate delle proprie interfacce locali e dei vicini connessi (\emph{Link-State Advertisements}, LSA, in OSPF). Tali messaggi vengono distribuiti a tutti i router mediante inondazione controllata (\emph{reliable flooding}). Al termine dello scambio, tutti i partecipanti condividono un medesimo \textbf{Link-State Database} (\textbf{LSDB}), che funge da grafo geometrico completo dell'area.

    \item \textbf{Fase 3 -- Calcolo del cammino ottimale (\emph{Path Computation})}: Ciascun router esegue in locale l'algoritmo di Dijkstra (\emph{Shortest Path First}, SPF) posizionando se stesso come radice di un albero di copertura minimo (\emph{Shortest Path Tree}, SPT).

    \item \textbf{Fase 4 -- Generazione e caricamento delle rotte nella RIB}: I percorsi a costo minimo ricavati dall'albero vengono convertiti in voci di instradamento e caricati nella tabella di instradamento principale (\textbf{Routing Information Base}, \textbf{RIB}), da cui verranno poi trasferiti nella FIB per l'inoltro in hardware.
\end{enumerate}

\newpage

\section{Livello Di Trasporto: I Protocolli TCP E UDP}

Il livello di trasporto (Layer 4) instaura una comunicazione logica \emph{end-to-end} tra processi applicativi eseguiti su host remoti. Mentre il livello di rete garantisce unicamente la consegna \emph{host-to-host} su base \emph{best-effort}, il livello di trasporto introduce il concetto di \textbf{multiplexing} e \textbf{demultiplexing} attraverso l'adozione delle porte software.

\begin{theorem}[Spazio Delle Porte Di Trasporto]
Un punto terminale di comunicazione a livello 4 è identificato dalla coppia composta dall'indirizzo IP e dal numero di \textbf{porta} (lungo $16\text{ bit}$, con valori ammissibili da $0$ a $65535$). L'intervallo si divide convenzionalmente in:
\begin{itemize}[noitemsep]
    \item \textbf{Well-Known Ports} ($0$--$1023$): riservate a servizi di sistema standard (es. HTTP porta $80$, Telnet porta $23$).
    \item \textbf{Ephemeral Ports} ($1024$--$65535$): allocate dinamicamente dal sistema operativo client all'apertura delle sessioni effimere.
\end{itemize}
\end{theorem}

% [IMMAGINE DEL FILE DI SLIDE CON NOME "1) Ripasso_3"] Multiplexing tramite porte sorgente e destinazione tra Web browser e server HTTP/Telnet — Slide 21 / Pagina 21 del PDF

A livello di trasporto emergono due protocolli cardine:
\begin{itemize}[noitemsep]
    \item \textbf{UDP} (\emph{User Datagram Protocol}): leggero, privo di connessione (\emph{connectionless}), non orientato alla conferma, privo di controllo di flusso o congestione.
    \item \textbf{TCP} (\emph{Transmission Control Protocol}): orientato alla connessione (\emph{connection-oriented}), affidabile, con recupero delle perdite, controllo del flusso e regolazione della congestione.
\end{itemize}

\subsection{Confronto Degli Header Di Livello 4}

La complessità funzionale di TCP si riflette direttamente nella conformazione del relativo \emph{header} rispetto alla controparte UDP.

% [IMMAGINE DEL FILE DI SLIDE CON NOME "1) Ripasso_3"] Struttura comparata degli header TCP (20 byte base) e UDP (8 byte base) — Slide 18 / Pagina 18 del PDF

\subsubsection{Formato Dell'Header UDP ($8\text{ Byte}$)}
L'header UDP è elementare e richiede solo quattro campi da $16\text{ bit}$:
\begin{itemize}[noitemsep]
    \item \textbf{Source Port} ($16\text{ bit}$): identifica l'applicazione mittente.
    \item \textbf{Destination Port} ($16\text{ bit}$): identifica l'applicazione ricevente.
    \item \textbf{Length} ($16\text{ bit}$): specifica la lunghezza complessiva dell'header e del payload UDP (lunghezza minima $8\text{ byte}$, massima teorica $65535\text{ byte}$).
    \item \textbf{Checksum} ($16\text{ bit}$): codice di controllo d'errore calcolato sull'header, sul payload e su uno pseudo-header IP.
\end{itemize}

\subsubsection{Formato Dell'Header TCP ($20\text{--}60\text{ Byte}$)}
In assenza di opzioni, l'header TCP presenta una dimensione base di $20\text{ byte}$ ($160\text{ bit}$), strutturata nei seguenti campi:
\begin{itemize}
    \item \textbf{Source Port} e \textbf{Destination Port} ($16\text{ bit}$ ciascuna): identificano i processi di origine e destinazione.
    \item \textbf{Sequence Number} ($32\text{ bit}$): definisce il numero di sequenza assegnato al primo byte di dati trasportato nel segmento corrente all'interno del flusso logico continuo.
    \item \textbf{Acknowledgment Number} ($32\text{ bit}$): valido se il flag ACK è asserito, indica il numero di sequenza del byte successivo che il ricevente attende di ricevere. Funge da riscontro cumulativo positivo per tutti i byte precedenti.
    \item \textbf{Header Length} (\emph{Data Offset}, $4\text{ bit}$): esprime la lunghezza totale dell'header misurata in parole da $32\text{ bit}$ ($4\text{ byte}$). Il valore minimo è $5$ (equivalente a $5 \times 4 = 20\text{ byte}$), mentre il valore massimo ammissibile è $15$ ($60\text{ byte}$, consentendo fino a $40\text{ byte}$ di opzioni).
    \item \textbf{Reserved} ($6\text{ bit}$): campo riservato per usi futuri, da impostare rigorosamente a zero.
    \item \textbf{Control Bits} (\emph{Flags}, $6$ bit originari più estensioni): codificano lo stato e lo scopo semantico del segmento.
    \item \textbf{Window Size} ($16\text{ bit}$): esprime l'ampiezza della finestra di ricezione (\emph{receive window}, massimo $65535\text{ byte}$ senza opzione di scala), utilizzata per il controllo di flusso.
    \item \textbf{Checksum} ($16\text{ bit}$): campo obbligatorio per l'integrità end-to-end, calcolato includendo un pseudo-header IP da $12\text{ byte}$.
    \item \textbf{Urgent Pointer} ($16\text{ bit}$): valido solo se il flag URG è asserito, segnala l'offset finale dei dati urgenti collocati all'inizio del payload.
    \item \textbf{Options} ($0$--$40\text{ byte}$): parametri addizionali negoziabili (es. MSS, Window Scale, SACK, Timestamp).
\end{itemize}

\subsection{Analisi Dei Flag Di Controllo E Gestione ECN}

I bit di controllo determinano le transizioni di stato del motore TCP:
\begin{itemize}
    \item \textbf{URG}: indica che il campo \emph{Urgent Pointer} è valido e che il segmento contiene dati prioritari da elaborare fuori banda.
    \item \textbf{ACK}: certifica che il campo \emph{Acknowledgment Number} trasporta un valore valido di riscontro.
    \item \textbf{PSH}: istruisce il ricevente a svuotare immediatamente il buffer di ricezione passando i dati al livello applicativo, senza attendere il riempimento di segmenti ulteriori.
    \item \textbf{RST}: forza l'azzeramento o l'abbattimento immediato di una connessione a fronte di errori irreversibili o pacchetti inattesi.
    \item \textbf{SYN}: sincronizza i numeri di sequenza all'atto dell'apertura iniziale della connessione.
    \item \textbf{FIN}: notifica l'intenzione del mittente di chiudere la propria direzione di trasmissione (avvio della terminazione della sessione).
\end{itemize}

% [IMMAGINE DEL FILE DI SLIDE CON NOME "1) Ripasso_3"] Tabella dei flag per il controllo della congestione esplicita ECN (CWR, ECE, ECT, CE) — Slide 18 / Pagina 18 del PDF

L'estensione standard RFC~3168 ha riutilizzato bit del campo di controllo per introdurre l'\textbf{Explicit Congestion Notification} (\textbf{ECN}), un meccanismo che permette ai router intermedi di segnalare la saturazione delle code senza dover scartare pacchetti:
\begin{itemize}
    \item \textbf{ECE} (\emph{ECN-Echo}): asserito durante il three-way handshake dal ricevitore per confermare il supporto a ECN; durante il normale trasferimento, indica che il ricevitore ha ricevuto un pacchetto IP con marcatura di congestione (\emph{Congestion Experienced}, CE).
    \item \textbf{CWR} (\emph{Congestion Window Reduced}): asserito dal mittente per confermare al ricevitore di aver ricevuto l'avviso ECE e di aver dimezzato la propria finestra di trasmissione.
\end{itemize}

\newpage

\section{Ciclo Di Vita Delle Connessioni TCP Ed Elaborazione Del Flusso Di Dati}

\subsection{Instaurazione Della Connessione: Il Three-Way Handshake}

Prima di scambiare dati applicativi, le due entità TCP devono sincronizzare i rispettivi numeri di sequenza iniziali (\textbf{ISN}, \emph{Initial Sequence Number}) e negoziare i buffer tramite una procedura in tre passi:

% [IMMAGINE DEL FILE DI SLIDE CON NOME "1) Ripasso_3"] Sequenza temporale del Three-Way Handshake TCP tra PC1 e PC2 — Slide 22 / Pagina 22 del PDF

\begin{enumerate}
    \item \textbf{Passo 1 (SYN da PC1 a PC2)}:
    Il client PC1 seleziona uno pseudocasuale $a$ come proprio ISN e invia un segmento con flag impostato a $\text{SYN}=1$:
    \[
        \text{Seq} = a, \quad \text{Ack} = 0 \quad (\text{ACK bit non attivo})
    \]

    \item \textbf{Passo 2 (SYN-ACK da PC2 a PC1)}:
    Il server PC2 alloca le risorse, sceglie il proprio ISN indipendente $b$, asserisce contestualmente i flag $\text{SYN}=1$ e $\text{ACK}=1$, confermando l'ISN del client tramite l'incremento unitario:
    \[
        \text{Seq} = b, \quad \text{Ack} = a + 1
    \]

    \item \textbf{Passo 3 (ACK da PC1 a PC2)}:
    Il client risponde asserendo il solo flag $\text{ACK}=1$. Il segmento reca:
    \[
        \text{Seq} = a + 1, \quad \text{Ack} = b + 1
    \]
    Alla ricezione di questo pacchetto, la connessione è formalmente stabilita in entrambi i versi (\emph{ESTABLISHED}).
\end{enumerate}

\subsection{Numerazione Del Flusso Di Byte E Calcolo Degli Acknowledgment}

A differenza dei protocolli orientati al messaggio, TCP modella il canale come un flusso continuo non strutturato di singoli byte (\textbf{Byte Stream}). Ciascun byte trasmesso riceve un identificatore progressivo:

% [IMMAGINE DEL FILE DI SLIDE CON NOME "1) Ripasso_3"] Calcolo sequenziale dei campi Sequence Number ed Ack Number durante il trasferimento dati tra PC1 e PC2 — Slide 23 / Pagina 23 del PDF

\begin{theorem}[Relazione Fondamentale Di Avanzamento Sequenza]
Dato un segmento TCP trasmesso con numero di sequenza iniziale $\text{Seq} = S$ e recante un carico utile applicativo di dimensione $L\text{ byte}$ ($L = \text{Payload Length}$), l'indice dell'ultimo byte del segmento è pari a $S + L - 1$. Conseguentemente, il valore del riscontro atteso dal ricevente è calcolato analiticamente come:
\[
    \boxed{\text{Ack} = \text{Seq} + L}
\]
\end{theorem}

Analizziamo l'interazione riportata nel diagramma operativo tra PC1 ($1.1.1.1$) e PC2 ($2.2.2.2$):
\begin{enumerate}[noitemsep]
    \item PC1 invia il primo segmento dati contenente $12\text{ byte}$ a partire da $\text{Seq} = a + 1$. Il segmento copre l'intervallo di byte $[a+1, \, a+12]$. Il riscontro inviato è $\text{Ack} = b + 1$.
    \item PC2 conferma la ricezione dei $12\text{ byte}$ inviando un segmento di riscontro puro ($L = 0\text{ byte}$) con:
    \[
        \text{Seq} = b + 1, \quad \text{Ack} = (a + 1) + 12 = a + 13
    \]
    \item PC1 trasmette successivamente un blocco di $66\text{ byte}$. Il numero di sequenza coincide esattamente con il riscontro ricevuto: $\text{Seq} = a + 13$. Poiché PC2 non ha inviato dati applicativi nel frattempo, il valore di acknowledgment del segmento inviato da PC1 rimane immutato su $\text{Ack} = b + 1$.
    \item PC2 riscontra il blocco emettendo:
    \[
        \text{Ack} = (a + 13) + 66 = a + 79
    \]
\end{enumerate}

\subsection{Chiusura Della Connessione: Il Four-Way Handshake}

Poiché il canale TCP è intrinsecamente \textbf{full-duplex}, ciascuna delle due direzioni di comunicazione è indipendente e deve essere terminata separatamente. La chiusura ordinata avviene tramite quattro segmenti (\emph{Four-Way Handshake}):

% [IMMAGINE DEL FILE DI SLIDE CON NOME "1) Ripasso_3"] Sequenza di abbattimento connessione Four-Way Handshake con gestione dello stato Half-Closed — Slide 30 / Pagina 30 del PDF

\begin{enumerate}
    \item \textbf{FIN da PC1 a PC2}: PC1 esaurisce i dati da trasmettere e invia un segmento con flag $\text{FIN}=1$ ($\text{Seq}=101, \text{Ack}=301$). Da questo istante, PC1 non può più originare dati, ma resta in ascolto per ricevere eventuale traffico residuo.
    \item \textbf{ACK da PC2 a PC1}: PC2 riscontra la richiesta con $\text{Ack}=102$. La connessione entra in uno stato di chiusura parziale (\textbf{Half-Closed}): il canale da PC1 a PC2 è disattivato, mentre il canale inverso rimane operativo.
    \item \textbf{FIN da PC2 a PC1}: Quando anche l'applicazione su PC2 decide di terminare la sessione, invia un proprio segmento con flag $\text{FIN}=1$ ($\text{Seq}=301, \text{Ack}=102$).
    \item \textbf{ACK finale da PC1 a PC2}: PC1 conferma la chiusura inviando $\text{Ack}=302$. Al completamento dei timeout di salvaguardia (\emph{TIME-WAIT}), i descrittori e le strutture di memoria su entrambi i nodi vengono rilasciati.
\end{enumerate}

\newpage

\section{Meccanismi Di Controllo Del Flusso E Della Congestione In TCP}

Un protocollo di trasporto robusto deve operare due forme distinte di controllo della velocità trasmissiva:
\begin{itemize}[noitemsep]
    \item \textbf{Controllo del Flusso (\emph{Flow Control})}: meccanismo \emph{end-to-end} che impedisce al mittente rapido di saturare i buffer di memoria del ricevitore lento.
    \item \textbf{Controllo della Congestione (\emph{Congestion Control})}: meccanismo distribuito che impedisce alla sorgente di iniettare nella rete una quantità di traffico eccedente la capacità di smaltimento dei router intermedi.
\end{itemize}

\subsection{Controllo Del Flusso Tramite Finestra Scorrevole (\emph{Sliding Window})}

Il ricevitore alloca in memoria un'area dedicata alla memorizzazione dei segmenti in ingresso, indicata come \textbf{RcvBuffer}. Lo spazio effettivamente libero e disponibile ad accogliere nuovi byte è definito come finestra di ricezione (\textbf{RcvWindow}):

% [IMMAGINE DEL FILE DI SLIDE CON NOME "1) Ripasso_3"] Modello a buffer del ricevitore con partizione tra dati in attesa per l'applicazione e spare room (RcvWindow) — Slide 25 / Pagina 25 del PDF

\[
    \text{RcvWindow} = \text{RcvBuffer} - \left( \text{LastByteReceived} - \text{LastByteRead} \right)
\]
Il valore istantaneo di $\text{RcvWindow}$ viene inserito nel campo \emph{Window} di ciascun segmento di ritorno trasmesso verso il mittente.

% [IMMAGINE DEL FILE DI SLIDE CON NOME "1) Ripasso_3"] Meccanismo a finestra scorrevole: trasmissione di 3 segmenti e contrazione della finestra ad 1 — Slide 24 / Pagina 24 del PDF

L'evoluzione dinamica della finestra si articola come segue:
\begin{enumerate}[noitemsep]
    \item Durante il three-way handshake iniziale, entrambe le stazioni dichiarano la dimensione dei propri buffer (es. $\text{win}=3$ segmenti).
    \item Il mittente trasmette pacchetti fino a saturare il credito concesso dalla finestra ($\text{seq}=101, 102, 103$).
    \item Se l'applicazione ricevente non preleva i dati con sufficiente rapidità, lo spazio disponibile si riduce; il ricevente emette un ACK con valore di finestra ridotto ($\text{Ack}=104, \text{win}=1$).
    \item Il mittente adegua istantaneamente il proprio limite superiore di trasmissione al nuovo credito dichiarato.
\end{enumerate}

\subsubsection{Gestione Della Finestra Zero (\emph{Zero-Window Probing})}
Nel caso estremo in cui il buffer del ricevitore sia completamente saturo, quest'ultimo invia un segmento con $\text{Window} = 0$. Il mittente arresta immediatamente l'inoltro di nuovi dati.

Tuttavia, per scongiurare una situazione di stallo permanente dovuta allo smarrimento dell'ACK con cui il ricevitore riapre successivamente la finestra, il mittente attiva un timer di persistenza (\emph{persist timer}, tipicamente con cadenza di ritrasmissione a $2\text{ minuti}$). Alla scadenza del timer, il mittente invia un segmento di test da $1\text{ byte}$ (\textbf{Zero-Window Probe}); il ricevitore è obbligato a replicare confermando il proprio stato corrente e annunciando la nuova disponibilità di finestra.

\subsection{Finestra Trasmissiva Effettiva E Throughput Massimo}

La reale capacità di iniezione di dati nel canale è governata contemporaneamente dal vincolo del ricevitore ($\text{RcvWindow}$) e dal vincolo della rete, espresso dalla finestra di congestione (\textbf{CongWin} o $\text{cwnd}$):

% [IMMAGINE DEL FILE DI SLIDE CON NOME "1) Ripasso_3"] Definizione analitica della Transmission Window come minimo tra CongWin e RcvWindow — Slide 26 / Pagina 26 del PDF

\begin{theorem}[Finestra Di Trasmissione Effettiva (\emph{TrWindow})]
La dimensione della finestra effettiva di trasmissione (\textbf{TrWindow}) entro cui il mittente può inoltrare byte consecutivi senza attendere riscontro è data dal valore minimo tra la finestra di congestione interna e la finestra pubblicizzata dal ricevitore:
\[
    \boxed{\text{TrWindow} = \min\left(\text{CongWin}, \, \text{RcvWindow}\right)}
\]
\end{theorem}

Assumendo condizioni ideali stazionarie in cui le perdite siano nulle e l'intero credito della finestra venga trasmesso in un intervallo pari al tempo di andata e ritorno (\emph{Round Trip Delay}, $\text{RTD}$ o $\text{RTT}$), il throughput massimo teorico conseguibile dalla connessione è espresso da:
\[
    \boxed{\text{Throughput}_{\max} = \frac{\text{TrWindow}}{\text{RTD}} \quad \left[\frac{\text{byte}}{\text{sec}}\right]}
\]

\subsection{Struttura Del Buffer Di Trasmissione}

Nel trasmettitore, lo spazio sequenziale della finestra di congestione è ripartito in quattro regioni dinamiche distinte:

% [IMMAGINE DEL FILE DI SLIDE CON NOME "1) Ripasso_3"] Ripartizione a blocchi colorati della finestra del mittente: sent acked, sent not acked, usable not yet sent, not usable — Slide 27 / Pagina 27 del PDF

\begin{itemize}[noitemsep]
    \item \textbf{Already ack'ed}: byte già trasmessi e per i quali è già pervenuto un riscontro valido (regione a monte di \texttt{send\_base}).
    \item \textbf{Sent, not yet ack'ed}: byte già inviati sul mezzo fisico ma ancora in attesa di ricezione di ACK (intervallo compreso tra \texttt{send\_base} e \texttt{nextseqnum}).
    \item \textbf{Usable, not yet sent}: byte per i quali esiste credito trasmissivo disponibile all'interno della finestra corrente, ma che non sono ancora stati materialmente iniettati nella rete.
    \item \textbf{Not usable}: byte del flusso applicativo situati oltre il limite superiore di $\text{CongWin}$, la cui trasmissione è interdetta fino all'avanzamento della finestra scorrevole.
\end{itemize}

\subsection{Algoritmi Di Controllo Della Congestione: Slow Start E Congestion Avoidance}

In assenza di segnalazione ECN esplicita, TCP stima lo stato di congestione dei collegamenti intermedi in modo inferenziale, interpretando la mancata ricezione dei riscontri o lo scadere dei timer come indice di saturazione dei nodi.

\subsubsection{Fase 1: Slow Start}
All'avvio della connessione, il trasmettitore ignora quale sia l'effettiva capacità di banda del percorso. L'algoritmo di \textbf{Slow Start} avvia la trasmissione partendo da una finestra minima e incrementandola in modo esponenziale per ciascun RTT:

% [IMMAGINE DEL FILE DI SLIDE CON NOME "1) Ripasso_3"] Diagramma temporale della crescita esponenziale della finestra durante la fase di TCP Slow Start — Slide 28 / Pagina 28 del PDF

L'algoritmo opera secondo il seguente pseudocodice concettuale:

\begin{algorithm}[H]
\caption{Algoritmo Di TCP Slow Start}
\Dati{Inizializzazione sessione o ripartenza da timeout}
\Risultato{Crescita esponenziale di $\text{CongWin}$}
$\text{CongWin} \gets 1$\;
\Mentre{$\text{nessun evento di perdita} \;\mathbf{e}\; (\text{CongWin} \le \text{threshold})$}{
    \Per{ciascun segmento riscontrato con successo}{
        $\text{CongWin} \gets \text{CongWin} + 1$\;
    }
}
\end{algorithm}

Poiché per ciascun ACK ricevuto la finestra si amplia di una unità, ad ogni tempo di propagazione completo $\text{RTD}$ la dimensione della finestra raddoppia sistematicamente ($1 \to 2 \to 4 \to 8 \dots$), determinando una rapida crescita esponenziale.

\subsubsection{Fase 2: Congestion Avoidance}
Non appena $\text{CongWin}$ supera il valore di soglia prefissato (\textbf{threshold} o \textbf{ssthresh}), il protocollo passa alla modalità di accrescimento conservativo denominata \textbf{Congestion Avoidance}:

% [IMMAGINE DEL FILE DI SLIDE CON NOME "1) Ripasso_3"] Grafico cartesiano della finestra di congestione con transizione da Slow Start a Congestion Avoidance e dimezzamento della soglia — Slide 29 / Pagina 29 del PDF

In questa fase, l'incremento non avviene più per singolo segmento riscontrato, ma la finestra si espande di un solo segmento per ciascun intero blocco di segmenti ($W = \text{CongWin}$) confermati:
\begin{align*}
    \Delta \text{CongWin} &\approx \frac{1}{\text{CongWin}} \quad \text{per ciascun ACK} \\
    \implies \Delta \text{CongWin} &= +1 \quad \text{per ciascun RTD}
\end{align*}
Tale comportamento realizza la nota dinamica lineare ad accrescimento additivo (\emph{Additive Increase}).

\subsubsection{Gestione Delle Anomalie: Tahoe Vs Reno}

La reazione del sistema alla perdita di pacchetti dipende dalla variante implementativa del protocollo:

\begin{enumerate}
    \item \textbf{Scadenza del Timer (\emph{Timeout})}:
    Denota una congestione severa nella rete. Sia in \textbf{TCP Tahoe} sia in \textbf{TCP Reno}:
    \begin{align*}
        \text{threshold} &\gets \max\left(\frac{\text{CongWin}}{2}, \, 2\right) \\
        \text{CongWin} &\gets 1
    \end{align*}
    Il sistema ricomincia la trasmissione da capo operando in regime di \emph{Slow Start}.

    \item \textbf{Ricezione Di 3 ACK Duplicati (\emph{Triple Duplicate ACK})}:
    La ricezione di tre riscontri ripetuti per il medesimo byte indica che il canale continua a recapitare pacchetti a destinazione (generando riscontri duplicati per i segmenti fuori ordine), pur essendosi verificata la perdita isolata di un segmento intermedio:
    \begin{itemize}
        \item \textbf{TCP Tahoe}: ignora la distinzione tra eventi; imposta $\text{threshold} = \text{CongWin}/2$, fa crollare $\text{CongWin} = 1$ e riparte in Slow Start.
        \item \textbf{TCP Reno}: introduce le ottimizzazioni di \textbf{Fast Retransmit} e \textbf{Fast Recovery}. Il segmento mancante viene ritrasmesso immediatamente senza attendere la scadenza del timer di ritrasmissione e la finestra non viene azzerata, bensì dimezzata:
        \[
            \boxed{\text{threshold} \gets \frac{\text{CongWin}}{2}, \qquad \text{CongWin} \gets \text{threshold}}
        \]
        La trasmissione prosegue direttamente all'interno della fase di \emph{Congestion Avoidance}, garantendo un throughput medio marcatamente superiore ed evitando lo stallo della sessione.
    \end{itemize}
\end{enumerate}

% AA -----------------------------------------------------
% BB -----------------------------------------------------

\section{Fondamenti Di Sicurezza Nelle Reti E Nei Sistemi Distribuiti}

Nel contesto dei sistemi distribuiti moderni e delle infrastrutture di calcolo su scala geografica, la sicurezza delle comunicazioni di rete rappresenta un requisito architetturale imprescindibile. Nei sistemi tradizionali, le interazioni avvenivano prevalentemente all'interno di ambienti controllati, dove la topologia di rete fisica offriva una naturale barriera protettiva; tuttavia, con la decentralizzazione dei carichi di lavoro e l'interconnessione globale su canali intrinsecamente insicuri (come la rete Internet), l'isolamento e la protezione dei flussi informativi devono essere garantiti tramite rigorosi costrutti algoritmici e logici.

\subsection{I Quattro Pilastri Della Sicurezza Di Rete}

La sicurezza delle comunicazioni poggia tradizionalmente su quattro proprietà fondamentali, che indirizzano le diverse vulnerabilità a cui sono esposti i dati in transito:

\begin{theorem}[Proprietà Fondamentali Della Sicurezza Di Rete]
Una rete di calcolatori si definisce sicura rispetto a una specifica transazione se e solo se garantisce congiuntamente le seguenti proprietà:
\begin{enumerate}[leftmargin=*]
    \item \textbf{Riservatezza (\emph{Confidentiality})}: garantisce che solo il mittente legittimo e il destinatario designato siano in grado di comprendere e decodificare il contenuto del messaggio trasmesso. Dal punto di vista implementativo, il mittente trasforma il messaggio in chiaro (\emph{plaintext}) in un messaggio cifrato (\emph{ciphertext}) mediante un'operazione di cifratura; simmetricamente, il destinatario recupera il testo originario applicando un algoritmo di decifratura. Per un eventuale intercettore, il testo cifrato deve risultare computazionalmente indistinguibile da una sequenza casuale di bit.
    \item \textbf{Autenticazione (\emph{Authentication})}: assicura che entrambi gli estremi della sessione di comunicazione (mittente e ricevente) siano in grado di verificare e confermare in modo inequivocabile l'identità dichiarata dalla controparte, prevenendo attacchi di impersonificazione (\emph{spoofing}) o replicazione illegittima delle credenziali.
    \item \textbf{Integrità Del Messaggio (\emph{Message Integrity})}: garantisce che il contenuto del messaggio non sia stato alterato, corrotto, parzialmente cancellato o iniettato arbitrariamente durante il transito (o successivamente in fase di stoccaggio) senza che tale manipolazione venga rilevata istantaneamente dal destinatario.
    \item \textbf{Accesso E Disponibilità (\emph{Access and Availability})}: impone che le risorse computazionali, i servizi applicativi e i canali di comunicazione rimangano costantemente accessibili e operativi per gli utenti autorizzati, contrastando tentativi di saturazione o interruzione intenzionale delle risorse di sistema.
\end{enumerate}
\end{theorem}

\subsection{Il Modello Di Riferimento Formale: Alice, Bob E Trudy}

Nella letteratura crittografica e nella progettazione dei protocolli di sicurezza, le interazioni avversariali vengono formalizzate attraverso una serie di attori canonici che modellano lo scambio informativo lungo un canale di comunicazione insicuro.

% [IMMAGINE DEL FILE DI SLIDE CON NOME "1) Ripasso_2"] Rappresentazione dello scenario di sicurezza: Alice (mittente) e Bob (ricevente) comunicano su un canale insicuro intercettato da Trudy (intruso) — Slide 7 / Pagina 11

In questo modello analitico:
\begin{itemize}[noitemsep]
    \item \textbf{Alice}: rappresenta l'entità mittente sicura (\emph{secure sender}) che intende trasmettere dati applicativi o messaggi di controllo.
    \item \textbf{Bob}: rappresenta l'entità ricevente sicura (\emph{secure receiver}) a cui è indirizzata la comunicazione.
    \item \textbf{Trudy}: modella l'avversario o intruso (\emph{intruder}) posizionato sul canale di comunicazione condiviso.
\end{itemize}

L'avversario Trudy possiede facoltà operative che si dividono in due categorie principali:
\begin{itemize}[leftmargin=*]
    \item \textbf{Attacchi Passivi (\emph{Eavesdropping})}: Trudy intercetta e registra i pacchetti in transito senza alterarne il flusso. L'obiettivo primario è la violazione della riservatezza, estraendo informazioni sensibili (credenziali, dati industriali, chiavi).
    \item \textbf{Attacchi Attivi}: Trudy interagisce attivamente con il canale mediante:
    \begin{itemize}[noitemsep]
        \item \emph{Cancellazione}: soppressione arbitraria di pacchetti per causare disservizi o forzare ritrasmissioni;
        \item \emph{Modifica}: alterazione del carico utile (\emph{payload}) o delle intestazioni per compromettere l'integrità;
        \item \emph{Iniezione e Replay}: introduzione di pacchetti fraudolenti o reinvio di vecchi messaggi legittimi per violare l'autenticazione.
    \end{itemize}
\end{itemize}

\subsection{Linguaggio E Notazione Della Crittografia}

Per formalizzare matematicamente i meccanismi di protezione contro un avversario sul canale, si definiscono gli operatori crittografici fondamentali.

% [IMMAGINE DEL FILE DI SLIDE CON NOME "1) Ripasso_2"] Schema a blocchi della crittografia: Alice cifra il messaggio m con la chiave K_A ottenendo K_A(m), Bob decifra con la chiave K_B riottenendo m, con Trudy sul canale — Slide 8 / Pagina 12

Sia $m$ il messaggio originale in chiaro (\emph{plaintext}), appartenente allo spazio dei messaggi $\mathcal{M}$. Quando Alice intende spedire $m$ a Bob, non lo immette direttamente sul canale, ma lo sottopone a una funzione di cifratura parametrizzata da una specifica chiave crittografica $K_A$:
\[
c = K_A(m)
\]
dove $c$ rappresenta il testo cifrato (\emph{ciphertext}), appartenente allo spazio dei testi cifrati $\mathcal{C}$. Il messaggio $c$ viaggia sul canale pubblico esposto a Trudy.

All'estremo di ricezione, Bob riceve $c$ e applica un algoritmo di decifratura parametrizzato dalla propria chiave crittografica $K_B$:
\[
m = K_B(c) = K_B\bigl(K_A(m)\bigr)
\]
Affinché il sistema crittografico sia corretto, la composizione della trasformazione di cifratura e di decifratura deve equivalere all'operatore identità sullo spazio dei messaggi:
\begin{equation}
\boxed{K_B\bigl(K_A(m)\bigr) = m}
\end{equation}
A seconda della relazione algebrica esistente tra la chiave di cifratura $K_A$ e la chiave di decifratura $K_B$, la crittografia si divide tassonomicamente in due grandi paradigmi:
\begin{itemize}[leftmargin=*]
    \item \textbf{Sistemi a Chiave Simmetrica}: in cui $K_A = K_B = K_S$. Mittente e destinatario condividono un unico segreto a priori.
    \item \textbf{Sistemi a Chiave Asimmetrica (o Chiave Pubblica)}: in cui $K_A \neq K_B$. Una chiave è resa pubblicamente accessibile, mentre l'altra è custodita rigorosamente segreta dal legittimo proprietario.
\end{itemize}

\newpage

\section{Evoluzione Del Perimetro Di Sicurezza E Paradigma Zero-Trust}

\subsection{Il Collasso Del Perimetro Tradizionale}

Nelle architetture di rete aziendali classiche, la sicurezza era strutturata secondo un modello cosiddetto \emph{castle-and-moat} (castello e fossato): si assumeva l'esistenza di un perimetro fisico ben definito, coincidente con la rete locale aziendale (\emph{on-premises LAN}). All'interno del perimetro, le entità (dispositivi, utenti, server) godevano di un elevato livello di fiducia implicita, mentre il confine esterno era presidiato da apparati di filtraggio statici come firewall perimetrali e sistemi IDS/IPS.

Tale modello è divenuto obsoleto e inadeguato a causa di profonde trasformazioni tecnologiche e organizzative:
\begin{itemize}[leftmargin=*]
    \item \textbf{Evoluzione della collaborazione digitale}: la forza lavoro accede alle risorse aziendali da qualsiasi posizione geografica e da qualsiasi connettività Internet, con una massiccia accelerazione del lavoro da remoto (\emph{smart working}).
    \item \textbf{Adozione di servizi Software as a Service (SaaS)}: carichi di lavoro mission-critical risiedono al di fuori del datacenter privato, ospitati su infrastrutture cloud distribuite e multi-tenant.
    \item \textbf{Paradigma BYOD (\emph{Bring Your Own Device})}: dipendenti e consulenti utilizzano dispositivi personali per accedere a dati e servizi aziendali, introducendo endpoint eterogenei e privi di configurazioni di conformità omogenee.
    \item \textbf{Dispositivi non gestiti (\emph{Unmanaged Devices})}: partner industriali, fornitori esterni e clienti interagiscono regolarmente con le piattaforme aziendali tramite dispositivi su cui l'organizzazione ospitante non ha alcun controllo amministrativo.
    \item \textbf{Dispositivi Internet of Things (IoT)}: apparati a bassa potenza di calcolo, sensori industriali e attuatori intelligenti vengono distribuiti capillarmente sia nella rete interna sia presso i clienti, costituendo vettori di vulnerabilità hardware e firmware privi di supporto per agenti software avanzati.
\end{itemize}

Di conseguenza, il perimetro di rete basato su indirizzi IP interni e segmentazione fisica non esiste più: \textbf{l'identità è divenuta il nuovo perimetro di sicurezza primario}.

% [IMMAGINE DEL FILE DI SLIDE CON NOME "1) Ripasso_2"] Architettura con l'Identità al centro del perimetro di sicurezza che orchestra l'accesso tra Cloud apps, On-premises apps, Devices, Partners/Customer ed Employee — Slide 2 / Pagina 3

In questo nuovo contesto, per \textbf{identità} si intende il meccanismo formale attraverso il quale una qualsiasi entità (un utente umano, un'applicazione software, un microservizio o un dispositivo hardware) può essere univocamente verificata e autenticata per dimostrare di essere effettivamente chi dichiara di essere.

\subsection{La Metodologia Zero-Trust}

Per indirizzare la scomparsa del confine perimetrale, le moderne architetture adottano la metodologia **Zero-Trust**, il cui assioma fondamentale è: \emph{"Trust no one, verify everything"} (non fidarsi di nessuno, verificare ogni cosa). A differenza del modello tradizionale, nel paradigma Zero-Trust la collocazione topologica all'interno della rete privata non conferisce alcun privilegio di fiducia implicito.

% [IMMAGINE DEL FILE DI SLIDE CON NOME "1) Ripasso_2"] Schema architetturale a tempio della metodologia Zero-Trust: base fondata su "Verify explicitly", "Least privileged access", "Assume breach" e sei colonne portanti (Identities, Devices, Applications, Data, Infrastructure, Networks) — Slide 3 / Pagina 4

\subsubsection{I Tre Principi Guida}

L'approccio Zero-Trust si basa su tre postulati operativi:
\begin{enumerate}[leftmargin=*]
    \item \textbf{Verifica Esplicita (\emph{Verify Explicitly})}: ogni richiesta di accesso deve essere sempre autenticata e autorizzata esplicitamente basandosi su tutti i dati di contesto disponibili (identità dell'utente, posizione, stato di salute del dispositivo, classificazione del servizio e anomalie comportamentali).
    \item \textbf{Accesso Al Minimo Privilegio (\emph{Least Privileged Access})}: l'accesso degli utenti e dei servizi deve essere limitato rigorosamente al minimo insieme di permessi necessari per completare lo specifico compito operativo, ricorrendo a paradigmi dinamici quali il \emph{Just-In-Time} (JIT) e il \emph{Just-Enough-Access} (JEA).
    \item \textbf{Assunzione Della Violazione (\emph{Assume Breach})}: l'architettura deve essere progettata assumendo che un attaccante sia già penetrato all'interno del sistema o che un'identità sia già stata compromessa. Ciò impone la minimizzazione del raggio d'impatto (\emph{blast radius}) tramite microsegmentazione, crittografia end-to-end e monitoraggio continuo della telemetria per individuare tempestivamente comportamenti anomali.
\end{enumerate}

\subsubsection{I Sei Pilastri Fondamentali}

La realizzazione pratica del modello Zero-Trust poggia su sei pilastri architetturali:
\begin{itemize}[leftmargin=*]
    \item \textbf{Identità (\emph{Identities})}: rappresentano utenti umani, account di servizio o dispositivi. Quando un'identità tenta di accedere a una risorsa, deve superare un'autenticazione forte (es. autenticazione a più fattori, MFA) e vedersi applicate policy di autorizzazione contestuali e adattive.
    \item \textbf{Dispositivi (\emph{Devices})}: costituiscono un'ampia superficie d'attacco poiché i dati fluiscono continuamente dagli endpoint verso i carichi di lavoro on-premises e cloud. I dispositivi devono essere costantemente monitorati per verificarne lo stato di salute, l'integrità del sistema operativo e la conformità alle policy di sicurezza prima di autorizzare l'accesso.
    \item \textbf{Applicazioni (\emph{Applications})}: costituiscono il canale attraverso cui i dati vengono consumati. Questo pilastro richiede la visibilità completa su tutti gli applicativi in uso, contrastando il fenomeno dello \emph{Shadow IT} (l'utilizzo di software non autorizzato o non gestito centralmente dal reparto IT), oltre alla regolamentazione rigorosa di permessi, API e sessioni operative.
    \item \textbf{Dati (\emph{Data})}: rappresentano l'asset primario da proteggere. I dati devono essere classificati, etichettati e cifrati in base al loro livello di confidenzialità, garantendo che rimangano protetti anche qualora dovessero fuoriuscire dai confini dei dispositivi o delle reti gestite dall'organizzazione.
    \item \textbf{Infrastruttura (\emph{Infrastructure})}: include macchine virtuali, container, server fisici e componenti cloud, configurandosi come un vettore di minaccia critico. Il pilastro impone l'analisi continua delle versioni software, della configurazione di sistema, l'impiego di accessi temporanei JIT e l'elaborazione di telemetria avanzata per rilevare anomalie e innescare azioni di blocco automatico.
    \item \textbf{Reti (\emph{Networks})}: la rete non deve più essere considerata una zona di transito fidata. Deve essere strutturata mediante una profonda microsegmentazione interna, protezione dalle minacce in tempo reale, cifratura end-to-end di tutto il traffico, ispezione avanzata e analisi continua dei pattern di flusso.
\end{itemize}

\newpage

\section{Il Modello Di Responsabilità Condivisa Nel Cloud Computing}

Nelle architetture tradizionali interamente basate su datacenter privati (\emph{on-premises}), l'organizzazione cliente detiene la responsabilità esclusiva (al 100\%) dell'intera catena tecnologica, dalla sicurezza fisica degli stabili e dell'alimentazione elettrica, fino alla cifratura dei dati applicativi e alla configurazione dei firewall logici.

Con la transizione verso il cloud computing, tale onere viene suddiviso tra il fornitore del servizio cloud (\emph{Cloud Service Provider}, CSP) e l'organizzazione cliente secondo il **Modello Di Responsabilità Condivisa** (\emph{Shared Responsibility Model}). Questo modello definisce formalmente i confini operativi e legali della sicurezza, evitando ambiguità che potrebbero generare falle di protezione.

% [IMMAGINE DEL FILE DI SLIDE CON NOME "1) Ripasso_2"] Matrice del Modello di Responsabilità Condivisa tra fornitore (Microsoft/CSP) e cliente nei modelli SaaS, PaaS, IaaS e On-premises — Slide 4 / Pagina 5

\subsection{Ripartizione Delle Responsabilità Per Modello Di Servizio}

L'esatta demarcazione tra le competenze del fornitore e quelle del cliente varia in funzione del modello di distribuzione adottato:

\begin{table}[H]
\centering
\small
\begin{tabularx}{\textwidth}{lcccc}
\toprule
\textbf{Ambito di Responsabilità} & \textbf{SaaS} & \textbf{PaaS} & \textbf{IaaS} & \textbf{On-Premises} \\
\midrule
Informazioni e Dati & Cliente & Cliente & Cliente & Cliente \\
Dispositivi (Mobile e PC) & Cliente & Cliente & Cliente & Cliente \\
Account e Gestione Identità & Cliente & Cliente & Cliente & Cliente \\
Infrastruttura Identità e Directory & Condivisa & Condivisa & Condivisa & Cliente \\
Applicazioni & Provider & Condivisa / Cliente & Cliente & Cliente \\
Controlli di Rete & Provider & Condivisa & Cliente & Cliente \\
Sistema Operativo & Provider & Provider & Cliente & Cliente \\
Host Fisici di Calcolo & Provider & Provider & Provider & Cliente \\
Rete Fisica & Provider & Provider & Provider & Cliente \\
Datacenter Fisico & Provider & Provider & Provider & Cliente \\
\bottomrule
\end{tabularx}
\caption{Matrice di demarcazione della sicurezza nel Modello di Responsabilità Condivisa.}
\end{table}

\subsubsection{On-Premises Datacenter}
Nel modello locale, l'organizzazione ha la proprietà e il pieno controllo di tutti gli strati dell'infrastruttura. È pertanto interamente responsabile della sicurezza fisica (controllo accessi ai rack, videosorveglianza, prevenzione incendi), della ridondanza elettrica, dell'hardware di calcolo e storage, del cablaggio, degli hypervisor, dei sistemi operativi, delle patch software e delle policy sui dati.

\subsubsection{Infrastructure as a Service (IaaS)}
Nel modello IaaS, il fornitore cloud si assume l'onere della gestione dei componenti fisici del datacenter: server fisici, apparati di rete, alimentazione e virtualizzazione di base. Tra tutti i modelli cloud, l'IaaS è quello che richiede il massimo livello di gestione e competenza tecnica da parte del cliente. Il cliente è infatti responsabile della scelta, installazione, configurazione e aggiornamento dei sistemi operativi delle macchine virtuali, della gestione dei controlli di rete virtuali (routing, subnetting, Security Group, tabelle di instradamento), della sicurezza degli applicativi distribuiti e della protezione dei dati.

\subsubsection{Platform as a Service (PaaS)}
Il modello PaaS astrae l'infrastruttura sottostante per fornire agli sviluppatori un ambiente pronto per la compilazione, il testing e il dispiegamento rapido di software. Il cloud provider gestisce interamente l'hardware, i cluster fisici e i sistemi operativi sottostanti, facendosi carico dell'applicazione automatica delle patch di sicurezza a livello OS. La responsabilità del cliente si concentra sulla corretta configurazione e implementazione del codice applicativo e sulla gestione sicura dei database e dei dati residenti.

\subsubsection{Software as a Service (SaaS)}
Nel modello SaaS, l'intero stack applicativo e infrastrutturale è erogato e gestito direttamente dal fornitore cloud (come nei casi di Microsoft 365, Skype o Dynamics CRM Online), tipicamente accessibile mediante abbonamento periodico. L'utente o l'azienda cliente delega al fornitore la gestione della disponibilità, delle patch del sistema operativo, dell'infrastruttura di rete e del codice applicativo. Il carico di gestione per il cliente è minimo, ma restano sotto la sua esclusiva responsabilità la governance dei dati, la configurazione degli accessi e la sicurezza degli endpoint.

\subsection{Gli Elementi A Responsabilità Esclusiva Del Cliente}

Indipendentemente dal modello di cloud adottato (da IaaS a SaaS), esistono tre componenti la cui responsabilità di protezione rimane **sempre e in ogni circostanza a carico del cliente**:
\begin{enumerate}[noitemsep]
    \item \textbf{Informazioni e Dati}: classificazione, cifratura e controllo dei permessi di accesso ai contenuti archiviati o elaborati;
    \item \textbf{Dispositivi (Mobile, Laptop, Desktop)}: sicurezza e stato di conformità dei terminali utilizzati per accedere ai servizi;
    \item \textbf{Account e Identità}: gestione del ciclo di vita delle credenziali, autenticazione multi-fattore e segregazione dei ruoli.
\end{enumerate}

\newpage

\section{Tassonomia Delle Minacce Informatiche Comuni}

La comprensione delle modalità con cui gli attaccanti tentano di compromettere l'architettura IT è fondamentale per progettare opportune contromisure difensive.

\subsection{Violazione Dei Dati (Data Breach)}

Una \textbf{violazione dei dati} (\emph{data breach}) si verifica quando dati confidenziali, riservati o protetti vengono esfiltrati, copiati o sottratti da un ambiente aziendale senza autorizzazione. Questa minaccia riguarda principalmente i **dati personali**, definiti come qualsiasi informazione relativa a una persona fisica identificata o identificabile, direttamente o indirettamente (dati anagrafici, numeri di documenti, credenziali bancarie, record sanitari).

Le tecniche di attacco comunemente sfruttate per provocare un data breach includono:
\begin{itemize}[leftmargin=*]
    \item \textbf{Phishing}: campagne massive basate su comunicazioni ingannevoli (email, messaggi) che simulano comunicazioni legittime per indurre gli utenti a rivelare credenziali o dati sensibili.
    \item \textbf{Spear Phishing}: variante mirata e personalizzata del phishing, diretta verso individui specifici all'interno di un'organizzazione (spesso dotati di privilegi elevati), basata su informazioni preliminari raccolte sulla vittima.
    \item \textbf{Truffe di Supporto Tecnico (\emph{Tech Support Scams})}: tattiche di ingegneria sociale in cui l'attaccante si spaccia per un operatore di assistenza tecnica legittimo per ottenere accesso remoto non autorizzato al dispositivo della vittima.
    \item \textbf{SQL Injection (SQLi)}: tecnica di sfruttamento delle vulnerabilità a livello applicativo in cui stringhe di codice SQL malevolo vengono inserite nei campi di input dell'utente, permettendo all'attaccante di eseguire query arbitrarie sul database di backend, eludendo l'autenticazione ed esfiltrando interi archivi.
    \item \textbf{Malware di Furto Credenziali (\emph{Info-stealer})}: agenti software nocivi espressamente concepiti per sottrarre password salvate nei browser, cookie di sessione o coordinate bancarie.
\end{itemize}

\subsection{Attacchi All'Identità: Attacco A Dizionario E Forza Bruta}

L'**attacco a dizionario** (\emph{dictionary attack}) è una forma specializzata di attacco a forza bruta (\emph{brute-force}) contro i meccanismi di autenticazione. L'attaccante tenta di violare un account provando in modo automatizzato una vasta sequenza di parole, frasi comuni o credenziali precedentemente compromesse e catalogate all'interno di elenchi predefiniti (\emph{wordlist} o dizionari). Ciascuna stringa candidata viene testata a fronte di un nome utente noto fino all'individuazione della corrispondenza corretta.

\subsection{Ransomware E Software Malevolo}

Il termine generico **malware** identifica qualsiasi codice o applicazione progettata per recare danno, sottrarre dati o interrompere il corretto funzionamento di un dispositivo di calcolo. La principale motivazione dei criminali informatici è il tornaconto finanziario, perseguito tramite l'uso di botnet per sferrare attacchi coordinati, il furto di identità bancarie, la vendita di accessi o l'estorsione diretta.

Il **Ransomware** è una tipologia particolarmente pervasiva di malware che, una volta penetrata nel sistema target:
\begin{enumerate}[noitemsep]
    \item Esegue la scansione dei dischi locali e delle condivisioni di rete;
    \item Cifra in modo irreversibile file, cartelle e database della vittima mediante algoritmi crittografici robusti;
    \item Rende inaccessibile l'ambiente operativo e presenta una richiesta di riscatto (\emph{ransom}), generalmente richiesta in criptovalute in cambio della chiave crittografica di decifratura.
\end{enumerate}

\subsection{Altre Categorie Di Minacce Informatiche}

\begin{itemize}[leftmargin=*]
    \item \textbf{Coin Miner (Cryptojacker)}: malware silenti che abusano delle risorse computazionali (CPU e GPU) del sistema ospite all'insaputa dell'utente per effettuare il calcolo distribuito di criptovalute.
    \item \textbf{Rootkit}: pacchetti di software malevolo a basso livello progettati per intercettare e manomettere i processi standard del sistema operativo (spesso a livello kernel o bootloader). Una volta installato un rootkit, l'intero strato di telemetria e diagnostica del sistema cessa di essere attendibile, poiché il malware maschera la propria presenza e altera i report di stato del sistema stesso.
    \item \textbf{Trojan}: applicazioni malevole mascherate da software lecito e utile per indurre l'utente a scaricarle ed eseguirle. I trojan non possiedono la capacità di propagarsi autonomamente da un host all'altro: richiedono l'interazione umana o l'intervento di un altro componente malevolo per l'installazione iniziale.
    \item \textbf{Worm}: programmi malevoli capaci di auto-replicarsi e propagarsi autonomamente lungo la rete senza alcuna necessità di interazione umana o di legarsi a file eseguibili ospiti. Sfruttano attivamente vulnerabilità di rete, protocolli di condivisione file, allegati email e dischi rimovibili.
    \item \textbf{Vulnerabilità Ed Exploit (ed Exploit Kit)}: una \textbf{vulnerabilità} è un difetto logico o una debolezza insita nel codice di un software. Un \textbf{exploit} è una sequenza di comandi o un frammento di codice che sfrutta tale debolezza per aggirare i controlli di sicurezza, permettendo l'esecuzione di codice arbitrario da parte di pacchetti automatizzati noti come \textbf{exploit kit}.
\end{itemize}

\subsection{Attacchi Di Interruzione Del Servizio (DoS E DDoS)}

Un attacco **Denial of Service (DoS)** mira a compromettere la disponibilità di una risorsa informatica esaurendone intenzionalmente le capacità computazionali, di memoria o di banda passante.

Nella variante distribuita, nota come **Distributed Denial of Service (DDoS)**, l'attaccante coordina simultaneamente decine o centinaia di migliaia di macchine precedentemente compromesse (organizzate in una \emph{botnet}) per inondare di traffico un target specifico. Qualsiasi interfaccia o servizio applicativo pubblicamente raggiungibile attraverso la rete Internet può costituire il bersaglio di un attacco DDoS. L'effetto finale è la completa saturazione dei buffer di rete, della capacità di elaborazione delle connessioni o dei socket di trasporto, rendendo l'applicazione irraggiungibile per gli utenti legittimi.

\newpage

\section{Meccanismi Di Protezione: Crittografia E Funzioni Hash}

La protezione dell'integrità e della confidenzialità delle informazioni poggia su due tecniche distinte ma complementari: la crittografia (simmetrica e asimmetrica) e l'hashing.

% [IMMAGINE DEL FILE DI SLIDE CON NOME "1) Ripasso_2"] Rappresentazione visiva di crittografia simmetrica (chiave condivisa), crittografia asimmetrica (chiave pubblica e privata) e funzione di hashing unidirezionale — Slide 6 / Pagina 9

\subsection{Dati A Riposo (\emph{At Rest}) E Dati In Transito (\emph{In Transit})}

Le strategie di sicurezza distinguono lo stato dei dati in base alla loro condizione operativa:

\subsubsection{Cifratura Dei Dati A Riposo (\emph{Encryption at Rest})}
I dati a riposo sono memorizzati in modo persistente su dispositivi fisici, quali hard disk, unità a stato solido, database relazionali o bucket di storage nel cloud. La cifratura a riposo garantisce che il dato sia convertito in una forma non intellegibile in assenza delle opportune chiavi crittografiche. Se un attaccante sottrae fisicamente un'unità di memorizzazione o ottiene l'accesso non autorizzato a snapshot di dischi virtuali, la mancanza della chiave rende impossibile la decifratura dei contenuti.

\subsubsection{Cifratura Dei Dati In Transito (\emph{Encryption in Transit})}
I dati in transito sono informazioni in movimento lungo una rete fisica o logica (attraverso la rete pubblica Internet o all'interno di reti private tra microservizi). La protezione in transito è tipicamente implementata a diversi livelli dello stack protocollare:
\begin{itemize}[noitemsep]
    \item A livello di trasporto e applicativo: tramite protocolli sicuri quali **HTTPS** (basato su **TLS**, \emph{Transport Layer Security});
    \item A livello di rete: tramite tunnel **IPsec** o protocolli overlay cifrati.
\end{itemize}
Tale meccanismo impedisce ad osservatori esterni sul canale di ispezionare il payload o manipolare le sessioni attive.

\subsection{Funzioni Hash Crittografiche}

A differenza degli algoritmi di cifratura, una funzione di hash opera in modo completamente unidirezionale:
\begin{equation}
h = H(m)
\end{equation}
dove $H$ accetta in ingresso una stringa $m$ di lunghezza arbitraria e produce in uscita un digest $h$ avente una lunghezza fissa determinata (es. 256 bit per SHA-256).

Le caratteristiche formali di una funzione di hash sicura includono:
\begin{itemize}[leftmargin=*]
    \item \textbf{Determinismo}: a parità di messaggio $m$ in ingresso e a parità di algoritmo, viene generato invariabilmente il medesimo valore di hash $h$.
    \item \textbf{Non Invertibilità (Unidirezionalità)}: dato $h$, deve essere computazionalmente impossibile ricavare il messaggio originario $m$.
    \item \textbf{Assenza di Chiavi}: l'algoritmo non impiega chiavi crittografiche per il calcolo; il valore $h$ non viene mai ``decifrato''.
    \item \textbf{Resistenza alle Collisioni}: deve risultare computazionalmente proibitivo individuare due messaggi distinti $m_1 \neq m_2$ tali per cui $H(m_1) = H(m_2)$.
\end{itemize}

\subsection{Memorizzazione Delle Password E Meccanismo Del Salting}

Le funzioni di hash trovano applicazione fondamentale nella gestione e persistenza delle credenziali di accesso. Un sistema non deve mai memorizzare le password degli utenti in chiaro: al momento della creazione o modifica, il sistema calcola l'hash della password e memorizza esclusivamente tale digest. Durante l'autenticazione, la password inserita dall'utente viene nuovamente processata dalla medesima funzione di hash; se il digest ottenuto coincide con quello registrato nel database, l'identità è confermata.

\subsubsection{La Vulnerabilità Del Determinismo E Gli Attacchi A Dizionario}
Poiché le funzioni di hash sono deterministiche e gli algoritmi sono noti a livello pubblico, un attaccante che acquisisca un dump del database delle password può calcolare preventivamente l'hash di miliardi di parole di dizionario (eventualmente precompilando tali corrispondenze in strutture note come \emph{Rainbow Tables}). Non appena l'hash di una parola comune corrisponde a un valore memorizzato, l'attaccante scopre immediatamente la password della vittima.

\subsubsection{Mitigazione Tramite Salting (\emph{Salt})}
Per neutralizzare gli attacchi a dizionario e le tabelle precalcolate, le architetture moderne introducono una stringa casuale a lunghezza fissa denominata **salt** prima del calcolo dell'hash.

Dato un salt casuale $s$ generato univocamente per ciascun utente, la memorizzazione avviene calcolando:
\begin{equation}
\boxed{h = H(m \parallel s)}
\end{equation}
dove $\parallel$ indica l'operazione di concatenazione tra la password $m$ e il salt $s$.

Nel database vengono memorizzati sia il digest $h$ sia il valore in chiaro del salt $s$. I vantaggi ingegneristici di tale approccio sono:
\begin{itemize}[noitemsep]
    \item Anche se due utenti scelgono la medesima password identica $m$, avendo due valori di salt distinti ($s_1 \neq s_2$) i rispettivi hash risulteranno totalmente diversi: $H(m \parallel s_1) \neq H(m \parallel s_2)$.
    \item L'attaccante non può più sfruttare dizionari generici precalcolati, ma è costretto a ricomputare l'intero dizionario da zero specificamente per il salt di ogni singolo utente, rendendo l'attacco di massa computazionalmente impraticabile.
\end{itemize}

\newpage

\section{Crittografia Simmetrica}

Nel modello a chiave simmetrica, Alice e Bob condividono la medesima chiave segreta $K_S$. Tale chiave regola sia la trasformazione del testo in chiaro in testo cifrato, sia l'operazione inversa:

% [IMMAGINE DEL FILE DI SLIDE CON NOME "1) Ripasso_2"] Diagramma del protocollo a chiave simmetrica: Alice cifra il messaggio m con la chiave condivisa K_S e Bob lo decifra con la medesima K_S ottenendo m = K_S(K_S(m)) — Slide 9 / Pagina 13

\subsection{Cifrari A Sostituzione Monoalfabetica}

Il prototipo concettuale della crittografia simmetrica è il cifrario a sostituzione monoalfabetica, in cui ciascun carattere del testo in chiaro viene sostituito da un corrispondente carattere dell'alfabeto cifrato secondo una permutazione prestabilita.

% [IMMAGINE DEL FILE DI SLIDE CON NOME "1) Ripasso_2"] Esempio di cifrario monoalfabetico con mappatura delle 26 lettere e cifratura della frase di prova — Slide 10 / Pagina 14

Considerando l'alfabeto inglese di 26 caratteri, la chiave di cifratura consiste in una mappatura biiettiva dall'insieme dei 26 caratteri su se stesso. Ad esempio, data la permutazione:
\begin{center}
\small
\texttt{plaintext:  a b c d e f g h i j k l m n o p q r s t u v w x y z}\\
\texttt{ciphertext: m n b v c x z a s d f g h j k l p o i u y t r e w q}
\end{center}
Il messaggio in chiaro:
\[
\texttt{bob. i love you. alice}
\]
viene convertito nel messaggio cifrato:
\[
\texttt{nkn. s gktc wky. mgsbc}
\]
Sebbene lo spazio delle chiavi sia pari a $26! \approx 4 \times 10^{26}$ (rendendo insostenibile un attacco a forza bruta cieco), questo schema è debolissimo a fronte di un'analisi delle frequenze statistiche dei caratteri nella lingua di destinazione.

\subsection{Il Problema Fondamentale Dello Scambio Delle Chiavi}

Il limite architetturale principale della crittografia simmetrica risiede nel **problema della distribuzione delle chiavi**: come possono due entità remote (Alice e Bob) accordarsi in modo sicuro sul valore numerico del segreto $K_S$, in particolare se non hanno mai avuto un incontro fisico preliminare o se operano attraverso un canale pubblico completamente esposto all'intercettazione di Trudy?

\subsection{Data Encryption Standard (DES)}

Introdotto come standard federale statunitense dal NIST nel 1977 e riconfermato nel 1993, il **Data Encryption Standard (DES)** è stato per decenni il riferimento per la cifratura simmetrica a blocchi.

Le specifiche operative di DES comprendono:
\begin{itemize}[noitemsep]
    \item Struttura: cifrario a blocchi (\emph{block cipher}) con elaborazione a blocchi di dimensione fissa pari a 64 bit di testo in chiaro;
    \item Lunghezza della chiave: chiave nominale a 64 bit, di cui però solo 56 bit effettivi sono usati dall'algoritmo (gli 8 bit restanti fungono da controllo di parità);
    \item Modalità operativa: impiego standard con \emph{Cipher Block Chaining} (CBC).
\end{itemize}

\subsubsection{Architettura E Funzionamento Interno Del DES}
DES adotta una classica struttura a rete di Feistel a 16 iterazioni (\emph{rounds}).

% [IMMAGINE DEL FILE DI SLIDE CON NOME "1) Ripasso_2"] Architettura di Feistel del DES: blocco a 64 bit, permutazione iniziale, 16 round identici con sottochiavi a 48 bit, permutazione finale e blocco a 64 bit in uscita — Slide 12 / Pagina 16

La dinamica operativa si articola nei seguenti passaggi:
\begin{enumerate}[leftmargin=*]
    \item \textbf{Permutazione Iniziale (IP)}: il blocco di 64 bit di ingresso subisce un riordinamento fisso dei bit;
    \item \textbf{Scomposizione in semiblocchi}: il blocco viene diviso in due metà da 32 bit, denotate con $L_0$ e $R_0$;
    \item \textbf{16 Round di Feistel}: per ogni round $i \in \{1, \dots, 16\}$, la trasformazione calcola:
    \begin{align*}
        L_i &= R_{i-1} \\
        R_i &= L_{i-1} \oplus f(R_{i-1}, K_i)
    \end{align*}
    dove $\oplus$ indica l'operatore XOR bit a bit, $f$ è una funzione non lineare basata su box di espansione, box di sostituzione (S-Box) e permutazioni fisse, e $K_i$ è una sottochiave di 48 bit derivata dai 56 bit della chiave principale tramite un algoritmo di \emph{key schedule};
    \item \textbf{Permutazione Finale (FP)}: al termine del sedicesimo round, le due metà vengono invertite e sottoposte alla permutazione finale $IP^{-1}$, restituendo il blocco cifrato di 64 bit.
\end{enumerate}

\subsubsection{Sicurezza E Superamento Del DES}
Nonostante non siano mai stati scoperti attacchi analitici efficienti capaci di scardinare la struttura del cifrario, la dimensione ridotta dello spazio delle chiavi ($2^{56} \approx 7.2 \times 10^{16}$ combinazioni) ha condannato il DES all'obsolescenza. Con l'evoluzione della potenza di calcolo, la sfida \emph{DES Challenge} ha dimostrato la capacità di forzare una chiave DES mediante brute force in meno di un giorno.

Per estendere la vita utile dei sistemi basati su DES senza riprogettare l'hardware, è stato introdotto il **Triple-DES (3DES)**, che applica l'algoritmo tre volte consecutive impiegando tre chiavi distinte da 56 bit:
\[
c = E_{K_3}\bigl(D_{K_2}(E_{K_1}(m))\bigr)
\]
pur aumentando la sicurezza computazionale a 112 o 168 bit effettivi, il 3DES sconta un'intrinseca lentezza computazionale dovuta alla triplice esecuzione.

\subsection{Advanced Encryption Standard (AES)}

A fronte dei limiti del DES, nel novembre 2001 il NIST ha standardizzato l'**Advanced Encryption Standard (AES)**, basato sull'algoritmo Rijndael.

Le proprietà fondamentali di AES sono:
\begin{itemize}[leftmargin=*]
    \item Cifrario a blocchi con dimensione fissa del blocco di testo pari a \textbf{128 bit};
    \item Supporto per tre lunghezze di chiave: \textbf{128 bit}, \textbf{192 bit} e \textbf{256 bit};
    \item Resistenza computazionale: mentre la ricerca esaustiva di una chiave DES a 56 bit può essere completata in un secondo su cluster dedicati moderni, la violazione a forza bruta di una chiave AES a 128 bit richiederebbe un tempo stimato nell'ordine di circa \textbf{149 trilioni di anni} ($1.49 \times 10^{14}$ anni).
\end{itemize}

\newpage

\section{Crittografia Asimmetrica E L'Algoritmo RSA}

\subsection{Il Cambio Di Paradigma}

La crittografia asimmetrica o a chiave pubblica ha rivoluzionato il campo della sicurezza informatica a seguito delle pubblicazioni storiche di Whitfield Diffie e Martin Hellman (1976) e di Ron Rivest, Adi Shamir e Leonard Adleman (1978).

% [IMMAGINE DEL FILE DI SLIDE CON NOME "1) Ripasso_2"] Confronto concettuale tra modello simmetrico e modello asimmetrico a chiave pubblica — Slide 14 / Pagina 18

Nel modello asimmetrico:
\begin{itemize}[noitemsep]
    \item Mittente e destinatario non condividono alcun segreto a priori;
    \item Ciascuna entità (es. Bob) genera una coppia di chiavi complementari:
    \begin{itemize}
        \item Una **chiave pubblica** $K_B^+$, resa nota a chiunque sulla rete;
        \item Una **chiave privata** $K_B^-$, custodita rigorosamente segreta dal possessore.
    \end{itemize}
\end{itemize}

% [IMMAGINE DEL FILE DI SLIDE CON NOME "1) Ripasso_2"] Diagramma del protocollo a chiave pubblica: cifratura del messaggio m con la chiave pubblica di Bob K_B^+ e decifratura con la chiave privata K_B^- — Slide 15 / Pagina 19

Chiunque intenda inviare un messaggio riservato a Bob cifra il testo usando la chiave pubblica $K_B^+$. Solo ed esclusivamente Bob, in possesso della corrispondente chiave privata $K_B^-$, potrà decifrare il testo cifrato:
\begin{equation}
m = K_B^-\bigl(K_B^+(m)\bigr)
\end{equation}

\subsection{Requisiti Matematici Di Un Sistema A Chiave Pubblica}

Affinché un sistema di crittografia asimmetrica sia valido e sicuro, devono sussistere contemporaneamente due requisiti:
\begin{enumerate}[leftmargin=*]
    \item \textbf{Correttezza e Invertibilità}: devono esistere algoritmi efficienti per calcolare la coppia di funzioni $K_B^+(\cdot)$ e $K_B^-(\cdot)$ tali per cui l'inversione sia garantita su tutto lo spazio dei messaggi:
    \[
    K_B^-\bigl(K_B^+(m)\bigr) = m
    \]
    \item \textbf{Sicurezza Unidirezionale (\emph{Trapdoor})}: conoscendo pubblicamente la chiave di cifratura $K_B^+$ e l'algoritmo di cifratura, deve essere **computazionalmente impossibile ricavare la chiave privata** $K_B^-$.
\end{enumerate}

\subsection{Fondamenti Operativi Dell'Algoritmo RSA}

L'algoritmo **RSA** opera interpretando qualsiasi messaggio non come una mera sequenza di caratteri, ma come un pattern di bit univocamente corrispondente a un numero intero positivo.

Ad esempio, consideriamo il pattern di 8 bit:
\[
m = \texttt{10010001}_2
\]
Questo messaggio corrisponde al numero intero in base decimale:
\[
m = 1 \cdot 2^7 + 0 \cdot 2^6 + 0 \cdot 2^5 + 1 \cdot 2^4 + 0 \cdot 2^3 + 0 \cdot 2^2 + 0 \cdot 2^1 + 1 \cdot 2^0 = 128 + 16 + 1 = 145_{10}
\]
Cifrare il messaggio equivale quindi a eseguire un'operazione matematica sul numero associato per ottenere un nuovo valore numerico intero, che costituirà il testo cifrato.

\subsection{Generazione Della Coppia Di Chiavi RSA}

La procedura di inizializzazione per generare la chiave pubblica e la chiave privata in RSA si compone di 5 passaggi formali:

\begin{enumerate}[leftmargin=*]
    \item \textbf{Scelta dei numeri primi}: si scelgono due numeri primi molto grandi e distinti, $p$ e $q$ (nelle applicazioni moderne dell'ordine di 1024 o 2048 bit ciascuno).

    \item \textbf{Calcolo del modulo e del toziente}:
    si calcola il modulo pubblico $n$:
    \[
    n = p \cdot q
    \]
    e si determina il valore della funzione toziente di Eulero $\phi(n)$, denotata tradizionalmente con $z$:
    \[
    z = (p - 1)(q - 1)
    \]

    \item \textbf{Selezione dell'esponente pubblico $e$}:
    si sceglie un intero $e$ strettamente minore del modulo ($e < n$) tale che $e$ e $z$ siano coprimi (ovvero non abbiano fattori comuni oltre a 1):
    \[
    \gcd(e, z) = 1
    \]

    \item \textbf{Calcolo dell'esponente privato $d$}:
    si calcola l'inverso moltiplicativo modulare di $e$ rispetto al modulo $z$, determinando un intero $d$ tale che la quantità $(e \cdot d - 1)$ sia esattamente divisibile per $z$:
    \begin{equation}
    e \cdot d \equiv 1 \pmod z \quad \iff \quad e \cdot d \bmod z = 1
    \end{equation}

    \item \textbf{Definizione delle chiavi}:
    \begin{itemize}[noitemsep]
        \item La \textbf{chiave pubblica} $K_B^+$ è costituita dalla coppia di interi:
        \[
        \boxed{K_B^+ = (n, e)}
        \]
        \item La \textbf{chiave privata} $K_B^-$ è costituita dalla coppia di interi:
        \[
        \boxed{K_B^- = (n, d)}
        \]
    \end{itemize}
\end{enumerate}

\subsection{Operazioni Di Cifratura E Decifratura}

Disponendo della chiave pubblica $(n, e)$ e della chiave privata $(n, d)$, le operazioni di cifratura e decifratura sono definite nell'aritmetica modulare come segue:

\subsubsection{Cifratura}
Dato un messaggio $m$ codificato come intero tale che $0 \le m < n$, il testo cifrato $c$ si calcola elevando il messaggio all'esponente pubblico $e$ modulo $n$:
\begin{equation}
\boxed{c = m^e \bmod n}
\end{equation}

\subsubsection{Decifratura}
Ricevuto il testo cifrato $c$, il destinatario recupera il messaggio originario elevando $c$ al proprio esponente segreto $d$ modulo $n$:
\begin{equation}
\boxed{m = c^d \bmod n}
\end{equation}

\subsubsection{Validità Algebrica}
La correttezza matematica del processo poggia sulla proprietà per cui:
\[
c^d \bmod n = \bigl(m^e \bmod n\bigr)^d \bmod n = m^{ed} \bmod n
\]
Poiché per costruzione $ed \equiv 1 \pmod{\phi(n)}$, per il Teorema di Eulero (generalizzazione del Piccolo Teorema di Fermat), risulta che per ogni intero $m$ coprimo con $n$:
\[
m^{ed} \equiv m^{k \cdot \phi(n) + 1} \equiv \bigl(m^{\phi(n)}\bigr)^k \cdot m \equiv (1)^k \cdot m \equiv m \pmod n
\]
garantendo così il perfetto ripristino del testo originale.

\subsection{Esempio Numerico Dettagliato Passo-Passo}

Di seguito viene illustrato il calcolo analitico completo dell'algoritmo RSA utilizzando i valori didattici proposti nelle slide:

% [IMMAGINE DEL FILE DI SLIDE CON NOME "1) Ripasso_2"] Esempio numerico RSA con p=5, q=7, e=5, d=29, cifratura del messaggio m=12 (00001100) in c=17 e successiva decifratura — Slide 20 / Pagina 24

\subsubsection{Fase 1: Generazione Delle Chiavi}
\begin{enumerate}[leftmargin=*]
    \item Scelta dei numeri primi:
    \[
    p = 5, \quad q = 7
    \]
    \item Calcolo del modulo $n$ e del parametro $z$:
    \[
    n = p \cdot q = 5 \cdot 7 = 35
    \]
    \[
    z = (p - 1)(q - 1) = (5 - 1)(7 - 1) = 4 \cdot 6 = 24
    \]
    \item Scelta dell'esponente pubblico $e$:
    Scegliamo $e = 5$. Verifichiamo che sia coprimo con $z = 24$:
    \[
    \gcd(5, 24) = 1 \quad \implies \quad \text{scelta valida}
    \]
    \item Calcolo dell'esponente privato $d$:
    Dobbiamo determinare un valore intero $d$ che soddisfi la congruenza:
    \[
    5 \cdot d \equiv 1 \pmod{24} \quad \iff \quad 5d - 1 = k \cdot 24 \quad (k \in \mathbb{Z}^+)
    \]
    Scegliendo $d = 29$:
    \[
    e \cdot d = 5 \cdot 29 = 145
    \]
    Dividendo per $z = 24$:
    \[
    145 = 6 \cdot 24 + 1 \quad \implies \quad 145 \bmod 24 = 1
    \]
    Pertanto $d = 29$ è un valore valido.
    \item Risultato della generazione delle chiavi:
    \[
    K_B^+ = (n=35, \, e=5), \qquad K_B^- = (n=35, \, d=29)
    \]
\end{enumerate}

\subsubsection{Fase 2: Cifratura Del Messaggio}
Si consideri il messaggio espresso come sequenza di 8 bit:
\[
\texttt{00001100}_2
\]
Convertito in rappresentazione decimale:
\[
m = 0 \cdot 2^7 + \dots + 1 \cdot 2^3 + 1 \cdot 2^2 + 0 \cdot 2^1 + 0 \cdot 2^0 = 8 + 4 = 12
\]
Si noti che $m = 12 < n = 35$, rispettando il vincolo imposto dall'algoritmo.

L'operazione di cifratura consiste nel calcolo di:
\[
c = m^e \bmod n = 12^5 \bmod 35
\]
Calcoliamo la potenza $12^5$:
\begin{align*}
    12^1 &= 12 \\
    12^2 &= 144 \\
    12^3 &= 144 \cdot 12 = 1728 \\
    12^4 &= 1728 \cdot 12 = 20736 \\
    12^5 &= 20736 \cdot 12 = 248832
\end{align*}
Eseguiamo ora la divisione modulare $248832 \bmod 35$:
\[
\frac{248832}{35} = 7109.4857\dots \implies 7109 \cdot 35 = 248815
\]
\[
c = 248832 - 248815 = 17
\]
Il testo cifrato trasmesso lungo il canale insicuro è pertanto:
\[
\boxed{c = 17}
\]

\subsubsection{Fase 3: Decifratura Del Testo Cifrato}
Bob riceve $c = 17$ e applica la propria chiave privata $(n=35, d=29)$:
\[
m = c^d \bmod n = 17^{29} \bmod 35
\]
Il valore numerico esplicito di $17^{29}$ è un numero di 36 cifre:
\[
17^{29} = 481968572106750915091411825223071697
\]
Eseguendo la riduzione modulare:
\[
\frac{481968572106750915091411825223071697}{35} = 13770530631621454716897480720659191 \text{ con resto } 12
\]
Si ottiene esattamente:
\[
\boxed{m = 12}
\]
riconvertibile nella sequenza binaria originaria $\texttt{00001100}_2$, a conferma dell'esattezza e dell'invertibilità del procedimento.

% AA -----------------------------------------------------------------
% BB -----------------------------------------------------------------
\section{Proprietà Di Commutatività Operativa Dell'Algoritmo RSA}

L'architettura matematica dell'algoritmo RSA presenta una proprietà algebrica fondamentale che ne estende le funzionalità ben oltre la semplice confidenzialità dei dati: l'invertibilità dell'ordine di applicazione della chiave pubblica e della chiave privata.

Nello schema di cifratura convenzionale, orientato a garantire la riservatezza, il mittente impiega la chiave pubblica del destinatario $K_B^+$ per cifrare il messaggio $m$, e il destinatario decifra il testo risultante mediante la propria chiave privata $K_B^-$:
\[
K_B^-\bigl(K_B^+(m)\bigr) = m
\]
Tuttavia, invertendo rigorosamente l'ordine temporale e funzionale degli operatori, si ottiene un'identità algebrica analoga:
\begin{equation}
\boxed{K_B^-\bigl(K_B^+(m)\bigr) = m = K_B^+\bigl(K_B^-(m)\bigr)}
\end{equation}

% [IMMAGINE DEL FILE DI SLIDE CON NOME "1) Ripasso_2"] Proprietà fondamentale di RSA: uguaglianza del risultato applicando prima la chiave pubblica e poi la privata, oppure prima la chiave privata e poi la pubblica — Slide 21 / Pagina 1

\subsection{Fondamento Algebrico Della Commutatività}

La validità di tale proprietà discende in modo diretto dall'aritmetica modulare su cui poggia RSA. Date le due chiavi appartenenti alla medesima coppia:
\[
K_B^+ = (n, e), \qquad K_B^- = (n, d)
\]
le due operazioni di trasformazione si traducono in elevamenti a potenza modulari:
\begin{align*}
    K_B^+(x) &\equiv x^e \pmod n \\
    K_B^-(x) &\equiv x^d \pmod n
\end{align*}
Applicando prima la chiave privata sul messaggio $m$ e successivamente la chiave pubblica sul risultato parziale, si ottiene la catena di congruenze:
\begin{equation}
K_B^+\bigl(K_B^-(m)\bigr) \equiv \bigl(m^d \bmod n\bigr)^e \bmod n \equiv \bigl(m^d\bigr)^e \equiv m^{d \cdot e} \pmod n
\end{equation}
Poiché il prodotto tra numeri interi gode della proprietà commutativa nello scalare dell'esponente ($d \cdot e = e \cdot d$), e ricordando che per costruzione della coppia di chiavi sussiste la congruenza $e \cdot d \equiv 1 \pmod{\phi(n)}$, applicando il Teorema di Eulero si ricava:
\begin{equation}
m^{d \cdot e} \equiv m^{e \cdot d} \equiv m^{k \cdot \phi(n) + 1} \equiv \bigl(m^{\phi(n)}\bigr)^k \cdot m \equiv (1)^k \cdot m \equiv m \pmod n
\end{equation}

\subsection{Implicazione Architetturale: Dalla Cifratura Alla Firma}

Questa simmetria operativa ha un impatto concettuale rilevantissimo:
\begin{itemize}[leftmargin=*]
    \item \textbf{Cifratura per Riservatezza}: chiunque può applicare la chiave pubblica $K_B^+$ (essendo di pubblico dominio) per cifrare un dato, ma un'unica entità (il detentore di $K_B^-$) è matematicamente in grado di decifrarlo.
    \item \textbf{Trasformazione per Autenticazione e Firma}: solo il detentore della chiave privata $K_B^-$ è in grado di calcolare la quantità $K_B^-(m)$. Chiunque nella rete, disponendo della chiave pubblica $K_B^+$, può verificare l'operazione calcolando $K_B^+\bigl(K_B^-(m)\bigr) = m$. Questa seconda modalità operativa costituisce la base algoritmica della **firma digitale** e dell'**autenticazione dell'entità**.
\end{itemize}

\newpage

\section{Progettazione Di Protocolli Di Autenticazione: Evoluzione Ed Attacchi}

L'obiettivo formale di un protocollo di autenticazione è consentire a un'entità verificatrice (\textbf{Bob}) di accertare senza ambiguità l'identità dichiarata da un'entità remota richiedente (\textbf{Alice}) attraverso un canale di comunicazione non fidato, presidiato da un avversario attivo (\textbf{Trudy}).

Per comprendere la robustezza dei moderni protocolli crittografici, è istruttivo analizzare l'evoluzione incrementale delle difese a fronte dei fallimenti riscontrati nei protocolli ingenui.

\subsection{Protocollo ap1.0: Dichiarazione Non Protetta}

Nel protocollo più elementare immaginabile, Alice dichiara semplicemente la propria identità trasmettendo un messaggio in chiaro lungo il canale.

% [IMMAGINE DEL FILE DI SLIDE CON NOME "1) Ripasso_2"] Flusso del protocollo ap1.0: Alice invia a Bob il messaggio "I am Alice" — Slide 22 / Pagina 2

La transazione si riassume nell'invio del messaggio:
\[
\text{Alice} \longrightarrow \text{Bob}: \quad \text{``I am Alice''}
\]

\subsubsection{Scenario di Fallimento}
In un contesto di rete distribuito, le entità comunicanti non sono visivamente collegate né condividono una linea punto-punto fisica isolata. Di conseguenza, Bob non può associare il flusso di bit a un volto o a una presenza fisica.

% [IMMAGINE DEL FILE DI SLIDE CON NOME "1) Ripasso_2"] Fallimento del protocollo ap1.0: Trudy invia a Bob "I am Alice" spacciandosi per lei — Slide 23 / Pagina 3

L'avversario Trudy può intercettare la sessione o semplicemente generare un pacchetto autonomo contenente la medesima stringa:
\[
\text{Trudy} \longrightarrow \text{Bob}: \quad \text{``I am Alice''}
\]
Bob non dispone di alcun criterio logico per discriminare la richiesta autentica di Alice da quella fraudolenta di Trudy. Il protocollo ap1.0 è totalmente vulnerabile all'impersonificazione diretta.

\subsection{Protocollo ap2.0: Autenticazione Tramite Indirizzo IP Sorgente}

Per colmare la lacuna di ap1.0, si potrebbe ipotizzare di ancorare l'identità applicativa alle informazioni fornite dal livello di rete (Network Layer), incapsulando la dichiarazione d'identità all'interno di un pacchetto IP recante l'indirizzo IP sorgente autentico di Alice ($\text{IP}_{\text{Alice}}$).

% [IMMAGINE DEL FILE DI SLIDE CON NOME "1) Ripasso_2"] Schema del protocollo ap2.0: Alice invia a Bob un pacchetto IP con indirizzo sorgente "Alice's IP address" e payload "I am Alice" — Slide 24 / Pagina 4

Il protocollo prevede che Bob accetti l'identità di Alice solo se il campo \emph{Source Address} dell'intestazione del pacchetto corrisponde all'indirizzo IP precedentemente associato ad Alice.

\subsubsection{Scenario di Fallimento: IP Spoofing}
Questo meccanismo ignora la natura non autenticata delle intestazioni standard nel protocollo IPv4/IPv6. Trudy, avendo accesso all'interfaccia di rete e operando tramite raw socket a livello di sistema operativo, può forgiare arbitrariamente i campi dell'header IP.

% [IMMAGINE DEL FILE DI SLIDE CON NOME "1) Ripasso_2"] Fallimento di ap2.0: Trudy genera un pacchetto con IP sorgente contraffatto ("spoofing" dell'indirizzo di Alice) verso Bob — Slide 25 / Pagina 5

L'attacco si concretizza nell'**IP Spoofing**: Trudy costruisce un datagramma impostando l'indirizzo sorgente su $\text{IP}_{\text{Alice}}$ e lo inoltra verso Bob. Bob esamina l'header IP, constata la corrispondenza con l'indirizzo legittimo di Alice e concede l'accesso, validando erroneamente un'identità contraffatta.

\subsection{Protocollo ap3.0: Autenticazione Mediante Password Segreta}

Poiché i dati di instradamento pubblici non sono fidati, si introduce un elemento di segretezza condiviso: Alice e Bob condividono a priori una password segreta. Per farsi riconoscere, Alice trasmette sia la dichiarazione sia la password segreta.

% [IMMAGINE DEL FILE DI SLIDE CON NOME "1) Ripasso_2"] Schema del protocollo ap3.0: Alice invia indirizzo IP, dichiarazione "I'm Alice" e la propria password segreta; Bob risponde con OK — Slide 26 / Pagina 6

I flussi prevedono:
\begin{align*}
    \text{Alice} \longrightarrow \text{Bob}&: \quad [\text{IP}_{\text{Alice}}, \, \text{``I am Alice''}, \, \text{Password}_{\text{Alice}}] \\
    \text{Bob} \longrightarrow \text{Alice}&: \quad [\text{IP}_{\text{Bob}}, \, \text{OK}]
\end{align*}

\subsubsection{Scenario di Fallimento: Attacco Di Riproduzione (\emph{Playback Attack})}
Sebbene questo schema impedisca a un attaccante privo di credenziali di inventare una richiesta valida, esso fallisce totalmente in presenza di un canale di ascolto passivo (\emph{sniffing}).

% [IMMAGINE DEL FILE DI SLIDE CON NOME "1) Ripasso_2"] Fallimento di ap3.0: Trudy intercetta e registra il pacchetto con la password e successivamente lo riproduce verso Bob — Slide 27 / Pagina 7

Trudy, monitorando il traffico sul canale comune, registra l'intero pacchetto trasmesso da Alice (comprensivo di credenziali in chiaro). Successivamente, Trudy avvia una nuova sessione verso Bob trasmettendo una copia esatta del pacchetto catturato. Questo schema offensivo è noto come **attacco di riproduzione** o **replay attack** (\emph{playback attack}). Bob confronta la password ricevuta con quella attesa nel proprio database e convalida la sessione a favore di Trudy.

\subsection{Protocollo ap3.1: Password Cifrata}

Un tentativo intuitivo per neutralizzare lo sniffing consiste nell'applicare la crittografia alle credenziali, trasmettendo la password in forma cifrata.

% [IMMAGINE DEL FILE DI SLIDE CON NOME "1) Ripasso_2"] Protocollo ap3.1: Alice invia la propria password cifrata a Bob — Slide 28 / Pagina 8

Alice trasmette:
\[
\text{Alice} \longrightarrow \text{Bob}: \quad [\text{IP}_{\text{Alice}}, \, \text{``I am Alice''}, \, E(\text{Password}_{\text{Alice}})]
\]

\subsubsection{Scenario di Fallimento: Inefficacia Della Sola Cifratura Statica}
L'approccio implementato in ap3.1 rivela un grave equivoco concettuale: Trudy non ha alcuna necessità di decifrare o comprendere il valore della password in chiaro per violare il protocollo.

% [IMMAGINE DEL FILE DI SLIDE CON NOME "1) Ripasso_2"] Fallimento di ap3.1: Trudy registra la sequenza binaria cifrata e la riproduce verso Bob ottenendo l'accesso ("record and playback still works!") — Slide 29 / Pagina 9

Poiché il testo cifrato risultante da una stringa statica (in assenza di parametri variabili nel tempo) rimane invariato ad ogni iterazione, Trudy registra la sequenza binaria del messaggio cifrato $E(\text{Password}_{\text{Alice}})$ e la riproduce identica in una sessione successiva verso Bob. Bob riceve il testo cifrato, lo decifra con successo estraendo la password corretta, e autentica Trudy. La sola confidenzialità del dato non protegge dall'attacco di riproduzione.

\subsection{Protocollo ap4.0: Meccanismo Sfida-Risposta (\emph{Challenge-Response}) Con Nonce Simmetrico}

Per sconfiggere definitivamente il playback attack, è indispensabile garantire la **freschezza** (\emph{liveness}) della sessione. Ciò si ottiene introducendo un **nonce**.

\begin{theorem}[Definizione Di Nonce]
Un \textbf{nonce} (\emph{number used once-in-a-lifetime}) è un valore numerico generato casualmente da una parte verificatrice, impiegato una sola volta nell'arco dell'intero ciclo di vita del protocollo crittografico, allo scopo di verificare che lo scambio informativo avvenga in tempo reale e prevenire attacchi di rigiocabilità.
\end{theorem}

% [IMMAGINE DEL FILE DI SLIDE CON NOME "1) Ripasso_2"] Flusso del protocollo ap4.0: Alice invia "I am Alice", Bob risponde inviando il nonce casuale R, Alice restituisce R cifrato con la chiave simmetrica condivisa K_{A-B} — Slide 30 / Pagina 10

Il protocollo **ap4.0** adotta uno schema **sfida-risposta** (\emph{challenge-response}) basato su una chiave crittografica simmetrica $K_{A-B}$ condivisa esclusivamente tra Alice e Bob:
\begin{enumerate}[leftmargin=*]
    \item \textbf{Inizializzazione}: Alice dichiara la propria intenzione di autenticarsi:
    \[
    \text{Alice} \longrightarrow \text{Bob}: \quad \text{``I am Alice''}
    \]
    \item \textbf{Sfida (\emph{Challenge})}: Bob genera un nonce casuale $R$ e lo invia ad Alice in chiaro:
    \[
    \text{Bob} \longrightarrow \text{Alice}: \quad R
    \]
    \item \textbf{Risposta (\emph{Response})}: Alice cifra il nonce $R$ impiegando la chiave simmetrica segreta $K_{A-B}$ e restituisce il blocco cifrato a Bob:
    \[
    \text{Alice} \longrightarrow \text{Bob}: \quad K_{A-B}(R)
    \]
    \item \textbf{Verifica}: Bob decifra il blocco ricevuto utilizzando $K_{A-B}$ e confronta il valore ottenuto con il nonce $R$ generato al passo 2.
\end{enumerate}

Poiché il valore $R$ è unico per ogni sessione, un eventuale replay da parte di Trudy del messaggio $K_{A-B}(R)$ registrato in sessioni precedenti risulterà non valido, in quanto Bob avrà generato una sfida $R' \neq R$. Inoltre, poiché solo la vera Alice possiede la chiave segreta $K_{A-B}$, nessun altro soggetto può calcolare la risposta corretta.

\subsubsection{Limiti Operativi Di ap4.0}
Il limite intrinseco di ap4.0 risiede nella necessità di gestire **chiavi simmetriche condivise a priori** tra tutte le coppie di comunicanti. In una rete con $N$ entità eterogenee che devono potersi autenticare a vicenda, il numero totale di chiavi simmetriche da generare, scambiare e proteggere scala in modo quadratico:
\[
\binom{N}{2} = \frac{N(N - 1)}{2} = \mathcal{O}(N^2)
\]
Ciò rende lo schema ap4.0 non scalabile per ambienti cloud distribuiti ad alta dinamicità.

\subsection{Protocollo ap5.0: Autenticazione Asimmetrica A Chiave Pubblica}

Per eliminare l'onere delle chiavi segrete condivise a coppie, il protocollo **ap5.0** adotta la crittografia a chiave asimmetrica coniugata al meccanismo del nonce.

% [IMMAGINE DEL FILE DI SLIDE CON NOME "1) Ripasso_2"] Flusso del protocollo ap5.0: scambio sfida-risposta in cui Alice cifra il nonce R con la propria chiave privata K_A^-, e Bob verifica applicando la chiave pubblica K_A^+ — Slide 31 / Pagina 11

I passaggi del protocollo ap5.0 si articolano come segue:
\begin{enumerate}[leftmargin=*]
    \item Alice notifica a Bob la propria identità:
    \[
    \text{Alice} \longrightarrow \text{Bob}: \quad \text{``I am Alice''}
    \]
    \item Bob estrae un nonce casuale $R$ e lo invia ad Alice:
    \[
    \text{Bob} \longrightarrow \text{Alice}: \quad R
    \]
    \item Alice cifra la sfida $R$ impiegando la propria **chiave privata** $K_A^-$ e trasmette il valore ottenuto a Bob:
    \[
    \text{Alice} \longrightarrow \text{Bob}: \quad K_A^-(R)
    \]
    \item Bob richiede o recupera la **chiave pubblica** di Alice $K_A^+$:
    \[
    \text{Bob} \longrightarrow \text{Alice}: \quad \text{``send me your public key''}
    \]
    \[
    \text{Alice} \longrightarrow \text{Bob}: \quad K_A^+
    \]
    \item Bob applica la chiave pubblica $K_A^+$ al messaggio cifrato ricevuto:
    \[
    K_A^+\bigl(K_A^-(R)\bigr) = R
    \]
    Se il risultato coincide esattamente con il nonce $R$ originariamente trasmesso, Bob deduce che l'interlocutore è in possesso della chiave privata $K_A^-$, e dunque deve necessariamente trattarsi di Alice.
\end{enumerate}

\subsection{La Falla Di Sicurezza Di ap5.0: Attacco Man-In-The-Middle (MITM)}

Nonostante l'eleganza teorica, il protocollo ap5.0 presenta una vulnerabilità critica causata dall'**assenza di autenticazione del canale di distribuzione della chiave pubblica**: Bob non ha alcuna prova oggettiva che la chiave $K_A^+$ ricevuta appartenga effettivamente ad Alice.

Questa debolezza espone il protocollo al classico attacco **Man-In-The-Middle (MITM)**, in cui Trudy si interpone in modo invisibile tra Alice e Bob, spacciandosi per Alice agli occhi di Bob e per Bob agli occhi di Alice.

% [IMMAGINE DEL FILE DI SLIDE CON NOME "1) Ripasso_2"] Rappresentazione completa dell'attacco Man-In-The-Middle su ap5.0: Trudy intercetta e sostituisce le chiavi pubbliche K_A^+ e K_B^+ con le proprie chiavi K_T^+ e K_T^-, decifrando e ricifrando ogni messaggio — Slide 32 / Pagina 12

L'attacco si articola nelle seguenti fasi:
\begin{enumerate}[leftmargin=*]
    \item Alice invia l'inizializzazione: $\text{Alice} \to \text{Trudy}: \text{``I am Alice''}$. Trudy intercetta il pacchetto e lo inoltra a Bob: $\text{Trudy} \to \text{Bob}: \text{``I am Alice''}$.
    \item Bob risponde con il nonce $R$, destinato ad Alice. Trudy intercetta $R$ e lo invia ad Alice.
    \item Alice firma il nonce con la propria chiave privata calcolando $K_A^-(R)$ e lo trasmette. Trudy intercetta la risposta.
    \item Bob richiede la chiave pubblica per validare la firma. In questo istante, **Trudy trasmette a Bob la propria chiave pubblica** $K_T^+$, dichiarando falsamente che appartiene ad Alice.
    \item Trudy sostituisce il nonce firmato da Alice con il proprio nonce firmato $K_T^-(R)$. Bob verifica la consistenza applicando la chiave ricevuta:
    \[
    K_T^+\bigl(K_T^-(R)\bigr) = R
    \]
    Il controllo ha esito positivo e Bob si convince di dialogare autenticamente con Alice, mentre sta interagendo con Trudy.
    \item Simmetricamente, quando in seguito Bob vuole inviare un messaggio confidenziale $m$ ad Alice, lo cifra con la chiave pubblica che crede di Alice (ma che è in realtà $K_T^+$):
    \[
    c = K_T^+(m)
    \]
    \item Trudy intercetta il messaggio $c$, lo decifra istantaneamente con la propria chiave privata:
    \[
    m = K_T^-\bigl(K_T^+(m)\bigr)
    \]
    Trudy ottiene il testo in chiaro $m$, può leggerlo, archiviarlo o alterarlo a piacimento. Successivamente, per non destare sospetti, Trudy ricifra $m$ utilizzando la vera chiave pubblica di Alice $K_A^+$:
    \[
    c' = K_A^+(m)
    \]
    e lo recapita ad Alice. Alice decifra il messaggio mediante la propria chiave privata $K_A^-$:
    \[
    K_A^-\bigl(K_A^+(m)\bigr) = m
    \]
    Sia Alice che Bob credono di scambiare dati in un canale protetto ed esclusivo, mentre Trudy controlla, intercetta e manipola l'intera sessione di comunicazione.
\end{enumerate}

\newpage

\section{Firme Digitali E Integrità Dei Messaggi}

La firma digitale rappresenta la trasposizione algoritmica nel mondo digitale della firma autografa tradizionale, fornendo garanzie legali e crittografiche superiori rispetto alla controparte cartacea.

% [IMMAGINE DEL FILE DI SLIDE CON NOME "1) Ripasso_2"] Concetto fondamentale di firma digitale: analogia con la firma autografa, verificabile e infalsificabile — Slide 33 / Pagina 13

\subsection{Proprietà Fondamentali Della Firma Digitale}

Un meccanismo di firma digitale valido deve soddisfare due requisiti imprescindibili:
\begin{enumerate}[leftmargin=*]
    \item \textbf{Verificabilità}: chiunque riceva il documento deve essere in grado di verificare in modo oggettivo e algoritmico che l'autore dichiarato sia effettivamente colui che ha apposto la firma.
    \item \textbf{Non Falsificabilità (\emph{Nonforgeability})}: deve essere computazionalmente impossibile per chiunque altro (inclusa la parte ricevente) generare una firma valida attribuibile al mittente o modificare il documento firmato senza invalidare la firma stessa.
\end{enumerate}

Da tali proprietà discende direttamente il principio di **non ripudio** (\emph{non-repudiation}): il firmatario non può successivamente negare di aver generato e sottoscritto il documento. Qualora insorga una controversia legale, il destinatario può esibire il documento $m$ congiuntamente alla firma associata di fronte a un'autorità giudiziaria, dimostrando in modo inconfutabile la paternità della firma.

\subsection{Firma Digitale Elementare Su Un Messaggio Completo}

Nel modello più semplice, la firma digitale si ottiene applicando la trasformazione a chiave asimmetrica all'intero corpo del messaggio.

% [IMMAGINE DEL FILE DI SLIDE CON NOME "1) Ripasso_2"] Generazione di una firma digitale semplice: Bob cifra l'intero messaggio m con la propria chiave privata K_B^-, producendo la coppia (m, K_B^-(m)) — Slide 34 / Pagina 14

Dato un messaggio $m$, il mittente Bob produce la firma cifrando $m$ con la propria chiave privata $K_B^-$:
\[
\text{Firma} = K_B^-(m)
\]
Bob trasmette quindi al destinatario la coppia composta dal testo del messaggio e dalla relativa firma:
\[
\text{Bob} \longrightarrow \text{Alice}: \quad \bigl(m, \, K_B^-(m)\bigr)
\]

% [IMMAGINE DEL FILE DI SLIDE CON NOME "1) Ripasso_2"] Verifica della firma digitale da parte di Alice mediante la chiave pubblica K_B^+ di Bob e verifica delle proprietà di integrità e non ripudio — Slide 35 / Pagina 15

All'atto della ricezione, Alice verifica la validità della firma applicando la chiave pubblica di Bob $K_B^+$ al testo cifrato ricevuto:
\begin{equation}
\text{Verifica: calcola } K_B^+\bigl(K_B^-(m)\bigr) \stackrel{?}{=} m
\end{equation}
Se l'uguaglianza è verificata, Alice ottiene tre conferme matematiche:
\begin{enumerate}[noitemsep]
    \item Solo il detentore di $K_B^-$ (Bob) ha potuto apporre la firma;
    \item Nessun altro soggetto sulla rete ha potuto falsificarla;
    \item Il messaggio firmato è esattamente $m$ e non un testo alterato $m'$. Qualora Trudy avesse modificato anche un singolo bit di $m$ trasformandolo in $m'$, la decifratura della firma non coinciderebbe con il nuovo testo ($K_B^+\bigl(K_B^-(m)\bigr) = m \neq m'$).
\end{enumerate}

\subsection{Digest Del Messaggio E Funzioni Hash}

La cifratura asimmetrica dell'intero corpo del messaggio descritta nel modello elementare presenta un collo di bottiglia prestazionale proibitivo. Gli algoritmi asimmetrici (come RSA) si basano su moltiplicazioni ed elevamenti a potenza tra numeri interi di grandi dimensioni (2048--4096 bit), risultando da tre a quattro ordini di grandezza più lenti rispetto ai cifrari simmetrici. Applicare RSA a documenti di grandi dimensioni (ad esempio file da centinaia di megabyte) provocherebbe un consumo inaccettabile di CPU e tempi di elaborazione incompatibili con le architetture real-time.

La soluzione ingegneristica consiste nel disaccoppiare l'operazione di firma dal volume del messaggio, generando una **impronta digitale** (\emph{digital fingerprint}) o **digest** a dimensione fissa, computazionalmente leggero da produrre.

% [IMMAGINE DEL FILE DI SLIDE CON NOME "1) Ripasso_2"] Schema del Message Digest: elaborazione di un messaggio di grande dimensione m tramite funzione di hash H per produrre un digest a dimensione fissa H(m) — Slide 36 / Pagina 16

Dato un messaggio $m$ di lunghezza arbitraria, si applica una funzione di hash crittografica $H$:
\begin{equation}
\text{Digest} = H(m)
\end{equation}
Le proprietà crittografiche imposte alla funzione $H$ sono:
\begin{itemize}[leftmargin=*]
    \item \textbf{Compressione (\emph{Many-to-one})}: la funzione proietta uno spazio di ingresso infinito su uno spazio discreto di uscite a dimensione prefissata.
    \item \textbf{Efficienza di calcolo}: dato $m$, il calcolo di $H(m)$ deve richiedere minime risorse di calcolo.
    \item \textbf{Resistenza alla preimmagine (Unidirezionalità)}: dato un qualsiasi digest $x$, deve risultare computazionalmente impossibile individuare un messaggio $m$ tale per cui $H(m) = x$.
    \item \textbf{Resistenza forte alle collisioni}: deve essere computazionalmente impossibile trovare due messaggi distinti $m_1 \neq m_2$ tali che $H(m_1) = H(m_2)$.
\end{itemize}

\subsection{Algoritmi Di Hash Noti}

Nel corso dell'evoluzione della sicurezza informatica sono stati standardizzati vari algoritmi di hash:

% [IMMAGINE DEL FILE DI SLIDE CON NOME "1) Ripasso_2"] Tassonomia degli algoritmi di hash: MD5 (RFC 1321 a 128 bit) e SHA-1 (FIPS PUB 180-1 a 160 bit) — Slide 37 / Pagina 17

\begin{itemize}[leftmargin=*]
    \item \textbf{MD5 (\emph{Message Digest 5})}: standardizzato da Ron Rivest nella RFC 1321, elabora il testo in un processo a 4 passi generando un digest a \textbf{128 bit}. Sebbene storicamente diffusissimo, MD5 è oggi crittograficamente compromesso a causa di vulnerabilità alle collisioni dimostrate sperimentalmente.
    \item \textbf{SHA-1 (\emph{Secure Hash Algorithm 1})}: sviluppato dal NIST e specificato nello standard federale FIPS PUB 180-1, produce un digest a \textbf{160 bit}. Anch'esso considerato oggi superato per impieghi di alta sicurezza (a favore delle famiglie SHA-2 e SHA-3), ha costituito a lungo il cardine dei protocolli distribuiti.
\end{itemize}

\subsection{Architettura Moderna Della Firma Digitale}

Integrando le funzioni di hash con la crittografia asimmetrica, si realizza il paradigma moderno: \textbf{la firma digitale coincide con la cifratura del solo digest del messaggio}.

% [IMMAGINE DEL FILE DI SLIDE CON NOME "1) Ripasso_2"] Schema architetturale completo della firma digitale: Bob calcola l'hash H(m), lo cifra con K_B^- e lo invia insieme a m; Alice calcola H(m) localmente, decifra la firma con K_B^+ e confronta i due digest — Slide 38 / Pagina 18

L'architettura operativa si articola in due fasi distinte:

\subsubsection{Fase 1: Generazione della Firma da parte di Bob}
\begin{enumerate}[leftmargin=*]
    \item Bob prende il messaggio in chiaro $m$ di dimensione arbitraria;
    \item Calcola il digest a lunghezza fissa mediante la funzione di hash: $h = H(m)$;
    \item Cifra il solo digest $h$ applicando la propria chiave privata:
    \[
    S = K_B^-\bigl(H(m)\bigr)
    \]
    \item Concatena il messaggio originario in chiaro con la firma calcolata e invia il pacchetto:
    \[
    \text{Bob} \longrightarrow \text{Alice}: \quad \Bigl(m, \, K_B^-\bigl(H(m)\bigr)\Bigr)
    \]
\end{enumerate}

\subsubsection{Fase 2: Verifica della Firma e dell'Integrità da parte di Alice}
All'arrivo del blocco $(m, S)$, Alice esegue una doppia elaborazione indipendente:
\begin{enumerate}[leftmargin=*]
    \item Estrae il messaggio in chiaro $m$ e calcola localmente il valore di hash applicando la medesima funzione standard:
    \[
    h_{\text{locale}} = H(m)
    \]
    \item Prende la firma $S = K_B^-(H(m))$ e applica la chiave pubblica di Bob $K_B^+$ per decifrare il digest originale:
    \[
    h_{\text{ricevuto}} = K_B^+\bigl(S\bigr) = K_B^+\Bigl(K_B^-\bigl(H(m)\bigr)\Bigr) = H(m)
    \]
    \item Esegue il confronto di uguaglianza:
    \begin{equation}
    \boxed{h_{\text{locale}} \stackrel{?}{=} h_{\text{ricevuto}}}
    \end{equation}
\end{enumerate}

Se $h_{\text{locale}} = h_{\text{ricevuto}}$, l'intero protocollo ha successo: viene attestata con certezza matematica l'autenticità dell'autore (solo Bob conosce $K_B^-$) e l'assoluta integrità del messaggio (qualsiasi alterazione di $m$ produrrebbe un valore $h_{\text{locale}}$ divergente da quello decifrato dalla firma).

\newpage

\section{Infrastrutture A Chiave Pubblica (PKI) E Autorità Di Certificazione}

La trattazione dei protocolli di firma e autenticazione asimmetrica presuppone che il ricevente conosca con certezza assoluta la chiave pubblica legittima della controparte. Come dimostrato dall'attacco Man-In-The-Middle al protocollo ap5.0, in assenza di un meccanismo di associazione fidata tra un'identità e la sua chiave pubblica, l'intero edificio crittografico crolla.

Questo problema viene risolto dall'introduzione delle **Autorità Di Certificazione** (\emph{Certification Authority}, CA) nell'ambito delle **Public Key Infrastructure (PKI)**.

\subsection{Ruolo Della Certification Authority}

Una Certification Authority è un'entità terza di fiducia (\emph{Trusted Third Party}) il cui scopo primario è **legare in modo crittograficamente inoppugnabile una specifica chiave pubblica all'identità di una determinata entità** $E$ (che può essere una persona fisica, un server web, un router o un microservizio cloud).

% [IMMAGINE DEL FILE DI SLIDE CON NOME "1) Ripasso_2"] Processo di rilascio del certificato digitale: l'entità Bob registra la propria chiave pubblica K_B^+ fornendo prove di identità alla CA, la quale genera il certificato firmandolo con la propria chiave privata K_{CA}^- — Slide 39 / Pagina 19

Il processo di accreditamento e rilascio del certificato digitale segue un iter strutturato:
\begin{enumerate}[leftmargin=*]
    \item L'entità $E$ (ad esempio Bob) genera autonomamente la propria coppia di chiavi $(K_B^+, K_B^-)$;
    \item Bob presenta la propria chiave pubblica $K_B^+$ alla CA, allegando una rigorosa prova di identità documentale o tecnica (\emph{proof of identity});
    \item La CA esegue la verifica della validità dell'identità di Bob;
    \item La CA crea una struttura dati standardizzata (un **certificato digitale**) contenente l'identità di Bob, la sua chiave pubblica $K_B^+$, date di validità ed estensioni di conformità;
    \item La CA appone la propria firma digitale sul certificato cifrandone il digest con la propria **chiave privata** $K_{\text{CA}}^-$:
    \begin{equation}
    \text{Certificato}_{\text{Bob}} = \Bigl(\text{Dati}_{\text{Bob}}, \, K_B^+, \, K_{\text{CA}}^-\bigl(H(\text{Dati}_{\text{Bob}} \parallel K_B^+)\bigr)\Bigr)
    \end{equation}
\end{enumerate}

La CA agisce da garante universale: il certificato firmato sancisce formalmente che \emph{``la chiave pubblica ivi contenuta appartiene a tutti gli effetti all'entità Bob''}.

\subsection{Verifica Del Certificato Digitale}

Quando un utente o un'applicazione (Alice) deve comunicare in modo sicuro con Bob, non richiede più la chiave pubblica direttamente a quest'ultimo in modo non verificato.

% [IMMAGINE DEL FILE DI SLIDE CON NOME "1) Ripasso_2"] Verifica del certificato: Alice ottiene il certificato di Bob, applica la chiave pubblica fidata della CA K_{CA}^+ e ne estrae la chiave pubblica autentica K_B^+ — Slide 40 / Pagina 20

La dinamica di verifica opera come segue:
\begin{enumerate}[leftmargin=*]
    \item Alice acquisisce il certificato digitale di Bob (ricevendolo direttamente da Bob all'avvio della sessione TLS o prelevandolo da un repository pubblico);
    \item Alice isola la firma apposta dalla CA sul certificato e vi applica la **chiave pubblica della CA** $K_{\text{CA}}^+$:
    \[
    K_{\text{CA}}^+\Bigl(K_{\text{CA}}^-\bigl(H(\text{Dati} \parallel K_B^+)\bigr)\Bigr) = H(\text{Dati} \parallel K_B^+)
    \]
    \item Alice calcola in modo autonomo il digest del corpo del certificato e lo confronta con il valore decifrato;
    \item Riscontrata l'uguaglianza dei valori di hash, Alice ha la certezza matematica che il certificato non è stato manipolato da terzi e che la chiave pubblica $K_B^+$ estratta appartiene inconfutabilmente a Bob.
\end{enumerate}

La radice di fiducia (\emph{Root of Trust}) si sposta così sulla conoscenza preventiva e sicura della chiave pubblica della CA $K_{\text{CA}}^+$, la quale viene preinstallata e certificata nativamente all'interno dei sistemi operativi e dei browser web accreditati.

\newpage

\section{Ingegneria Dei Sistemi Di Posta Elettronica Sicura (Secure E-Mail)}

I protocolli applicativi per la posta elettronica (come SMTP) sono stati originariamente concepiti senza meccanismi di protezione nativi, trasmettendo i dati in chiaro su reti non fidate. L'applicazione dei principi crittografici a tale ambito deve tenere conto della natura asincrona della posta elettronica: mittente e destinatario non sono contemporaneamente attivi (\emph{online}) in una sessione sincrona a bassa latenza, impedendo meccanismi di negoziazione interattiva bidirezionale in tempo reale.

Vengono analizzati i tre scenari di progettazione per implementare una trasmissione email protetta tra Alice e Bob.

\subsection{Scenario 1: Garanzia Di Sola Riservatezza}

Alice desidera inviare una mail confidenziale $m$ a Bob in modo tale che nessun intermediario (o server di posta di transito) possa leggerne il contenuto.

Per ottimizzare le prestazioni, non è pensabile cifrare un messaggio di dimensioni consistenti con la chiave pubblica di Bob $K_B^+$. Si adotta quindi un approccio a **cifrario ibrido** (\emph{envelope encryption}).

% [IMMAGINE DEL FILE DI SLIDE CON NOME "1) Ripasso_2"] Flusso di Secure E-Mail per sola confidenzialità: Alice cifra il messaggio m con una chiave simmetrica casuale K_S, cifra la chiave K_S con la chiave pubblica di Bob K_B^+, e invia la busta crittografica su Internet; Bob inverte il processo — Slide 41-42 / Pagine 21-22

Il procedimento implementato da Alice si suddivide nei seguenti passi:
\begin{enumerate}[leftmargin=*]
    \item Alice genera localmente una **chiave simmetrica segreta casuale e monouso** $K_S$ (chiave di sessione);
    \item Cifra il corpo del messaggio $m$ impiegando $K_S$ attraverso un cifrario simmetrico ad alte prestazioni (es. AES):
    \[
    c_{\text{msg}} = K_S(m)
    \]
    \item Cifra la sola chiave simmetrica $K_S$ impiegando la chiave pubblica di Bob $K_B^+$:
    \[
    c_{\text{chiave}} = K_B^+(K_S)
    \]
    \item Inoltra la composizione delle due porzioni cifrate lungo la rete verso Bob:
    \[
    \text{Alice} \longrightarrow \text{Bob}: \quad \Bigl(K_S(m), \, K_B^+(K_S)\Bigr)
    \]
\end{enumerate}

All'atto della ricezione, Bob esegue l'operazione complementare:
\begin{enumerate}[leftmargin=*]
    \item Estrae la componente $K_B^+(K_S)$ e vi applica la propria chiave privata $K_B^-$, recuperando la chiave di sessione simmetrica:
    \[
    K_S = K_B^-\bigl(K_B^+(K_S)\bigr)
    \]
    \item Utilizza la chiave simmetrica $K_S$ appena recuperata per decifrare il payload applicativo $K_S(m)$, estraendo il messaggio originale in chiaro:
    \[
    m = K_S\bigl(K_S(m)\bigr)
    \]
\end{enumerate}
In questo modo si ottiene la massima efficienza computazionale garantendo al contempo che solo ed esclusivamente Bob possa accedere alla chiave di sessione e al payload.

\subsection{Scenario 2: Autenticazione Del Mittente Ed Integrità Del Messaggio}

Nel secondo scenario, la confidenzialità non è richiesta (il contenuto del messaggio può viaggiare in chiaro), ma Alice desidera offrire a Bob la certezza dell'integrità del testo e l'autenticità della provenienza, impedendo a chiunque di impersonarla o alterare il testo.

% [IMMAGINE DEL FILE DI SLIDE CON NOME "1) Ripasso_2"] Architettura di Secure E-Mail per autenticazione e integrità: Alice genera il digest H(m), lo cifra con la propria chiave privata K_A^- e lo invia insieme a m in chiaro; Bob verifica decifrando con K_A^+ e confrontando con l'hash locale — Slide 43 / Pagina 23

Il flusso operativo è basato sulla firma digitale ottimizzata del digest:
\begin{enumerate}[leftmargin=*]
    \item Alice applica la funzione di hash $H$ al messaggio $m$ ottenendo il digest:
    \[
    h = H(m)
    \]
    \item Cifra il digest con la propria chiave privata $K_A^-$, generando la firma crittografica:
    \[
    S = K_A^-\bigl(H(m)\bigr)
    \]
    \item Trasmette a Bob il messaggio originale in chiaro concatenato alla firma:
    \[
    \text{Alice} \longrightarrow \text{Bob}: \quad \Bigl(m, \, K_A^-\bigl(H(m)\bigr)\Bigr)
    \]
    \item Bob riceve il blocco e conduce la verifica duale: calcola $H(m)$ dal testo in chiaro ricevuto, decifra la firma mediante la chiave pubblica di Alice $K_A^+$ calcolando $K_A^+\bigl(K_A^-(H(m))\bigr) = H(m)$, e ne verifica la coincidenza.
\end{enumerate}

\subsection{Scenario 3: Schema Ibrido Completo A Tre Chiavi}

Il terzo scenario rappresenta la configurazione ingegneristica completa: Alice intende garantire simultaneamente la **riservatezza** (\emph{secrecy}), l'**autenticazione del mittente** (\emph{sender authentication}) e l'**integrità del messaggio** (\emph{message integrity}), unita al non ripudio.

Per soddisfare congiuntamente tutte queste proprietà su un messaggio arbitrariamente lungo in modo computazionalmente efficiente, Alice deve coordinare l'impiego sinergico di tre chiavi crittografiche distinte:
\begin{itemize}[noitemsep]
    \item La propria **chiave privata** $K_A^-$ (garantisce autenticazione, integrità e non ripudio);
    \item La **chiave pubblica di Bob** $K_B^+$ (garantisce la confidenzialità della chiave simmetrica verso il solo destinatario);
    \item Una **nuova chiave simmetrica effimera** $K_S$ generata appositamente per la sessione (garantisce l'efficienza della cifratura massiva).
\end{itemize}

% [IMMAGINE DEL FILE DI SLIDE CON NOME "1) Ripasso_2"] Architettura di Secure E-Mail completa a tre chiavi: calcolo del digest H(m), firma con K_A^-, cifratura simmetrica del bundle con K_S, e incapsulamento della chiave K_S tramite la chiave pubblica di Bob K_B^+ — Slide 44 / Pagina 24

L'architettura di composizione crittografica segue un ordine rigoroso di incapsulamento:
\begin{enumerate}[leftmargin=*]
    \item \textbf{Generazione del Digest e Firma Digitale}:
    Alice calcola preliminarmente il digest del messaggio originario $m$ e lo firma impiegando la propria chiave privata $K_A^-$:
    \[
    S = K_A^-\bigl(H(m)\bigr)
    \]
    Viene così creato il blocco logicamente autenticato:
    \[
    B = \Bigl(m \parallel K_A^-\bigl(H(m)\bigr)\Bigr)
    \]

    \item \textbf{Cifratura Simmetrica del Blocco Combinato}:
    Alice genera la chiave simmetrica di sessione $K_S$ e la impiega per cifrare l'intero blocco $B$ (contenente sia il messaggio in chiaro sia la firma):
    \[
    C_1 = K_S\Bigl(m \parallel K_A^-\bigl(H(m)\bigr)\Bigr)
    \]

    \item \textbf{Cifratura Asimmetrica della Chiave di Sessione}:
    Per consentire al solo Bob di estrarre $K_S$, Alice cifra la chiave simmetrica utilizzando la chiave pubblica del destinatario $K_B^+$:
    \[
    C_2 = K_B^+(K_S)
    \]

    \item \textbf{Spedizione del Pacchetto Complessivo}:
    Alice unisce $C_1$ e $C_2$ e li trasmette sul canale di rete pubblico verso Bob:
    \begin{equation}
    \boxed{\text{Alice} \longrightarrow \text{Bob}: \quad \Biggl(K_S\Bigl(m \parallel K_A^-\bigl(H(m)\bigr)\Bigr), \quad K_B^+(K_S)\Biggr)}
    \end{equation}
\end{enumerate}

\subsubsection{Procedura di Elaborazione e Verifica da parte di Bob}
All'arrivo del messaggio composto $(C_1, C_2)$, Bob esegue il processo inverso procedendo dall'involucro esterno verso il dato interno:
\begin{enumerate}[leftmargin=*]
    \item \textbf{Recupero della Chiave Simmetrica}: Bob applica la propria chiave privata $K_B^-$ a $C_2$:
    \[
    K_S = K_B^-\bigl(C_2\bigr) = K_B^-\bigl(K_B^+(K_S)\bigr)
    \]
    \item \textbf{Decifratura del Carico Utile}: Bob impiega la chiave $K_S$ appena ottenuta per decifrare $C_1$, riottenendo in chiaro il messaggio e la firma:
    \[
    K_S(C_1) = \Bigl(m \parallel K_A^-\bigl(H(m)\bigr)\Bigr)
    \]
    \item \textbf{Verifica di Autenticità e Integrità}: Bob applica la chiave pubblica di Alice $K_A^+$ alla firma estratta:
    \[
    K_A^+\Bigl(K_A^-\bigl(H(m)\bigr)\Bigr) = H(m)
    \]
    Calcola autonomamente $H(m)$ a partire dal testo $m$ recuperato e confronta i due digest. La corrispondenza attesta congiuntamente che il testo non è stato alterato, che è stato originato indiscutibilmente da Alice, e che nessun'altra entità al di fuori di Bob è stata in grado di leggerlo durante il transito.
\end{enumerate}



\end{document}
